% !TeX program = xelatex
\documentclass[12pt,a4paper]{article}
\usepackage{fontspec}
\setmainfont{Liberation Serif}   % oder Times New Roman, Arial
\usepackage[ngerman]{babel}      % deutsche Silbentrennung und Lokalisierung

\usepackage{amsmath,amssymb,amsthm,mathtools}
\usepackage{geometry}
\usepackage{booktabs}
\usepackage{longtable}
\usepackage{hyperref}
\usepackage{rotating}
\usepackage{pdflscape}
\usepackage{caption}
\geometry{margin=2cm}

\newtheorem{definition}{Definition}
\newcommand{\Keff}{K_{\mathrm{eff}}}
\newcommand{\kappac}{\kappa_c}
\newcommand{\be}{\beta}
\newcommand{\ga}{\gamma}
\newcommand{\al}{\alpha}
\newcommand{\Gcoeff}{\Gamma}
\newcommand{\bracket}[1]{\left\langle #1 \right\rangle}

\begin{document}
	
	\begin{titlepage}
		\centering
		\vspace*{1cm}
		\Huge\textbf{YUCT: Der periodische Lagrange-Operator koordinierter Materie} \\
		\vspace{0.3cm}
		\Large\textbf{(Anhang PLCM)} \\
		\vspace{0.5cm}
		\large\textit{Ein einheitlicher Rahmen für Elementeigenschaften, Legierungsdesign und Materialinformatik}
		\vspace{2cm}
		
		\Large Alexey V. Yakushev \\
		\large Yakushev Research \\
		\url{https://yuct.org/}
		\vspace{2cm}
		
		\vspace{2cm}
		\textbf DOI: \url{https://doi.org/10.5281/zenodo.18444598}\\
		
		\vspace{1cm}
		
		\large Version 2.2, März 2026
		\vfill
	\end{titlepage}
	
	\tableofcontents
	\newpage
	
	\section{Einleitung}
	
	Der Periodische Lagrange-Operator Koordinierter Materie (PLCM) ist die materialwissenschaftliche Ausprägung der Yakushev Unified Coordination Theory (YUCT). Er integriert alle bekannten Eigenschaften chemischer Elemente, ihre Kombinationen zu Legierungen und Verarbeitungsmethoden in einen einzigen mathematischen Rahmen. Aufbauend auf der 120-Sektor-Architektur von YUCT führt PLCM ein:
	\begin{itemize}
		\item Einen vollständigen Satz von Observablen für jedes Element (atomare, physikalische, mechanische, thermische, magnetische, katalytische, gravitative).
		\item Kopplungskoeffizienten \(\kappa_{sr}\), die Wechselwirkungen zwischen Elementen, Kristallstrukturen und Verarbeitung codieren.
		\item Eine Vorhersageformel für Legierungseigenschaften basierend auf Elementdaten und Kopplung.
		\item Integration mit thermodynamischen Datenbanken (CALPHAD), um vorhandenes Wissen zu nutzen.
	\end{itemize}
	Der zentrale Parameter ist die Koordinationseffizienz \(\Keff\), die alle Eigenschaften durch den universellen Exponenten \(\be = 0,67\) skaliert (Anhang L). PLCM verwandelt so das Periodensystem in ein generatives Modell für das Materialdesign.
	
	\section{Der PLCM-Lagrange-Operator}
	
	Der Lagrange-Operator wird in derselben Funktionalform wie der allgemeine YUCT-Lagrange-Operator geschrieben:
	\[
	\mathcal{L}_{\text{PLCM}} = \int d^{19}X \sqrt{-G} \,
	\exp\Bigg[ \sum_{s=1}^{120} \lambda_s L_s + \sum_{1\le s<r\le 120} \kappa_{sr} \mathrm{Tr}\big(\Psi_{sr} \cdot O_s \cdot O_r^\dagger\big) \Bigg],
	\]
	wobei \(L_s\) die kinetischen Terme für die Observablen \(O_s\) des Sektors \(s\) enthält und \(\kappa_{sr}\) die Kopplungskoeffizienten sind. Die Sektorzerlegung wird unten detailliert beschrieben.
	
	\subsection{Sektorstruktur (120 Sektoren)}
	
	\begin{enumerate}
		\item \textbf{Elemente (Sektoren 1–80)}: Jedes Element \(Z\) belegt Sektor \(s=Z\).
		\item \textbf{Seltene Erden \& Actinoide (81–100)}: Eigene Sektoren für Lanthanoide und Actinoide.
		\item \textbf{Kristallstrukturen (101–115)}: Typen von Gittern (BCC, FCC, HCP, Diamant, ionisch, intermetallisch, HEA usw.).
		\item \textbf{Legierungen \& Verbindungen (116–125)}: Spezifische Materialien (Bronze, Stahl, Ni\(_{3}\)Al, YBCO usw.).
		\item \textbf{Verarbeitung (126–130)}: Thermische, Verformungs-, Legierungs-, Strahlungs-, additive Fertigung.
		\item \textbf{Thermodynamische Datenbank (131)}: CALPHAD-Integration.
	\end{enumerate}
	
	\subsection{Observablen in den Elementsektoren}
	
	Jeder Elementsektor speichert die folgenden Eigenschaften (Quellen: CRC Handbook, NIST, Pauling, metallurgische Datenbanken):
	
	\begin{longtable}{p{4cm} p{3cm} p{8cm}}
		\caption{Observablen in den Elementsektoren (1–80)} \\
		\toprule
		\textbf{Eigenschaft} & \textbf{Symbol} & \textbf{Einheiten / Anmerkungen} \\
		\midrule
		Ordnungszahl & \(Z\) & – \\
		Atommasse & \(A\) & u \\
		Kovalenter Radius & \(r_{\text{cov}}\) & pm \\
		Ionenradius (Shannon) & \(r_{\text{ion}}\) & pm \\
		Elektronegativität (Pauling) & \(\chi\) & – \\
		Erste Ionisierungsenergie & \(I_1\) & eV \\
		Elektronenaffinität & \(E_{\text{ea}}\) & eV \\
		Koordinationseffizienz & \(\Keff\) & aus Anhang C \\
		Gravitationskoeffizient & \(\Gcoeff\) & \(\Gcoeff = \kappa_c \alpha \Keff^{\beta}\) \\
		Schmelzpunkt & \(T_m\) & K \\
		Siedepunkt & \(T_b\) & K \\
		Curie-Temperatur (falls ferromagnetisch) & \(T_C\) & K \\
		Sprödbruchübergangstemperatur & \(T_{\text{brittle}}\) & K \\
		Oxidationsbeginn (in O\(_{2}\)) & \(T_{\text{ox}}\) & K \\
		Dichte & \(\rho\) & g/cm³ \\
		Elastizitätsmodul & \(E\) & GPa \\
		Schubmodul & \(G\) & GPa \\
		Poissonzahl & \(\nu\) & – \\
		Härte (Brinell) & \(H_B\) & – \\
		Zugfestigkeit & \(\sigma_B\) & MPa \\
		Elektrische Leitfähigkeit & \(\sigma\) & \% IACS \\
		Wärmeleitfähigkeit & \(\kappa\) & W/(m·K) \\
		Magnetische Suszeptibilität & \(\chi_m\) & – \\
		Katalytische Aktivität (TOF) & TOF & s\(^{-1}\) \\
		Gefahrenklasse & – & GHS / GOST \\
		\bottomrule
	\end{longtable}
	
	% --- BLOCK 1 ---
	\subsection{Elementeigenschaftsdaten (Beispiel: erste 20 Elemente)}
	
	Tabelle \ref{tab:elem-data-first20} listet numerische Werte der wichtigsten Observablen für Elemente \(Z=1\) bis \(20\). Der vollständige Datensatz für alle 118 Elemente ist im Repository \url{https://github.com/Alexey-Yakushev-YUCT/YPSDC/} verfügbar. Die Werte stammen, sofern nicht anders angegeben, aus Standardreferenzen (CRC Handbook, NIST, Pauling).
	
	\textbf{Anmerkungen:}
	\begin{itemize}
		\item \(\Keff\) (Koordinationseffizienz) wird nach der in Anhang C beschriebenen Methode berechnet. Die Werte für \(Z>80\) werden mit derselben Formel extrapoliert.
		\item \(\Gamma_{\text{rel}}\) (relativer Gravitationskoeffizient) wird berechnet als \(\Gamma_{\text{rel}} = \kappa_c \alpha \Keff^{\beta}\) mit \(\kappa_c=0,33\), \(\beta=0,67\). Die absolute Skala \(\alpha\) ist derzeit für relative Vergleiche auf 1 gesetzt; eine vollständige Kalibrierung wird durchgeführt, sobald experimentelle Daten verfügbar sind.
		\item Der Elastizitätsmodul \(E\) ist für die stabilste feste Phase bei Raumtemperatur angegeben. Für Elemente, die unter Standardbedingungen nicht fest sind, wird „–“ verwendet. Für Bor entspricht der Wert \(\beta\)-rhomboidem Bor (\(E \approx 460\) GPa). Für Kohlenstoff wird der Wert für Diamant (1050 GPa) aufgeführt; wenn Kohlenstoff in Legierungen (z. B. Stahl) verwendet wird, sollte der effektive Modul an die erwartete Phase (Graphit, Carbide) angepasst werden. Eine vollständige Tabelle der phasenabhängigen Eigenschaften ist im Repository verfügbar.
		\item Einheiten: \(A\) in u, \(r_{\text{cov}}\) in pm, \(I_1\) in eV, \(T_m\) in K, \(E\) in GPa.
	\end{itemize}
	
	\begin{longtable}{p{0.6cm}p{1.2cm}p{1.2cm}p{1.2cm}p{1.2cm}p{1.2cm}p{1.5cm}p{1.2cm}p{1.2cm}p{1.2cm}}
		\caption{Elementeigenschaften für Z=1–20} \label{tab:elem-data-first20} \\
		\toprule
		\(Z\) & Symbol & \(A\) (u) & \(r_{\text{cov}}\) (pm) & \(\chi\) & \(I_1\) (eV) & \(\Keff\) & \(\Gamma_{\text{rel}}\) & \(T_m\) (K) & \(E\) (GPa) \\
		\midrule
		\endfirsthead
		\toprule
		\(Z\) & Symbol & \(A\) & \(r_{\text{cov}}\) & \(\chi\) & \(I_1\) & \(\Keff\) & \(\Gamma_{\text{rel}}\) & \(T_m\) & \(E\) \\
		\midrule
		\endhead
		1  & H  & 1,008  & 37   & 2,20 & 13,60 & 1,00  & 0,33  & 14,0  & – \\
		2  & He & 4,002  & 32   & –    & 24,59 & 0,153 & 0,07  & 0,95  & – \\
		3  & Li & 6,94   & 128  & 0,98 & 5,39  & 1,85  & 0,48  & 453,7 & 4,9 \\
		4  & Be & 9,01   & 96   & 1,57 & 9,32  & 2,34  & 0,56  & 1560  & 287 \\
		5  & B  & 10,81  & 84   & 2,04 & 8,30  & 3,12  & 0,69  & 2573  & 460\footnotemark \\
		6  & C  & 12,01  & 77   & 2,55 & 11,26 & 5,67  & 1,02  & 3800  & 1050\footnotemark \\
		7  & N  & 14,01  & 75   & 3,04 & 14,53 & 4,23  & 0,85  & 63,2  & – \\
		8  & O  & 16,00  & 73   & 3,44 & 13,62 & 6,89  & 1,15  & 54,4  & – \\
		9  & F  & 19,00  & 71   & 3,98 & 17,42 & 7,45  & 1,21  & 53,5  & – \\
		10 & Ne & 20,18  & 69   & –    & 21,56 & 0,181 & 0,08  & 24,6  & – \\
		11 & Na & 22,99  & 154  & 0,93 & 5,14  & 2,13  & 0,52  & 371,0 & 10 \\
		12 & Mg & 24,31  & 130  & 1,31 & 7,65  & 2,57  & 0,59  & 923,0 & 45 \\
		13 & Al & 26,98  & 118  & 1,61 & 5,99  & 3,01  & 0,67  & 933,5 & 70 \\
		14 & Si & 28,09  & 111  & 1,90 & 8,15  & 4,33  & 0,86  & 1687  & 130 \\
		15 & P  & 30,97  & 106  & 2,19 & 10,49 & 3,79  & 0,78  & 317   & – \\
		16 & S  & 32,06  & 102  & 2,58 & 10,36 & 5,12  & 0,96  & 388   & – \\
		17 & Cl & 35,45  & 99   & 3,16 & 12,97 & 4,57  & 0,90  & 172   & – \\
		18 & Ar & 39,95  & 98   & –    & 15,76 & 0,213 & 0,08  & 83,8  & – \\
		19 & K  & 39,10  & 203  & 0,82 & 4,34  & 2,38  & 0,56  & 336,5 & 3,5 \\
		20 & Ca & 40,08  & 174  & 1,00 & 6,11  & 2,85  & 0,63  & 1115  & 20 \\
		\bottomrule
	\end{longtable}
	\footnotetext{Für Bor entspricht der Wert \(\beta\)-rhomboidem Bor. Für Kohlenstoff ist der Wert für Diamant angegeben; in Legierungen kann Kohlenstoff als Graphit oder Carbid vorliegen, was eine phasenspezifische Korrektur erfordert (siehe Repository).}
	\vspace{0.2cm}
	*Vollständiger Datensatz für alle Elemente verfügbar im Repository.
	% --- ENDE BLOCK 1 ---
	
	% --- FORTSETZUNG NACH BLOCK 1 ---
	
	\subsection*{Elementeigenschaften für Z=21–40}
	\addcontentsline{toc}{subsection}{Elementeigenschaften für Z=21–40}
	
	\begin{longtable}{p{0.6cm}p{1.2cm}p{1.2cm}p{1.2cm}p{1.2cm}p{1.2cm}p{1.5cm}p{1.2cm}p{1.2cm}p{1.2cm}}
		\caption{Elementeigenschaften für Z=21–40} \label{tab:elem-data-21to40} \\
		\toprule
		\(Z\) & Symbol & \(A\) (u) & \(r_{\text{cov}}\) (pm) & \(\chi\) & \(I_1\) (eV) & \(\Keff\) & \(\Gamma_{\text{rel}}\) & \(T_m\) (K) & \(E\) (GPa) \\
		\midrule
		\endfirsthead
		\toprule
		\(Z\) & Symbol & \(A\) & \(r_{\text{cov}}\) & \(\chi\) & \(I_1\) & \(\Keff\) & \(\Gamma_{\text{rel}}\) & \(T_m\) & \(E\) \\
		\midrule
		\endhead
		21 & Sc & 44,96 & 144 & 1,36 & 6,56 & 3,12 & 0,69 & 1814 & 74 \\
		22 & Ti & 47,87 & 132 & 1,54 & 6,83 & 3,57 & 0,75 & 1941 & 116 \\
		23 & V  & 50,94 & 122 & 1,63 & 6,75 & 3,79 & 0,78 & 2183 & 128 \\
		24 & Cr & 52,00 & 118 & 1,66 & 6,77 & 4,01 & 0,81 & 2130 & 279 \\
		25 & Mn & 54,94 & 117 & 1,55 & 7,43 & 3,85 & 0,79 & 1519 & 198 \\
		26 & Fe & 55,85 & 117 & 1,83 & 7,90 & 4,26 & 0,84 & 1811 & 211 \\
		27 & Co & 58,93 & 116 & 1,88 & 7,88 & 4,38 & 0,86 & 1768 & 209 \\
		28 & Ni & 58,69 & 115 & 1,91 & 7,64 & 4,50 & 0,87 & 1728 & 200 \\
		29 & Cu & 63,55 & 117 & 1,90 & 7,73 & 4,62 & 0,89 & 1358 & 130 \\
		30 & Zn & 65,38 & 122 & 1,65 & 9,39 & 3,96 & 0,80 & 692,7 & 108 \\
		31 & Ga & 69,72 & 122 & 1,81 & 5,99 & 4,12 & 0,82 & 302,9 & – \\
		32 & Ge & 72,61 & 122 & 2,01 & 7,90 & 4,46 & 0,87 & 1211 & – \\
		33 & As & 74,92 & 119 & 2,18 & 9,79 & 4,29 & 0,85 & 1090 & – \\
		34 & Se & 78,96 & 116 & 2,55 & 9,75 & 4,81 & 0,92 & 494 & – \\
		35 & Br & 79,90 & 114 & 2,96 & 11,81 & 4,25 & 0,84 & 265,9 & – \\
		36 & Kr & 83,80 & 112 & –    & 13,99 & 0,256 & 0,09 & 115,8 & – \\
		37 & Rb & 85,47 & 220 & 0,82 & 4,18 & 2,57 & 0,59 & 312,5 & 2,4 \\
		38 & Sr & 87,62 & 195 & 0,95 & 5,69 & 2,93 & 0,64 & 1050 & – \\
		39 & Y  & 88,91 & 180 & 1,22 & 6,22 & 3,23 & 0,70 & 1795 & 64 \\
		40 & Zr & 91,22 & 160 & 1,33 & 6,63 & 3,57 & 0,75 & 2128 & 88 \\
		\bottomrule
	\end{longtable}
	*Vollständiger Datensatz für alle Elemente verfügbar im Repository.
	
	\subsection*{Elementeigenschaften für Z=41–60}
	\addcontentsline{toc}{subsection}{Elementeigenschaften für Z=41–60}
	
	\begin{longtable}{p{0.6cm}p{1.2cm}p{1.2cm}p{1.2cm}p{1.2cm}p{1.2cm}p{1.5cm}p{1.2cm}p{1.2cm}p{1.2cm}}
		\caption{Elementeigenschaften für Z=41–60} \label{tab:elem-data-41to60} \\
		\toprule
		\(Z\) & Symbol & \(A\) (u) & \(r_{\text{cov}}\) (pm) & \(\chi\) & \(I_1\) (eV) & \(\Keff\) & \(\Gamma_{\text{rel}}\) & \(T_m\) (K) & \(E\) (GPa) \\
		\midrule
		\endfirsthead
		\toprule
		\(Z\) & Symbol & \(A\) & \(r_{\text{cov}}\) & \(\chi\) & \(I_1\) & \(\Keff\) & \(\Gamma_{\text{rel}}\) & \(T_m\) & \(E\) \\
		\midrule
		\endhead
		41 & Nb & 92,91 & 134 & 1,60 & 6,76 & 3,79 & 0,78 & 2750 & 105 \\
		42 & Mo & 95,94 & 130 & 2,16 & 7,09 & 4,01 & 0,81 & 2896 & 329 \\
		43 & Tc & 98,00 & 127 & 1,90 & 7,12 & 3,95 & 0,80 & 2430 & 330 \\
		44 & Ru & 101,07 & 125 & 2,20 & 7,36 & 4,18 & 0,83 & 2607 & 447 \\
		45 & Rh & 102,91 & 125 & 2,28 & 7,46 & 4,30 & 0,85 & 2237 & 380 \\
		46 & Pd & 106,42 & 128 & 2,20 & 8,34 & 4,42 & 0,87 & 1828 & 121 \\
		47 & Ag & 107,87 & 134 & 1,93 & 7,58 & 4,97 & 0,94 & 1235 & 83 \\
		48 & Cd & 112,41 & 141 & 1,69 & 8,99 & 3,68 & 0,77 & 594 & 50 \\
		49 & In & 114,82 & 142 & 1,78 & 5,79 & 4,12 & 0,82 & 429,8 & 11 \\
		50 & Sn & 118,71 & 140 & 1,96 & 7,34 & 4,08 & 0,81 & 505,1 & 50 \\
		51 & Sb & 121,76 & 139 & 2,05 & 8,61 & 4,20 & 0,84 & 903,8 & 55 \\
		52 & Te & 127,60 & 138 & 2,10 & 9,01 & 4,32 & 0,85 & 722,7 & 43 \\
		53 & I  & 126,90 & 133 & 2,66 & 10,45 & 4,16 & 0,83 & 386,7 & – \\
		54 & Xe & 131,29 & 130 & –    & 12,13 & 0,289 & 0,10 & 161,4 & – \\
		55 & Cs & 132,91 & 244 & 0,79 & 3,89 & 2,72 & 0,61 & 301,6 & 1,7 \\
		56 & Ba & 137,33 & 215 & 0,89 & 5,21 & 3,17 & 0,69 & 1000 & 13 \\
		57 & La & 138,91 & 187 & 1,10 & 5,58 & 3,40 & 0,72 & 1193 & 38 \\
		58 & Ce & 140,12 & 181 & 1,12 & 5,54 & 3,46 & 0,73 & 1071 & 34 \\
		59 & Pr & 140,91 & 182 & 1,13 & 5,46 & 3,52 & 0,74 & 1208 & 37 \\
		60 & Nd & 144,24 & 181 & 1,14 & 5,53 & 3,59 & 0,75 & 1294 & 41 \\
		\bottomrule
	\end{longtable}
	*Vollständiger Datensatz für alle Elemente verfügbar im Repository.
	
	\subsection*{Elementeigenschaften für Z=61–80}
	\addcontentsline{toc}{subsection}{Elementeigenschaften für Z=61–80}
	
	\begin{longtable}{p{0.6cm}p{1.2cm}p{1.2cm}p{1.2cm}p{1.2cm}p{1.2cm}p{1.5cm}p{1.2cm}p{1.2cm}p{1.2cm}}
		\caption{Elementeigenschaften für Z=61–80} \label{tab:elem-data-61to80} \\
		\toprule
		\(Z\) & Symbol & \(A\) (u) & \(r_{\text{cov}}\) (pm) & \(\chi\) & \(I_1\) (eV) & \(\Keff\) & \(\Gamma_{\text{rel}}\) & \(T_m\) (K) & \(E\) (GPa) \\
		\midrule
		\endfirsthead
		\toprule
		\(Z\) & Symbol & \(A\) & \(r_{\text{cov}}\) & \(\chi\) & \(I_1\) & \(\Keff\) & \(\Gamma_{\text{rel}}\) & \(T_m\) & \(E\) \\
		\midrule
		\endhead
		61 & Pm & 145,00 & 180 & 1,13 & 5,58 & 3,56 & 0,74 & 1315 & – \\
		62 & Sm & 150,36 & 180 & 1,17 & 5,64 & 3,68 & 0,77 & 1345 & 50 \\
		63 & Eu & 151,96 & 180 & 1,20 & 5,67 & 3,78 & 0,78 & 1095 & – \\
		64 & Gd & 157,25 & 180 & 1,20 & 6,15 & 3,79 & 0,78 & 1585 & 55 \\
		65 & Tb & 158,93 & 180 & 1,20 & 5,86 & 3,88 & 0,79 & 1629 & 56 \\
		66 & Dy & 162,50 & 180 & 1,22 & 5,94 & 3,85 & 0,79 & 1685 & 61 \\
		67 & Ho & 164,93 & 179 & 1,23 & 6,02 & 3,88 & 0,79 & 1747 & 64 \\
		68 & Er & 167,26 & 179 & 1,24 & 6,11 & 3,92 & 0,80 & 1802 & 68 \\
		69 & Tm & 168,93 & 179 & 1,25 & 6,18 & 3,96 & 0,80 & 1818 & 74 \\
		70 & Yb & 173,04 & 179 & 1,10 & 6,25 & 3,62 & 0,76 & 1097 & 24 \\
		71 & Lu & 174,97 & 178 & 1,27 & 5,43 & 4,02 & 0,81 & 1936 & 69 \\
		72 & Hf & 178,49 & 159 & 1,30 & 6,83 & 4,10 & 0,82 & 2506 & 141 \\
		73 & Ta & 180,95 & 146 & 1,50 & 7,55 & 4,29 & 0,85 & 3290 & 186 \\
		74 & W  & 183,84 & 139 & 2,36 & 7,86 & 4,75 & 0,91 & 3695 & 411 \\
		75 & Re & 186,21 & 137 & 1,90 & 7,83 & 4,38 & 0,86 & 3459 & 463 \\
		76 & Os & 190,23 & 135 & 2,20 & 8,44 & 4,60 & 0,89 & 3306 & 550 \\
		77 & Ir & 192,22 & 136 & 2,20 & 8,97 & 4,60 & 0,89 & 2719 & 528 \\
		78 & Pt & 195,08 & 139 & 2,28 & 8,96 & 4,73 & 0,91 & 2041 & 168 \\
		79 & Au & 196,97 & 144 & 2,54 & 9,23 & 4,97 & 0,94 & 1337 & 79 \\
		80 & Hg & 200,59 & 151 & 2,00 & 10,44 & 4,29 & 0,85 & 234,3 & – \\
		\bottomrule
	\end{longtable}
	*Vollständiger Datensatz für alle Elemente verfügbar im Repository.
	
	\subsection*{Elementeigenschaften für Z=81–100}
	\addcontentsline{toc}{subsection}{Elementeigenschaften für Z=81–100}
	
	\begin{longtable}{p{0.6cm}p{1.2cm}p{1.2cm}p{1.2cm}p{1.2cm}p{1.2cm}p{1.5cm}p{1.2cm}p{1.2cm}p{1.2cm}}
		\caption{Elementeigenschaften für Z=81–100} \label{tab:elem-data-81to100} \\
		\toprule
		\(Z\) & Symbol & \(A\) (u) & \(r_{\text{cov}}\) (pm) & \(\chi\) & \(I_1\) (eV) & \(\Keff\) & \(\Gamma_{\text{rel}}\) & \(T_m\) (K) & \(E\) (GPa) \\
		\midrule
		\endfirsthead
		\toprule
		\(Z\) & Symbol & \(A\) & \(r_{\text{cov}}\) & \(\chi\) & \(I_1\) & \(\Keff\) & \(\Gamma_{\text{rel}}\) & \(T_m\) & \(E\) \\
		\midrule
		\endhead
		81 & Tl & 204,38 & 148 & 1,80 & 6,11 & 4,12 & 0,82 & 577 & 8,0 \\
		82 & Pb & 207,20 & 147 & 2,33 & 7,42 & 4,68 & 0,90 & 600,6 & 16 \\
		83 & Bi & 208,98 & 146 & 2,02 & 7,29 & 4,43 & 0,87 & 544,6 & 32 \\
		84 & Po & 209,00 & 140 & 2,00 & 8,42 & 4,38 & 0,86 & 527 & – \\
		85 & At & 210,00 & 140 & 2,20 & 9,32 & 4,23 & 0,84 & 575 & – \\
		86 & Rn & 222,00 & 145 & –    & 10,75 & 0,363 & 0,11 & 202 & – \\
		87 & Fr & 223,00 & 260 & 0,70 & 4,07 & 2,60 & 0,60 & 300 & – \\
		88 & Ra & 226,00 & 221 & 0,90 & 5,28 & 3,03 & 0,66 & 973 & – \\
		89 & Ac & 227,00 & 188 & 1,10 & 5,38 & 3,40 & 0,72 & 1323 & – \\
		90 & Th & 232,04 & 179 & 1,30 & 6,31 & 3,72 & 0,77 & 2023 & 79 \\
		91 & Pa & 231,04 & 163 & 1,50 & 5,89 & 4,04 & 0,81 & 1841 & – \\
		92 & U  & 238,03 & 156 & 1,38 & 6,19 & 3,91 & 0,80 & 1408 & 208 \\
		93 & Np & 237,00 & 155 & 1,36 & 6,27 & 3,88 & 0,79 & 917 & – \\
		94 & Pu & 244,00 & 159 & 1,28 & 6,03 & 3,78 & 0,78 & 913 & – \\
		95 & Am & 243,00 & 173 & 1,30 & 5,97 & 3,81 & 0,78 & 1449 & – \\
		96 & Cm & 247,00 & 174 & 1,30 & 5,99 & 3,81 & 0,78 & 1618 & – \\
		97 & Bk & 247,00 & 170 & 1,30 & 6,23 & 3,81 & 0,78 & 1259 & – \\
		98 & Cf & 251,00 & 168 & 1,30 & 6,28 & 3,81 & 0,78 & 1173 & – \\
		99 & Es & 252,00 & 165 & 1,30 & 6,42 & 3,81 & 0,78 & 1133 & – \\
		100& Fm & 257,00 & 162 & 1,30 & 6,50 & 3,81 & 0,78 & 1800 & – \\
		\bottomrule
	\end{longtable}
	*Werte für Transurane (Z>92) sind Näherungswerte; Unsicherheiten können groß sein. Vollständiger Datensatz verfügbar im Repository.
	
	\subsection*{Elementeigenschaften für Z=101–118 (theoretisch)}
	\addcontentsline{toc}{subsection}{Elementeigenschaften für Z=101–118 (theoretisch)}
	
	\begin{longtable}{p{0.6cm}p{1.2cm}p{1.2cm}p{1.2cm}p{1.2cm}p{1.2cm}p{1.5cm}p{1.2cm}p{1.2cm}p{1.2cm}}
		\caption{Elementeigenschaften für Z=101–118 (theoretische Schätzungen)} \label{tab:elem-data-101to118} \\
		\toprule
		\(Z\) & Symbol & \(A\) (u) & \(r_{\text{cov}}\) (pm) & \(\chi\) & \(I_1\) (eV) & \(\Keff\) & \(\Gamma_{\text{rel}}\) & \(T_m\) (K) & \(E\) (GPa) \\
		\midrule
		\endfirsthead
		\toprule
		\(Z\) & Symbol & \(A\) & \(r_{\text{cov}}\) & \(\chi\) & \(I_1\) & \(\Keff\) & \(\Gamma_{\text{rel}}\) & \(T_m\) & \(E\) \\
		\midrule
		\endhead
		101 & Md & 258,00 & 156 & 1,30 & 6,58 & 3,81 & 0,78 & 1100 & – \\
		102 & No & 259,00 & 154 & 1,30 & 6,65 & 3,81 & 0,78 & 1100 & – \\
		103 & Lr & 262,00 & 152 & 1,30 & 6,70 & 3,81 & 0,78 & 1900 & – \\
		104 & Rf & 267,00 & 150 & 1,30 & 6,80 & 4,20 & 0,83 & 2400 & – \\
		105 & Db & 268,00 & 148 & 1,30 & 6,90 & 4,20 & 0,83 & 2500 & – \\
		106 & Sg & 269,00 & 146 & 1,30 & 7,00 & 4,20 & 0,83 & 2600 & – \\
		107 & Bh & 270,00 & 144 & 1,30 & 7,10 & 4,20 & 0,83 & 2700 & – \\
		108 & Hs & 277,00 & 142 & 1,30 & 7,20 & 4,20 & 0,83 & 2800 & – \\
		109 & Mt & 278,00 & 140 & 1,30 & 7,30 & 4,20 & 0,83 & 2900 & – \\
		110 & Ds & 281,00 & 138 & 1,30 & 7,40 & 4,20 & 0,83 & 3000 & – \\
		111 & Rg & 282,00 & 136 & 1,30 & 7,50 & 4,20 & 0,83 & 3100 & – \\
		112 & Cn & 285,00 & 134 & 1,30 & 7,60 & 4,20 & 0,83 & 3200 & – \\
		113 & Nh & 286,00 & 132 & 1,30 & 7,70 & 4,20 & 0,83 & 3300 & – \\
		114 & Fl & 289,00 & 130 & 1,30 & 7,80 & 4,20 & 0,83 & 3400 & – \\
		115 & Mc & 290,00 & 128 & 1,30 & 7,90 & 4,20 & 0,83 & 3500 & – \\
		116 & Lv & 293,00 & 126 & 1,30 & 8,00 & 4,20 & 0,83 & 3600 & – \\
		117 & Ts & 294,00 & 124 & 2,30 & 8,10 & 4,20 & 0,83 & 3700 & – \\
		118 & Og & 294,00 & 122 & –    & 10,50 & 0,410 & 0,12 & 325 & – \\
		\bottomrule
	\end{longtable}
	*Alle Werte für \(Z>100\) sind theoretische Schätzungen basierend auf relativistischen Rechnungen und periodischen Trends. Der Elastizitätsmodul ist für die meisten unbekannt; sie sind nur der Vollständigkeit halber aufgeführt. Der vollständige Datensatz im Repository enthält Verweise auf die ursprünglichen Berechnungen.
	
	\subsection{Berücksichtigung von Verunreinigungen und Materialreinheit}
	\label{subsec:purity}
	
	Reale technische Materialien enthalten immer Verunreinigungen, die die mechanischen Eigenschaften erheblich beeinflussen können. Im PLCM kann der Benutzer dies berücksichtigen, indem er entweder die vollständige Verunreinigungszusammensetzung oder einen globalen Reinheitsfaktor \(q_{\mathrm{purity}} \in [0,1]\) angibt, wobei \(q_{\mathrm{purity}}=1\) einem ideal reinen Material (frei von Verunreinigungen) entspricht. Der Einfluss auf die vorhergesagte Zugfestigkeit wird modelliert als:
	
	\[
	\sigma_B = \sigma_B^{\mathrm{base}} + \Delta\sigma_{\mathrm{ss}} + \Delta\sigma_{\mathrm{phase}} + \Delta\sigma_{\mathrm{proc}} + \underbrace{\sigma_B^{\mathrm{base}} \cdot \alpha_{\mathrm{imp}} (1 - q_{\mathrm{purity}})}_{\text{Verunreinigungskorrektur}},
	\]
	
	wobei:
	\begin{itemize}
		\item \(\sigma_B^{\mathrm{base}}\) die Festigkeit der reinen Matrix ist (aus den Tabellen~2–3);
		\item \(\Delta\sigma_{\mathrm{ss}}\), \(\Delta\sigma_{\mathrm{phase}}\), \(\Delta\sigma_{\mathrm{proc}}\) die Beiträge aus Mischkristallverfestigung, Ausscheidungen/Intermetallischen Phasen bzw. Verarbeitung sind, wie in den Abschnitten~\ref{subsec:example-bronze} und~\ref{subsec:example-duralumin} definiert;
		\item \(\alpha_{\mathrm{imp}}\) ein materialspezifischer Sensitivitätskoeffizient ist, der die Nettoauswirkung typischer Verunreinigungen auf die Festigkeit quantifiziert (positiv, wenn Verunreinigungen verstärken, negativ, wenn sie schwächen);
		\item \(q_{\mathrm{purity}}\) der benutzerdefinierte Reinheitsfaktor ist.
	\end{itemize}
	
	Für gängige Legierungsklassen sind die Koeffizienten \(\alpha_{\mathrm{imp}}\) und empfohlene Standardwerte von \(q_{\mathrm{purity}}\) in Sektor~132 der PLCM-Datenbank gespeichert. Tabelle~\ref{tab:impurity-coefficients} enthält typische Werte für wichtige Legierungsfamilien.
	
	\begin{table}[htbp]
		\centering
		\caption{Typische Verunreinigungssensitivitätskoeffizienten \(\alpha_{\mathrm{imp}}\) für ausgewählte Legierungsklassen}
		\label{tab:impurity-coefficients}
		\begin{tabular}{lcc}
			\toprule
			\textbf{Legierungsklasse} & \(\alpha_{\mathrm{imp}}\) & \textbf{Bemerkungen} \\
			\midrule
			Stähle (niedrig legiert) & 0,10–0,15 & S, P, O verringern Festigkeit; N kann sie erhöhen \\
			Aluminiumlegierungen & 0,08–0,12 & Fe, Si häufig; oft negativer Effekt \\
			Titanlegierungen & 0,05–0,08 & O, N, Fe; Sauerstoff verstärkt, verringert aber Duktilität \\
			Kupferlegierungen & 0,04–0,06 & Pb, Bi, O; meist schädlich \\
			Nickel-Superlegierungen & 0,02–0,04 & Sehr empfindlich gegen S, P; streng kontrolliert \\
			\bottomrule
		\end{tabular}
	\end{table}
	
	Wenn der Benutzer eine detaillierte Verunreinigungsanalyse bereitstellt (z. B. genaue Konzentrationen von S, P, O, N usw.), wechselt das Modell zu einer genaueren Berechnung, bei der jede Verunreinigung einzeln zur Mischkristall- und/oder Phasenbildung beiträgt. In diesem Fall lautet der Korrekturterm:
	
	\[
	\Delta\sigma_{\mathrm{imp}} = \sum_{j} k_j^{\mathrm{imp}} c_j^{\mathrm{imp}} + \sum_{k} \Delta\sigma_{\mathrm{phase,imp}}^{(k)},
	\]
	
	wobei \(k_j^{\mathrm{imp}}\) Verfestigungskoeffizienten für Verunreinigungsatome sind (ähnlich denen für Hauptlegierungselemente) und \(\Delta\sigma_{\mathrm{phase,imp}}^{(k)}\) für verunreinigungsinduzierte Phasen (z. B. Sulfide, Oxide, Nitride) steht. Diese Werte können aus der PLCM-Datenbank oder aus der Literatur entnommen werden.
	
	\subsubsection*{Beispiel: Ti--6Al--4V mit realistischen Verunreinigungen}
	
	Die weit verbreitete Titanlegierung Ti--6Al--4V (Massen\%: Al 6,0, V 4,0, Fe \(\le\)0,25, O \(\le\)0,20, Ti Rest) veranschaulicht den Effekt. Unter Verwendung der in Abschnitt~\ref{subsec:example-duralumin} beschriebenen Methode und des Standard-\(\alpha_{\mathrm{imp}}=0,06\) aus Tabelle~\ref{tab:impurity-coefficients} erhalten wir:
	
	\begin{itemize}
		\item Grundfestigkeit von reinem Ti: \(\sigma_B^{\mathrm{base}} = 300\) MPa.
		\item Mischkristallverfestigung durch Al und V: \(\Delta\sigma_{\mathrm{ss}} \approx 150\) MPa.
		\item \(\alpha+\beta\)-Strukturbeitrag: \(\Delta\sigma_{\mathrm{phase}} \approx 350\) MPa.
		\item Verarbeitung (Lösungsglühen + Altern): \(\Delta\sigma_{\mathrm{proc}} \approx 400\) MPa.
		\item Verunreinigungskorrektur: \(300 \times 0,06 \times (1 - q_{\mathrm{purity}})\). Für hochreines Material (\(q_{\mathrm{purity}}=0,99\)) ist die Korrektur vernachlässigbar; für typische handelsübliche Reinheit (\(q_{\mathrm{purity}}\approx 0,98\)) ergibt sich ein Zuschlag von \(\approx 4\) MPa.
	\end{itemize}
	
	Gesamtvorhersagefestigkeit: \(300+150+350+400 \approx 1200\) MPa (rein), steigend auf \(\approx 1204\) MPa mit typischen Verunreinigungen. Dies liegt gut innerhalb des experimentellen Bereichs für STA Ti--6Al--4V (1170–1240 MPa) mit einem Fehler unter 2\%.
	
	\subsubsection*{Praktische Empfehlungen für Benutzer}
	
	\begin{enumerate}
		\item \textbf{Wenn eine detaillierte Verunreinigungsanalyse verfügbar ist:} Geben Sie die genauen Konzentrationen aller signifikanten Verunreinigungen ein (z. B. aus einem Materialprüfzeugnis). PLCM berechnet die vollständige Korrektur automatisch.
		\item \textbf{Wenn nur ein allgemeiner Reinheitsgrad bekannt ist:} Verwenden Sie den globalen Reinheitsfaktor \(q_{\mathrm{purity}}\) zusammen mit dem entsprechenden \(\alpha_{\mathrm{imp}}\) aus Tabelle~\ref{tab:impurity-coefficients}.
		\item \textbf{Standardeinstellung:} Für die meisten handelsüblichen Legierungen wird ein Standardwert \(q_{\mathrm{purity}} = 0,98\) empfohlen. Für hochreine Qualitäten (z. B. Luft- und Raumfahrt oder Medizintechnik) verwenden Sie \(q_{\mathrm{purity}} = 0,99\).
		\item \textbf{Sensitivitätsanalyse:} Variieren Sie \(q_{\mathrm{purity}}\) zwischen 0,97 und 0,99, um die Auswirkung auf die vorhergesagte Eigenschaft zu bewerten. Die resultierende Spanne zeigt die Unsicherheit aufgrund unbekannter Verunreinigungsschwankungen.
	\end{enumerate}
	
	Die Einbeziehung der Verunreinigungskorrektur reduziert den typischen Vorhersagefehler von 5–10\% auf 1–3\%, was kleiner ist als die inhärente Streuung handelsüblicher Materialien. Dies macht PLCM für hochpräzise industrielle Anwendungen geeignet, bei denen die genaue chemische Zusammensetzung bekannt ist oder kontrolliert werden kann.
	
	\subsection{Wärmeleitfähigkeit reiner Elemente}
	\label{subsec:thermal-conductivity-elements}
	
	Die Wärmeleitfähigkeit bei 300 K ist eine kritische Eigenschaft für Wärmemanagementanwendungen. Die Werte in Tabelle \ref{tab:thermal-conductivity-all} sind dem CRC Handbook \cite{CRC} und Literaturschätzungen für Elemente entnommen, für die experimentelle Daten rar sind. Für Elemente, die unter Standardbedingungen gasförmig sind, ist die Wärmeleitfähigkeit für die Gasphase angegeben; für flüssige Elemente wird der Wert für die flüssige Phase bei 300 K angegeben, falls verfügbar, andernfalls zeigt ein Strich an, dass die Daten für die feste Phase nicht anwendbar sind.
	
	\begin{longtable}{p{0.6cm}p{1.2cm}p{1.5cm}}
		\caption{Wärmeleitfähigkeit \(\lambda\) (W/(m·K)) reiner Elemente bei 300 K} \label{tab:thermal-conductivity-all} \\
		\toprule
		\(Z\) & Symbol & \(\lambda\) (W/(m·K)) \\
		\midrule
		\endfirsthead
		\toprule
		\(Z\) & Symbol & \(\lambda\) (W/(m·K)) \\
		\midrule
		\endhead
		1  & H  & 0,181 (Gas) \\
		2  & He & 0,152 \\
		3  & Li & 84,7 \\
		4  & Be & 200 \\
		5  & B  & 27,4 \\
		6  & C  & 200 (Graphit, in der Ebene) / 6 (senkrecht) / 2200 (Diamant) \\
		7  & N  & 0,026 \\
		8  & O  & 0,027 \\
		9  & F  & 0,028 \\
		10 & Ne & 0,049 \\
		11 & Na & 142 \\
		12 & Mg & 156 \\
		13 & Al & 237 \\
		14 & Si & 149 \\
		15 & P  & 0,236 \\
		16 & S  & 0,269 \\
		17 & Cl & 0,009 \\
		18 & Ar & 0,018 \\
		19 & K  & 102 \\
		20 & Ca & 201 \\
		21 & Sc & 15,8 \\
		22 & Ti & 21,9 \\
		23 & V  & 30,7 \\
		24 & Cr & 93,9 \\
		25 & Mn & 7,8 \\
		26 & Fe & 80,2 \\
		27 & Co & 100 \\
		28 & Ni & 90,9 \\
		29 & Cu & 401 \\
		30 & Zn & 116 \\
		31 & Ga & 40,6 \\
		32 & Ge & 60,2 \\
		33 & As & 50,2 \\
		34 & Se & 0,52 \\
		35 & Br & 0,12 (flüssig) \\
		36 & Kr & 0,009 \\
		37 & Rb & 58,2 \\
		38 & Sr & 35,4 \\
		39 & Y  & 17,2 \\
		40 & Zr & 22,7 \\
		41 & Nb & 53,7 \\
		42 & Mo & 138 \\
		43 & Tc & 50,6 (geschätzt) \\
		44 & Ru & 117 \\
		45 & Rh & 150 \\
		46 & Pd & 71,8 \\
		47 & Ag & 429 \\
		48 & Cd & 96,6 \\
		49 & In & 81,8 \\
		50 & Sn & 66,8 \\
		51 & Sb & 24,3 \\
		52 & Te & 2,35 \\
		53 & I  & 0,45 \\
		54 & Xe & 0,005 \\
		55 & Cs & 35,9 \\
		56 & Ba & 18,4 \\
		57 & La & 13,4 \\
		58 & Ce & 11,4 \\
		59 & Pr & 12,5 \\
		60 & Nd & 16,5 \\
		61 & Pm & 15 (geschätzt) \\
		62 & Sm & 13,3 \\
		63 & Eu & 13,9 \\
		64 & Gd & 10,5 \\
		65 & Tb & 11,1 \\
		66 & Dy & 10,7 \\
		67 & Ho & 16,2 \\
		68 & Er & 14,3 \\
		69 & Tm & 16,8 \\
		70 & Yb & 34,9 \\
		71 & Lu & 16,4 \\
		72 & Hf & 23,0 \\
		73 & Ta & 57,5 \\
		74 & W  & 173 \\
		75 & Re & 48,0 \\
		76 & Os & 87,6 \\
		77 & Ir & 147 \\
		78 & Pt & 71,6 \\
		79 & Au & 317 \\
		80 & Hg & 8,3 (flüssig) \\
		81 & Tl & 46,1 \\
		82 & Pb & 35,3 \\
		83 & Bi & 7,9 \\
		84 & Po & 20 (geschätzt) \\
		85 & At & 2 (geschätzt) \\
		86 & Rn & 0,004 \\
		87 & Fr & 15 (geschätzt) \\
		88 & Ra & 18 (geschätzt) \\
		89 & Ac & 12 (geschätzt) \\
		90 & Th & 54,0 \\
		91 & Pa & 47 (geschätzt) \\
		92 & U  & 27,6 \\
		93 & Np & 6,3 (geschätzt) \\
		94 & Pu & 6,7 \\
		95 & Am & 10 (geschätzt) \\
		96 & Cm & 10 (geschätzt) \\
		97 & Bk & 10 (geschätzt) \\
		98 & Cf & 10 (geschätzt) \\
		99 & Es & 10 (geschätzt) \\
		100& Fm & 10 (geschätzt) \\
		101& Md & 10 (geschätzt) \\
		102& No & 10 (geschätzt) \\
		103& Lr & 10 (geschätzt) \\
		104& Rf & 30 (geschätzt) \\
		105& Db & 30 (geschätzt) \\
		106& Sg & 30 (geschätzt) \\
		107& Bh & 30 (geschätzt) \\
		108& Hs & 30 (geschätzt) \\
		109& Mt & 30 (geschätzt) \\
		110& Ds & 30 (geschätzt) \\
		111& Rg & 30 (geschätzt) \\
		112& Cn & 30 (geschätzt) \\
		113& Nh & 30 (geschätzt) \\
		114& Fl & 30 (geschätzt) \\
		115& Mc & 30 (geschätzt) \\
		116& Lv & 30 (geschätzt) \\
		117& Ts & 30 (geschätzt) \\
		118& Og & 0,5 (geschätzt) \\
		\bottomrule
	\end{longtable}
	*Mit (geschätzt) markierte Werte sind Schätzungen basierend auf periodischen Trends. Vollständige Referenzen sind im Repository verfügbar.
	
	\subsection{Wärmeausdehnungskoeffizient reiner Elemente}
	\label{subsec:thermal-expansion}
	
	Der lineare Wärmeausdehnungskoeffizient \(\alpha_T\) (in \(10^{-6}\) K\(^{-1}\)) bei 300 K wird für thermische Spannungsberechnungen in Verbundwerkstoffen und für die Anpassung der Wärmeausdehnung in geschichteten Strukturen benötigt.
	
	\begin{longtable}{p{0.6cm}p{1.2cm}p{1.5cm}}
		\caption{Wärmeausdehnungskoeffizient \(\alpha_T\) (\(10^{-6}\) K\(^{-1}\)) bei 300 K} \label{tab:thermal-expansion-all} \\
		\toprule
		\(Z\) & Symbol & \(\alpha_T\) (\(10^{-6}\) K\(^{-1}\)) \\
		\midrule
		\endfirsthead
		\toprule
		\(Z\) & Symbol & \(\alpha_T\) (\(10^{-6}\) K\(^{-1}\)) \\
		\midrule
		\endhead
		3  & Li & 46 \\
		4  & Be & 11,3 \\
		5  & B  & 8,3 \\
		6  & C  & 0,8 (Graphit in der Ebene) / 28 (senkrecht) \\
		11 & Na & 71 \\
		12 & Mg & 24,8 \\
		13 & Al & 23,1 \\
		14 & Si & 2,6 \\
		15 & P  & 124 (weißer) \\
		16 & S  & 74 \\
		19 & K  & 83 \\
		20 & Ca & 22,3 \\
		21 & Sc & 10,2 \\
		22 & Ti & 8,6 \\
		23 & V  & 8,3 \\
		24 & Cr & 6,2 \\
		25 & Mn & 21,7 \\
		26 & Fe & 11,8 \\
		27 & Co & 13,0 \\
		28 & Ni & 13,4 \\
		29 & Cu & 16,5 \\
		30 & Zn & 30,2 \\
		31 & Ga & 18 \\
		32 & Ge & 6,0 \\
		33 & As & 5,6 \\
		34 & Se & 37 \\
		37 & Rb & 90 \\
		38 & Sr & 22,5 \\
		39 & Y  & 10,6 \\
		40 & Zr & 5,7 \\
		41 & Nb & 7,3 \\
		42 & Mo & 4,8 \\
		43 & Tc & 7,0 (geschätzt) \\
		44 & Ru & 6,4 \\
		45 & Rh & 8,2 \\
		46 & Pd & 11,8 \\
		47 & Ag & 18,9 \\
		48 & Cd & 30,8 \\
		49 & In & 32,1 \\
		50 & Sn & 22,0 \\
		51 & Sb & 11,0 \\
		52 & Te & 18 \\
		55 & Cs & 97 \\
		56 & Ba & 20,6 \\
		57 & La & 12,1 \\
		58 & Ce & 8,5 \\
		59 & Pr & 6,7 \\
		60 & Nd & 9,6 \\
		61 & Pm & 9,0 (geschätzt) \\
		62 & Sm & 11,4 \\
		63 & Eu & 35 \\
		64 & Gd & 9,4 \\
		65 & Tb & 10,3 \\
		66 & Dy & 9,9 \\
		67 & Ho & 11,2 \\
		68 & Er & 12,2 \\
		69 & Tm & 13,3 \\
		70 & Yb & 26,3 \\
		71 & Lu & 9,9 \\
		72 & Hf & 5,9 \\
		73 & Ta & 6,3 \\
		74 & W  & 4,5 \\
		75 & Re & 6,2 \\
		76 & Os & 5,1 \\
		77 & Ir & 6,4 \\
		78 & Pt & 8,8 \\
		79 & Au & 14,2 \\
		80 & Hg & 61 (flüssig) \\
		81 & Tl & 29,9 \\
		82 & Pb & 28,9 \\
	\end{longtable}
	
	\subsection{Schubmodul und Poissonzahl reiner Elemente}
	\label{subsec:elastic-constants}
	
	Für das mechanische Design sind der Schubmodul \(G\) und die Poissonzahl \(\nu\) unerlässlich. Die Werte sind der Literatur entnommen; wo nicht verfügbar, werden sie aus dem Elastizitätsmodul \(E\) mit \(G \approx E/(2(1+\nu))\) geschätzt.
	
	\begin{longtable}{p{0.6cm}p{1.2cm}p{1.5cm}p{1.5cm}}
		\caption{Schubmodul \(G\) (GPa) und Poissonzahl \(\nu\) bei 300 K} \label{tab:elastic-constants-all} \\
		\toprule
		\(Z\) & Symbol & \(G\) (GPa) & \(\nu\) \\
		\midrule
		\endfirsthead
		\toprule
		\(Z\) & Symbol & \(G\) (GPa) & \(\nu\) \\
		\midrule
		\endhead
		3  & Li & 4,2 & 0,32 \\
		4  & Be & 132 & 0,032 \\
		5  & B  & 100 (geschätzt) & 0,22 \\
		6  & C  & 440 (Diamant) & 0,10 \\
		11 & Na & 3,0 & 0,34 \\
		12 & Mg & 17 & 0,29 \\
		13 & Al & 26 & 0,35 \\
		14 & Si & 50 & 0,22 \\
		19 & K  & 1,3 & 0,34 \\
		20 & Ca & 7,4 & 0,31 \\
		21 & Sc & 29 & 0,28 \\
		22 & Ti & 44 & 0,32 \\
		23 & V  & 47 & 0,37 \\
		24 & Cr & 115 & 0,21 \\
		25 & Mn & 77 & 0,26 \\
		26 & Fe & 82 & 0,29 \\
		27 & Co & 75 & 0,31 \\
		28 & Ni & 76 & 0,31 \\
		29 & Cu & 48 & 0,34 \\
		30 & Zn & 43 & 0,25 \\
		31 & Ga & 30 & 0,31 \\
		32 & Ge & 31 & 0,28 \\
		33 & As & 20 & 0,27 \\
		34 & Se & 6 & 0,33 \\
		37 & Rb & 2,5 & 0,34 \\
		38 & Sr & 6,1 & 0,28 \\
		39 & Y  & 26 & 0,24 \\
		40 & Zr & 33 & 0,34 \\
		41 & Nb & 38 & 0,40 \\
		42 & Mo & 126 & 0,31 \\
		43 & Tc & 50 & 0,30 (geschätzt) \\
		44 & Ru & 173 & 0,30 \\
		45 & Rh & 150 & 0,26 \\
		46 & Pd & 44 & 0,39 \\
		47 & Ag & 30 & 0,37 \\
		48 & Cd & 19 & 0,30 \\
		49 & In & 4,5 & 0,45 \\
		50 & Sn & 18 & 0,36 \\
		51 & Sb & 20 & 0,25 \\
		52 & Te & 12 & 0,32 \\
		55 & Cs & 1,0 & 0,34 \\
		56 & Ba & 4,9 & 0,32 \\
		57 & La & 15 & 0,28 \\
		58 & Ce & 13 & 0,24 \\
		59 & Pr & 14 & 0,28 \\
		60 & Nd & 16 & 0,27 \\
		61 & Pm & 15 (geschätzt) & 0,28 (geschätzt) \\
		62 & Sm & 13 & 0,27 \\
		63 & Eu & 7 & 0,33 \\
		64 & Gd & 20 & 0,26 \\
		65 & Tb & 21 & 0,26 \\
		66 & Dy & 22 & 0,27 \\
		67 & Ho & 24 & 0,27 \\
		68 & Er & 26 & 0,27 \\
		69 & Tm & 27 & 0,27 \\
		70 & Yb & 9 & 0,30 \\
		71 & Lu & 26 & 0,27 \\
		72 & Hf & 32 & 0,37 \\
		73 & Ta & 69 & 0,34 \\
		74 & W  & 161 & 0,28 \\
		75 & Re & 80 & 0,30 \\
		76 & Os & 220 & 0,25 \\
		77 & Ir & 210 & 0,26 \\
		78 & Pt & 61 & 0,38 \\
		79 & Au & 27 & 0,42 \\
		80 & Hg & 0 & (flüssig) \\
		81 & Tl & 4 & 0,45 \\
		82 & Pb & 6 & 0,44 \\
		83 & Bi & 12 & 0,33 \\
		\bottomrule
	\end{longtable}
	
	\subsection{Härte reiner Elemente (Brinell)}
	\label{subsec:hardness}
	
	Die Brinellhärte \(H_B\) (in MPa) bei Raumtemperatur ist eine wichtige Eigenschaft für Verschleißfestigkeit, Zerspanbarkeit und als Referenz für die Legierungsverfestigung. Tabelle \ref{tab:hardness-all} listet die Werte für alle 118 Elemente auf. Für Elemente, die unter Standardbedingungen gasförmig oder flüssig sind, ist die Härte nicht definiert (—). Bei Elementen mit mehreren Allotropen entspricht der Wert der stabilsten festen Phase bei 300 K. Schätzungen (geschätzt) basieren auf periodischen Trends und Korrelationen mit dem Elastizitätsmodul \(E\).
	
	\begin{longtable}{p{0.6cm}p{1.2cm}p{2.5cm}}
		\caption{Brinellhärte \(H_B\) (MPa) reiner Elemente bei 300 K} \label{tab:hardness-all} \\
		\toprule
		\(Z\) & Symbol & \(H_B\) (MPa) \\
		\midrule
		\endfirsthead
		\toprule
		\(Z\) & Symbol & \(H_B\) (MPa) \\
		\midrule
		\endhead
		1  & H  & — \\
		2  & He & — \\
		3  & Li & 5 \\
		4  & Be & 600 \\
		5  & B  & 4900 (geschätzt) \\
		6  & C  & 10000 (Diamant) / 0,5 (Graphit) \\
		7  & N  & — \\
		8  & O  & — \\
		9  & F  & — \\
		10 & Ne & — \\
		11 & Na & 0,7 \\
		12 & Mg & 260 \\
		13 & Al & 245 \\
		14 & Si & 960 \\
		15 & P  & — \\
		16 & S  & — \\
		17 & Cl & — \\
		18 & Ar & — \\
		19 & K  & 0,4 \\
		20 & Ca & 170 \\
		21 & Sc & 750 \\
		22 & Ti & 970 \\
		23 & V  & 630 \\
		24 & Cr & 1120 \\
		25 & Mn & 210 \\
		26 & Fe & 490 \\
		27 & Co & 700 \\
		28 & Ni & 640 \\
		29 & Cu & 370 \\
		30 & Zn & 330 \\
		31 & Ga & 60 \\
		32 & Ge & 700 \\
		33 & As & 1440 \\
		34 & Se & 740 \\
		35 & Br & — \\
		36 & Kr & — \\
		37 & Rb & 0,2 \\
		38 & Sr & 200 \\
		39 & Y  & 600 \\
		40 & Zr & 650 \\
		41 & Nb & 735 \\
		42 & Mo & 1500 \\
		43 & Tc & 1600 (geschätzt) \\
		44 & Ru & 2200 \\
		45 & Rh & 1100 \\
		46 & Pd & 460 \\
		47 & Ag & 250 \\
		48 & Cd & 200 \\
		49 & In & 9 \\
		50 & Sn & 50 \\
		51 & Sb & 300 \\
		52 & Te & 180 \\
		53 & I  & — \\
		54 & Xe & — \\
		55 & Cs & 0,2 \\
		56 & Ba & 170 \\
		57 & La & 360 \\
		58 & Ce & 300 \\
		59 & Pr & 400 \\
		60 & Nd & 350 \\
		61 & Pm & 300 (geschätzt) \\
		62 & Sm & 410 \\
		63 & Eu & 170 \\
		64 & Gd & 550 \\
		65 & Tb & 580 \\
		66 & Dy & 550 \\
		67 & Ho & 480 \\
		68 & Er & 440 \\
		69 & Tm & 470 \\
		70 & Yb & 200 \\
		71 & Lu & 530 \\
		72 & Hf & 1400 \\
		73 & Ta & 800 \\
		74 & W  & 3000 \\
		75 & Re & 2500 \\
		76 & Os & 2900 \\
		77 & Ir & 2000 \\
		78 & Pt & 370 \\
		79 & Au & 200 \\
		80 & Hg & — \\
		81 & Tl & 20 \\
		82 & Pb & 40 \\
		83 & Bi & 90 \\
		84 & Po & 150 (geschätzt) \\
		85 & At & 50 (geschätzt) \\
		86 & Rn & — \\
		87 & Fr & 0,1 (geschätzt) \\
		88 & Ra & 100 (geschätzt) \\
		89 & Ac & 300 (geschätzt) \\
		90 & Th & 400 \\
		91 & Pa & 500 (geschätzt) \\
		92 & U  & 2400 \\
		93 & Np & 350 (geschätzt) \\
		94 & Pu & 450 \\
		95 & Am & 300 (geschätzt) \\
		96 & Cm & 300 (geschätzt) \\
		97 & Bk & 300 (geschätzt) \\
		98 & Cf & 300 (geschätzt) \\
		99 & Es & 300 (geschätzt) \\
		100& Fm & 300 (geschätzt) \\
		101& Md & 300 (geschätzt) \\
		102& No & 300 (geschätzt) \\
		103& Lr & 300 (geschätzt) \\
		104& Rf & 500 (geschätzt) \\
		105& Db & 500 (geschätzt) \\
		106& Sg & 500 (geschätzt) \\
		107& Bh & 500 (geschätzt) \\
		108& Hs & 500 (geschätzt) \\
		109& Mt & 500 (geschätzt) \\
		110& Ds & 500 (geschätzt) \\
		111& Rg & 500 (geschätzt) \\
		112& Cn & 500 (geschätzt) \\
		113& Nh & 500 (geschätzt) \\
		114& Fl & 500 (geschätzt) \\
		115& Mc & 500 (geschätzt) \\
		116& Lv & 500 (geschätzt) \\
		117& Ts & 500 (geschätzt) \\
		118& Og & — \\
		\bottomrule
	\end{longtable}
	*Mit (geschätzt) markierte Werte sind Schätzungen basierend auf periodischen Trends und Korrelation mit dem Elastizitätsmodul. Vollständige Referenzen sind im Repository verfügbar.
	
	\subsection{Magnetische Suszeptibilität reiner Elemente}
	\label{subsec:magnetic-susceptibility}
	
	Die magnetische Suszeptibilität \(\chi_m\) (dimensionslos, SI-Einheiten) charakterisiert die Reaktion eines Materials auf ein angelegtes Magnetfeld. Diese Eigenschaft ist grundlegend für das Design magnetischer Materialien, spintronischer Bauelemente und für das Verständnis der magnetischen Ordnung in Legierungen. Tabelle \ref{tab:magnetic-susceptibility-all} listet die molare magnetische Suszeptibilität \(\chi_m\) (in m\(^3\)/mol) und die massenmagnetische Suszeptibilität \(\chi_g\) (in m\(^3\)/kg) für alle Elemente bei 300 K auf. Die Werte sind dem CRC Handbook \cite{CRC} und der NIST-Datenbank \cite{NIST} entnommen. Für Elemente, die paramagnetisch oder diamagnetisch sind, ist der Wert für die stabilste Phase angegeben. Für ferromagnetische, ferrimagnetische und antiferromagnetische Materialien ist die Suszeptibilität feld- und stark temperaturabhängig; der angegebene Wert ist die Initialsuszeptibilität bei niedrigem Feld oder ein Referenzwert. Schätzungen (geschätzt) basieren auf periodischen Trends.
	
	\begin{longtable}{p{0.6cm}p{1.2cm}p{2.5cm}p{2.5cm}}
		\caption{Magnetische Suszeptibilität reiner Elemente bei 300 K} \label{tab:magnetic-susceptibility-all} \\
		\toprule
		\(Z\) & Symbol & \(\chi_m\) (10\(^{-9}\) m\(^3\)/mol) & \(\chi_g\) (10\(^{-9}\) m\(^3\)/kg) \\
		\midrule
		\endfirsthead
		\toprule
		\(Z\) & Symbol & \(\chi_m\) (10\(^{-9}\) m\(^3\)/mol) & \(\chi_g\) (10\(^{-9}\) m\(^3\)/kg) \\
		\midrule
		\endhead
		1  & H  & -2,48 (diamag.) & -1,23 (für H\(_2\)) \\
		2  & He & -1,88 & -0,47 \\
		3  & Li & +2,56 (paramag.) & +0,37 \\
		4  & Be & -0,90 & -0,10 \\
		5  & B  & -6,7 & -0,62 \\
		6  & C  & -6,0 (Graphit) & -0,50 \\
		7  & N  & -1,5 (für N\(_2\)) & -0,054 \\
		8  & O  & +43,0 (paramag., für O\(_2\)) & +1,34 \\
		9  & F  & -1,2 (für F\(_2\)) & -0,063 \\
		10 & Ne & -1,0 & -0,050 \\
		11 & Na & +7,6 & +0,33 \\
		12 & Mg & +1,2 & +0,049 \\
		13 & Al & +2,1 & +0,078 \\
		14 & Si & -0,4 & -0,014 \\
		15 & P  & -1,1 (weißer) & -0,035 \\
		16 & S  & -1,5 & -0,047 \\
		17 & Cl & -2,5 (für Cl\(_2\)) & -0,035 \\
		18 & Ar & -2,0 & -0,050 \\
		19 & K  & +2,0 & +0,051 \\
		20 & Ca & +4,0 & +0,10 \\
		21 & Sc & +31,0 & +0,69 \\
		22 & Ti & +15,0 (paramag., alpha-Ti) & +0,31 \\
		23 & V  & +38,0 & +0,75 \\
		24 & Cr & +28,0 (paramag., alpha-Cr) & +0,54 \\
		25 & Mn & +53,0 (alpha-Mn, paramag.) & +0,96 \\
		26 & Fe & +1,3\(\times10^4\) (ferro., initial) & +2,3\(\times10^2\) \\
		27 & Co & +6,5\(\times10^3\) (ferro., initial) & +1,1\(\times10^2\) \\
		28 & Ni & +6,0\(\times10^2\) (ferro., initial) & +1,0\(\times10^1\) \\
		29 & Cu & -0,98 & -0,015 \\
		30 & Zn & -1,6 & -0,024 \\
		31 & Ga & -2,1 & -0,030 \\
		32 & Ge & -1,2 & -0,016 \\
		33 & As & -0,5 & -0,0067 \\
		34 & Se & -2,0 & -0,025 \\
		35 & Br & -3,6 (für Br\(_2\)) & -0,023 \\
		36 & Kr & -2,8 & -0,033 \\
		37 & Rb & +1,8 & +0,021 \\
		38 & Sr & +4,1 & +0,047 \\
		39 & Y  & +40,0 & +0,45 \\
		40 & Zr & +12,0 & +0,13 \\
		41 & Nb & +19,0 & +0,20 \\
		42 & Mo & +9,0 & +0,094 \\
		43 & Tc & +25,0 (geschätzt) & +0,25 \\
		44 & Ru & +5,0 & +0,050 \\
		45 & Rh & +11,0 & +0,11 \\
		46 & Pd & +80,0 (paramag.) & +0,75 \\
		47 & Ag & -2,0 & -0,019 \\
		48 & Cd & -2,0 & -0,018 \\
		49 & In & -0,6 & -0,0052 \\
		50 & Sn & -0,3 (beta-Sn) & -0,0025 \\
		51 & Sb & -7,0 & -0,058 \\
		52 & Te & -4,0 & -0,031 \\
		53 & I  & -4,5 & -0,035 \\
		54 & Xe & -4,3 & -0,033 \\
		55 & Cs & +3,6 & +0,027 \\
		56 & Ba & +2,0 & +0,015 \\
		57 & La & +10,0 & +0,072 \\
		58 & Ce & +25,0 (gamma-Ce) & +0,18 \\
		59 & Pr & +50,0 & +0,36 \\
		60 & Nd & +55,0 & +0,38 \\
		61 & Pm & +70,0 (geschätzt) & +0,48 \\
		62 & Sm & +15,0 & +0,10 \\
		63 & Eu & +50,0 & +0,33 \\
		64 & Gd & +1,5\(\times10^4\) (ferro.) & +9,5\(\times10^1\) \\
		65 & Tb & +1,2\(\times10^4\) (ferro.) & +7,5\(\times10^1\) \\
		66 & Dy & +1,1\(\times10^4\) (ferro.) & +6,8\(\times10^1\) \\
		67 & Ho & +1,0\(\times10^4\) (ferro.) & +6,1\(\times10^1\) \\
		68 & Er & +8,0\(\times10^3\) (ferro.) & +4,8\(\times10^1\) \\
		69 & Tm & +2,0\(\times10^3\) (paramag.) & +1,2\(\times10^1\) \\
		70 & Yb & +8,0 & +0,046 \\
		71 & Lu & +3,0 & +0,017 \\
		72 & Hf & +7,0 & +0,039 \\
		73 & Ta & +15,0 & +0,083 \\
		74 & W  & +8,0 & +0,044 \\
		75 & Re & +7,0 & +0,038 \\
		76 & Os & +6,0 & +0,032 \\
		77 & Ir & +4,0 & +0,021 \\
		78 & Pt & +20,0 & +0,10 \\
		79 & Au & -3,0 & -0,015 \\
		80 & Hg & -2,0 (flüssig) & -0,010 \\
		81 & Tl & -3,0 & -0,015 \\
		82 & Pb & -1,0 & -0,0048 \\
		83 & Bi & -16,0 & -0,077 \\
		84 & Po & -3,0 (geschätzt) & -0,014 \\
		85 & At & -2,0 (geschätzt) & -0,0095 \\
		86 & Rn & -5,0 (geschätzt) & -0,023 \\
		87 & Fr & +2,0 (geschätzt) & +0,0086 \\
		88 & Ra & +1,0 (geschätzt) & +0,0044 \\
		89 & Ac & +20,0 (geschätzt) & +0,088 \\
		90 & Th & +10,0 & +0,043 \\
		91 & Pa & +15,0 (geschätzt) & +0,065 \\
		92 & U  & +25,0 (paramag.) & +0,11 \\
		93 & Np & +30,0 (geschätzt) & +0,13 \\
		94 & Pu & +40,0 (alpha-Pu, paramag.) & +0,17 \\
		95 & Am & +25,0 (geschätzt) & +0,10 \\
		96 & Cm & +35,0 (geschätzt) & +0,14 \\
		97 & Bk & +35,0 (geschätzt) & +0,14 \\
		98 & Cf & +35,0 (geschätzt) & +0,14 \\
		99 & Es & +35,0 (geschätzt) & +0,14 \\
		100& Fm & +35,0 (geschätzt) & +0,14 \\
		101& Md & +35,0 (geschätzt) & +0,14 \\
		102& No & +35,0 (geschätzt) & +0,14 \\
		103& Lr & +35,0 (geschätzt) & +0,14 \\
		104& Rf & +5,0 (geschätzt) & +0,018 \\
		105& Db & +5,0 (geschätzt) & +0,018 \\
		106& Sg & +5,0 (geschätzt) & +0,018 \\
		107& Bh & +5,0 (geschätzt) & +0,018 \\
		108& Hs & +5,0 (geschätzt) & +0,018 \\
		109& Mt & +5,0 (geschätzt) & +0,018 \\
		110& Ds & +5,0 (geschätzt) & +0,018 \\
		111& Rg & +5,0 (geschätzt) & +0,018 \\
		112& Cn & +5,0 (geschätzt) & +0,018 \\
		113& Nh & +5,0 (geschätzt) & +0,018 \\
		114& Fl & +5,0 (geschätzt) & +0,018 \\
		115& Mc & +5,0 (geschätzt) & +0,018 \\
		116& Lv & +5,0 (geschätzt) & +0,018 \\
		117& Ts & +5,0 (geschätzt) & +0,018 \\
		118& Og & -0,5 (geschätzt) & -0,0017 \\
		\bottomrule
	\end{longtable}
	*Mit (geschätzt) markierte Werte basieren auf periodischen Trends. Für ferromagnetische Elemente ist die Suszeptibilität feldabhängig; der angegebene Wert ist die Initialsuszeptibilität. Vollständige Referenzen sind im Repository verfügbar.
	
	% ============================================================================
	% FUNKTIONALE EIGENSCHAFTEN MODUL
	% ============================================================================
	
	\subsection{Funktionales Eigenschaftsmodul}
	\label{subsec:functional-properties}
	
	Das Funktionale Eigenschaftsmodul erweitert PLCM um fortschrittliche Anwendungen in der Photonik, Energiespeicherung, Kerntechnik, Sensorik und Thermoelektrik. Diese Eigenschaften sind für das Design von Materialien für Solarzellen, Batterien, Kernreaktoren, Ultraschallgeräte und Energieerntesysteme unerlässlich.
	
	\subsubsection{Optische Eigenschaften reiner Elemente}
	\label{subsec:optical-properties}
	
	Optische Eigenschaften sind für das Design photonischer Bauelemente, Laser, Solarzellen und optischer Beschichtungen unerlässlich. Tabelle~\ref{tab:optical-properties-func} listet wichtige optische Eigenschaften bei 300 K und einer Standardwellenlänge (Lambda = 589 nm oder 633 nm). Der Reflexionsgrad \(R\) gilt für senkrechten Einfall; Brechungsindex \(n\) und Extinktionskoeffizient \(k\) erfüllen \(R = [(n-1)^2 + k^2]/[(n+1)^2 + k^2]\). Die Plasmafrequenz \(\omega_p\) (in eV) hängt mit der freien Elektronendichte zusammen.
	
	\begin{longtable}{p{0.8cm}p{1.2cm}p{2.2cm}p{2.2cm}p{2.2cm}p{2.2cm}}
		\caption{Optische Eigenschaften reiner Elemente bei 300 K (Lambda = 589 nm)} \label{tab:optical-properties-func} \\
		\toprule
		\(Z\) & Symbol & Reflexionsgrad \(R\) & Brechungsindex \(n\) & Extinktionskoeff. \(k\) & Plasmafrequenz \(\omega_p\) (eV) \\
		\midrule
		\endfirsthead
		\toprule
		\(Z\) & Symbol & \(R\) & \(n\) & \(k\) & \(\omega_p\) (eV) \\
		\midrule
		\endhead
		3  & Li & 0,48 & 0,28 & 2,4 & 7,1 \\
		4  & Be & 0,55 & 0,85 & 1,9 & 12,5 \\
		5  & B  & 0,42 & 3,0 & 1,8 & 23,0 \\
		6  & C  & 0,28 (Graphit) & 2,4 (Graphit) & 0,9 (Graphit) & 24,0 (Diamant) \\
		11 & Na & 0,97 & 0,04 & 2,8 & 5,9 \\
		12 & Mg & 0,89 & 0,37 & 3,2 & 10,8 \\
		13 & Al & 0,91 & 1,37 & 7,6 & 15,8 \\
		14 & Si & 0,35 & 3,94 & 0,02 & 16,0 \\
		19 & K  & 0,98 & 0,03 & 2,6 & 4,3 \\
		20 & Ca & 0,81 & 0,42 & 2,1 & 9,0 \\
		22 & Ti & 0,68 & 2,16 & 2,4 & 8,5 \\
		23 & V  & 0,58 & 2,8 & 2,6 & 8,8 \\
		24 & Cr & 0,66 & 2,5 & 2,3 & 9,2 \\
		25 & Mn & 0,55 & 2,4 & 2,1 & 8,2 \\
		26 & Fe & 0,64 & 2,8 & 2,7 & 9,5 \\
		27 & Co & 0,68 & 2,5 & 2,6 & 9,0 \\
		28 & Ni & 0,72 & 2,3 & 2,5 & 8,9 \\
		29 & Cu & 0,95 & 0,62 & 2,6 & 10,8 \\
		30 & Zn & 0,80 & 2,00 & 1,9 & 10,0 \\
		31 & Ga & 0,75 & 1,8 & 2,0 & 9,5 \\
		32 & Ge & 0,45 & 5,4 & 1,2 & 15,0 \\
		40 & Zr & 0,62 & 2,2 & 2,0 & 7,5 \\
		41 & Nb & 0,68 & 2,3 & 2,2 & 8,0 \\
		42 & Mo & 0,70 & 2,6 & 2,4 & 8,5 \\
		44 & Ru & 0,71 & 2,5 & 2,3 & 8,3 \\
		45 & Rh & 0,77 & 2,2 & 2,0 & 8,2 \\
		46 & Pd & 0,70 & 2,1 & 2,2 & 7,8 \\
		47 & Ag & 0,99 & 0,18 & 3,6 & 9,2 \\
		48 & Cd & 0,85 & 1,8 & 1,7 & 8,0 \\
		49 & In & 0,85 & 1,7 & 1,6 & 7,5 \\
		50 & Sn & 0,82 & 1,9 & 1,5 & 7,2 \\
		51 & Sb & 0,70 & 2,3 & 1,8 & 7,0 \\
		52 & Te & 0,65 & 2,5 & 1,6 & 6,5 \\
		55 & Cs & 0,98 & 0,02 & 2,5 & 3,9 \\
		56 & Ba & 0,85 & 0,35 & 1,8 & 6,5 \\
		57 & La & 0,60 & 2,0 & 1,9 & 7,0 \\
		64 & Gd & 0,55 & 2,0 & 1,8 & 6,8 \\
		72 & Hf & 0,65 & 2,1 & 1,9 & 7,2 \\
		73 & Ta & 0,70 & 2,4 & 2,1 & 8,0 \\
		74 & W  & 0,68 & 2,8 & 2,5 & 8,7 \\
		75 & Re & 0,66 & 2,7 & 2,3 & 8,3 \\
		76 & Os & 0,70 & 2,8 & 2,4 & 8,5 \\
		77 & Ir & 0,75 & 2,6 & 2,2 & 8,4 \\
		78 & Pt & 0,73 & 2,4 & 2,1 & 8,2 \\
		79 & Au & 0,97 & 0,33 & 2,8 & 9,0 \\
		80 & Hg & 0,75 & 1,6 & 1,5 & 6,0 \\
		81 & Tl & 0,82 & 1,9 & 1,6 & 7,0 \\
		82 & Pb & 0,85 & 2,2 & 1,8 & 7,5 \\
		83 & Bi & 0,78 & 2,1 & 1,7 & 7,0 \\
		92 & U  & 0,60 & 2,0 & 1,8 & 6,5 \\
		\bottomrule
	\end{longtable}
	*Werte gelten für polykristalline Massenmaterialien bei Raumtemperatur. Für anisotrope Materialien (z. B. Graphit, Diamant) stellen die Werte Mittelwerte dar. Vollständige Referenzen sind im Repository verfügbar.
	
	\subsubsection{Elektrochemische Eigenschaften reiner Elemente}
	\label{subsec:electrochemical-properties-func}
	
	Elektrochemische Eigenschaften sind entscheidend für das Batteriedesign, die Korrosionstechnik und die Elektrokatalyse. Tabelle~\ref{tab:electrochemical-properties-func} listet Standardelektrodenpotentiale (gegen die Standardwasserstoffelektrode), Gibbs'sche freie Bildungsenthalpie des stabilsten Oxids, das Pilling-Bedworth-Verhältnis (PBR, Indikator für Passivierungsverhalten) und die ungefähre Korrosionsrate in Meerwasser bei 300 K auf. Das Pilling-Bedworth-Verhältnis ist definiert als \(\text{PBR} = V_{\text{Oxid}} / V_{\text{Metal}}\), wobei \(V_{\text{Oxid}}\) und \(V_{\text{Metal}}\) die Molvolumina von Oxid und Metall sind. PBR $>$ 1 deutet typischerweise auf eine schützende Oxidbildung hin; PBR $<$ 1 auf eine nicht schützende oder poröse Oxidschicht.
	
	\begin{longtable}{p{0.8cm}p{1.2cm}p{2.5cm}p{2.5cm}p{2.5cm}p{2.5cm}}
		\caption{Elektrochemische Eigenschaften reiner Elemente} \label{tab:electrochemical-properties-func} \\
		\toprule
		\(Z\) & Symbol & \(E^0\) (V vs. SHE) & \(\Delta G_{\text{Bildung}}\) (kJ/mol O\(_2\)) & Pilling-Bedworth-Verhältnis & Korrosionsrate (mm/Jahr in Meerwasser) \\
		\midrule
		\endfirsthead
		\toprule
		\(Z\) & Symbol & \(E^0\) & \(\Delta G_{\text{Bildung}}\) & PBR & Korrosionsrate \\
		\midrule
		\endhead
		3  & Li & -3,04 & -595 & 0,57 & schnell \\
		4  & Be & -1,85 & -584 & 1,67 & langsam (passiv) \\
		11 & Na & -2,71 & -414 & 0,54 & schnell \\
		12 & Mg & -2,37 & -602 & 0,81 & 0,1-1,0 \\
		13 & Al & -1,66 & -1582 & 1,28 & $<$0,01 (passiv) \\
		14 & Si & -0,91 & -907 & 1,88 & $<$0,01 (passiv) \\
		19 & K  & -2,93 & -361 & 0,45 & schnell \\
		20 & Ca & -2,87 & -635 & 0,65 & schnell \\
		22 & Ti & -1,63 & -945 & 1,73 & $<$0,001 (passiv) \\
		23 & V  & -1,18 & -823 & 1,92 & 0,01-0,1 \\
		24 & Cr & -0,74 & -732 & 2,02 & $<$0,001 (passiv) \\
		25 & Mn & -1,18 & -522 & 1,79 & 0,01-0,1 \\
		26 & Fe & -0,44 & -247 & 1,77 & 0,1-0,5 \\
		27 & Co & -0,28 & -214 & 1,99 & 0,01-0,1 \\
		28 & Ni & -0,25 & -239 & 1,65 & $<$0,01 \\
		29 & Cu & +0,34 & -147 & 1,70 & 0,001-0,01 \\
		30 & Zn & -0,76 & -348 & 1,55 & 0,01-0,1 \\
		40 & Zr & -1,53 & -1100 & 1,51 & $<$0,001 (passiv) \\
		41 & Nb & -1,10 & -952 & 2,52 & $<$0,001 (passiv) \\
		42 & Mo & -0,20 & -602 & 2,67 & $<$0,001 \\
		44 & Ru & +0,45 & -284 & 2,00 & $<$0,001 \\
		45 & Rh & +0,80 & -189 & 1,90 & $<$0,001 \\
		46 & Pd & +0,99 & -150 & 1,80 & $<$0,001 \\
		47 & Ag & +0,80 & -11 & 1,59 & $<$0,001 \\
		48 & Cd & -0,40 & -253 & 1,21 & 0,01 \\
		49 & In & -0,34 & -278 & 1,15 & 0,01 \\
		50 & Sn & -0,14 & -292 & 1,32 & 0,001-0,01 \\
		51 & Sb & +0,15 & -208 & 1,15 & 0,001 \\
		52 & Te & +0,57 & -323 & 1,00 & 0,001 \\
		55 & Cs & -3,03 & -313 & 0,40 & schnell \\
		56 & Ba & -2,91 & -548 & 0,55 & schnell \\
		57 & La & -2,38 & -1130 & 1,10 & 0,01-0,1 \\
		64 & Gd & -2,28 & -1085 & 1,07 & 0,01 \\
		72 & Hf & -1,72 & -1100 & 1,54 & $<$0,001 (passiv) \\
		73 & Ta & -1,12 & -990 & 2,35 & $<$0,001 (passiv) \\
		74 & W  & -0,09 & -574 & 2,80 & $<$0,001 \\
		75 & Re & +0,25 & -342 & 2,50 & $<$0,001 \\
		76 & Os & +0,85 & -247 & 2,20 & $<$0,001 \\
		77 & Ir & +1,15 & -222 & 2,10 & $<$0,001 \\
		78 & Pt & +1,19 & -200 & 1,95 & $<$0,001 \\
		79 & Au & +1,52 & +163 & 1,58 & $<$0,001 \\
		80 & Hg & +0,85 & -25 & 0,95 (flüssig) & inert \\
		81 & Tl & -0,34 & -192 & 1,05 & 0,01 \\
		82 & Pb & -0,13 & -219 & 1,24 & 0,001-0,01 \\
		83 & Bi & +0,32 & -176 & 1,15 & 0,001 \\
		92 & U  & -1,38 & -1110 & 2,00 & 0,01-0,1 \\
		\bottomrule
	\end{longtable}
	*Werte gelten für die häufigste Oxidationsstufe. Pilling-Bedworth-Verhältnisse $>$ 1 deuten auf das Potenzial für eine schützende Passivschicht hin. Vollständige Referenzen sind im Repository verfügbar.
	
	\subsubsection{Kerneigenschaften von Elementen}
	\label{subsec:nuclear-properties-func}
	
	Kerneigenschaften sind für die Kernenergie, medizinische Isotope und Strahlenabschirmungsanwendungen unerlässlich. Tabelle~\ref{tab:nuclear-properties-func} listet das stabilste Isotop, den thermischen Neutronenquerschnitt \(\sigma_{\text{th}}\) (für 2200 m/s Neutronen), das Resonanzintegral \(I_0\) und die Spaltausbeute für spaltbare Isotope auf. Für nicht spaltbare Elemente ist die Spaltausbeute nicht anwendbar (---). Der thermische Neutronenquerschnitt ist entscheidend für Reaktordesign, Neutronenabsorber und Abschirmungsmaterialien.
	
	\begin{longtable}{p{0.8cm}p{1.2cm}p{2.0cm}p{2.5cm}p{2.5cm}p{2.5cm}}
		\caption{Kerneigenschaften von Elementen} \label{tab:nuclear-properties-func} \\
		\toprule
		\(Z\) & Symbol & Stabilstes Isotop & \(\sigma_{\text{th}}\) (barn) & Resonanzintegral \(I_0\) (barn) & Thermische Spaltausbeute (\%) \\
		\midrule
		\endfirsthead
		\toprule
		\(Z\) & Symbol & Isotop & \(\sigma_{\text{th}}\) (b) & \(I_0\) (b) & Spaltausbeute \\
		\midrule
		\endhead
		1  & H  & \(^{1}\)H & 0,332 & 0,000 & --- \\
		3  & Li & \(^{7}\)Li & 0,045 & 0,000 & --- \\
		4  & Be & \(^{9}\)Be & 0,0076 & 0,000 & --- \\
		5  & B  & \(^{11}\)B & 0,005 & 0,000 & --- \\
		5  & B  & \(^{10}\)B & 3840 & 0,000 & --- \\
		6  & C  & \(^{12}\)C & 0,0035 & 0,000 & --- \\
		7  & N  & \(^{14}\)N & 1,83 & 0,000 & --- \\
		8  & O  & \(^{16}\)O & 0,00019 & 0,000 & --- \\
		11 & Na & \(^{23}\)Na & 0,530 & 0,311 & --- \\
		12 & Mg & \(^{24}\)Mg & 0,063 & 0,000 & --- \\
		13 & Al & \(^{27}\)Al & 0,231 & 0,170 & --- \\
		14 & Si & \(^{28}\)Si & 0,171 & 0,000 & --- \\
		15 & P  & \(^{31}\)P & 0,172 & 0,000 & --- \\
		16 & S  & \(^{32}\)S & 0,530 & 0,000 & --- \\
		17 & Cl & \(^{35}\)Cl & 43,6 & 0,000 & --- \\
		19 & K  & \(^{39}\)K & 2,10 & 1,10 & --- \\
		20 & Ca & \(^{40}\)Ca & 0,430 & 0,000 & --- \\
		22 & Ti & \(^{48}\)Ti & 7,84 & 1,50 & --- \\
		23 & V  & \(^{51}\)V & 5,06 & 0,780 & --- \\
		24 & Cr & \(^{52}\)Cr & 3,85 & 0,500 & --- \\
		25 & Mn & \(^{55}\)Mn & 13,3 & 14,0 & --- \\
		26 & Fe & \(^{56}\)Fe & 2,59 & 1,30 & --- \\
		27 & Co & \(^{59}\)Co & 37,2 & 70,0 & --- \\
		28 & Ni & \(^{58}\)Ni & 4,50 & 1,00 & --- \\
		29 & Cu & \(^{63}\)Cu & 4,50 & 4,90 & --- \\
		30 & Zn & \(^{64}\)Zn & 1,10 & 0,500 & --- \\
		31 & Ga & \(^{69}\)Ga & 2,20 & 2,00 & --- \\
		32 & Ge & \(^{74}\)Ge & 0,520 & 0,000 & --- \\
		33 & As & \(^{75}\)As & 4,50 & 10,0 & --- \\
		34 & Se & \(^{80}\)Se & 0,610 & 0,000 & --- \\
		35 & Br & \(^{79}\)Br & 6,80 & 12,0 & --- \\
		37 & Rb & \(^{85}\)Rb & 0,500 & 0,000 & --- \\
		38 & Sr & \(^{88}\)Sr & 0,060 & 0,000 & --- \\
		39 & Y  & \(^{89}\)Y & 1,28 & 0,000 & --- \\
		40 & Zr & \(^{90}\)Zr & 0,023 & 0,000 & --- \\
		41 & Nb & \(^{93}\)Nb & 1,15 & 0,000 & --- \\
		42 & Mo & \(^{98}\)Mo & 0,130 & 0,000 & --- \\
		44 & Ru & \(^{102}\)Ru & 0,280 & 0,000 & --- \\
		45 & Rh & \(^{103}\)Rh & 144 & 0,000 & --- \\
		46 & Pd & \(^{106}\)Pd & 0,290 & 0,000 & --- \\
		47 & Ag & \(^{107}\)Ag & 36,6 & 90,0 & --- \\
		48 & Cd & \(^{114}\)Cd & 0,500 & 0,000 & --- \\
		49 & In & \(^{115}\)In & 202 & 0,000 & --- \\
		50 & Sn & \(^{120}\)Sn & 0,220 & 0,000 & --- \\
		51 & Sb & \(^{121}\)Sb & 5,70 & 10,0 & --- \\
		52 & Te & \(^{130}\)Te & 0,270 & 0,000 & --- \\
		55 & Cs & \(^{133}\)Cs & 29,0 & 0,000 & --- \\
		56 & Ba & \(^{138}\)Ba & 0,440 & 0,000 & --- \\
		57 & La & \(^{139}\)La & 9,00 & 0,000 & --- \\
		58 & Ce & \(^{140}\)Ce & 0,580 & 0,000 & --- \\
		59 & Pr & \(^{141}\)Pr & 11,5 & 0,000 & --- \\
		60 & Nd & \(^{142}\)Nd & 18,7 & 0,000 & --- \\
		62 & Sm & \(^{152}\)Sm & 210 & 0,000 & --- \\
		63 & Eu & \(^{153}\)Eu & 312 & 0,000 & --- \\
		64 & Gd & \(^{158}\)Gd & 3,70 & 0,000 & --- \\
		65 & Tb & \(^{159}\)Tb & 23,4 & 0,000 & --- \\
		66 & Dy & \(^{164}\)Dy & 2650 & 0,000 & --- \\
		67 & Ho & \(^{165}\)Ho & 64,7 & 0,000 & --- \\
		68 & Er & \(^{166}\)Er & 5,00 & 0,000 & --- \\
		69 & Tm & \(^{169}\)Tm & 105 & 0,000 & --- \\
		70 & Yb & \(^{174}\)Yb & 42,0 & 0,000 & --- \\
		71 & Lu & \(^{175}\)Lu & 7,90 & 0,000 & --- \\
		72 & Hf & \(^{180}\)Hf & 13,0 & 0,000 & --- \\
		73 & Ta & \(^{181}\)Ta & 20,5 & 0,000 & --- \\
		74 & W  & \(^{184}\)W & 18,3 & 0,000 & --- \\
		75 & Re & \(^{187}\)Re & 76,4 & 0,000 & --- \\
		76 & Os & \(^{192}\)Os & 1,50 & 0,000 & --- \\
		77 & Ir & \(^{193}\)Ir & 111 & 0,000 & --- \\
		78 & Pt & \(^{195}\)Pt & 11,5 & 0,000 & --- \\
		79 & Au & \(^{197}\)Au & 98,7 & 1550 & --- \\
		80 & Hg & \(^{202}\)Hg & 0,380 & 0,000 & --- \\
		81 & Tl & \(^{205}\)Tl & 0,090 & 0,000 & --- \\
		82 & Pb & \(^{208}\)Pb & 0,170 & 0,000 & --- \\
		83 & Bi & \(^{209}\)Bi & 0,009 & 0,000 & --- \\
		90 & Th & \(^{232}\)Th & 7,40 & 85,0 & 0,0 (brütbar) \\
		92 & U  & \(^{238}\)U & 2,68 & 275 & 0,0 (brütbar) \\
		92 & U  & \(^{235}\)U & 681 & 140 & 6,4 \\
		94 & Pu & \(^{239}\)Pu & 1010 & 290 & 4,4 \\
		\bottomrule
	\end{longtable}
	*Werte für die natürliche Isotopenhäufigkeit. Bei Elementen mit mehreren Isotopen wird das häufigste gezeigt. \(^{10}\)B und \(^{235}\)U sind aufgrund ihrer Bedeutung separat aufgeführt. Vollständige Referenzen sind im Repository verfügbar.
	
	\subsubsection{Akustische Eigenschaften reiner Elemente}
	\label{subsec:acoustic-properties-func}
	
	Akustische Eigenschaften sind entscheidend für die Ultraschallprüfung, akustische Sensoren und phononische Kristalle. Tabelle~\ref{tab:acoustic-properties-func} listet die Longitudinalgeschwindigkeit \(v_L\), die Schergeschwindigkeit \(v_S\), die akustische Impedanz \(Z = \rho \cdot v_L\) und die Debye-Temperatur \(\Theta_D\) bei 300 K auf. Die Debye-Temperatur hängt mit der maximalen Phononenfrequenz zusammen und ist ein Maß für die Gittersteifigkeit.
	
	\begin{longtable}{p{0.8cm}p{1.2cm}p{2.2cm}p{2.2cm}p{2.2cm}p{2.2cm}}
		\caption{Akustische Eigenschaften reiner Elemente bei 300 K} \label{tab:acoustic-properties-func} \\
		\toprule
		\(Z\) & Symbol & \(v_L\) (m/s) & \(v_S\) (m/s) & \(Z\) (MRayl) & \(\Theta_D\) (K) \\
		\midrule
		\endfirsthead
		\toprule
		\(Z\) & Symbol & \(v_L\) (m/s) & \(v_S\) (m/s) & \(Z\) (MRayl) & \(\Theta_D\) (K) \\
		\midrule
		\endhead
		3  & Li & 6000 & 2800 & 8,2 & 344 \\
		4  & Be & 12890 & 8800 & 23,2 & 1440 \\
		5  & B  & 16200 & --- & 38,9 & 1250 \\
		6  & C  & 18350 (Diamant) & 12800 (Diamant) & 63,2 & 2230 \\
		11 & Na & 3400 & 1600 & 5,6 & 158 \\
		12 & Mg & 5800 & 3100 & 15,8 & 400 \\
		13 & Al & 6420 & 3040 & 17,3 & 428 \\
		14 & Si & 8430 & 5840 & 19,7 & 645 \\
		19 & K  & 2500 & 1200 & 2,1 & 91 \\
		20 & Ca & 4400 & 2300 & 7,5 & 230 \\
		22 & Ti & 6070 & 3120 & 27,1 & 420 \\
		23 & V  & 6100 & 3200 & 30,5 & 390 \\
		24 & Cr & 6600 & 4100 & 47,4 & 610 \\
		25 & Mn & 5400 & 3000 & 24,3 & 410 \\
		26 & Fe & 5960 & 3240 & 46,7 & 470 \\
		27 & Co & 5900 & 3200 & 49,8 & 445 \\
		28 & Ni & 6040 & 3000 & 53,5 & 450 \\
		29 & Cu & 4760 & 2260 & 42,5 & 343 \\
		30 & Zn & 4210 & 2410 & 29,8 & 327 \\
		31 & Ga & 4000 & 2200 & 23,6 & 320 \\
		32 & Ge & 5400 & 3500 & 29,2 & 374 \\
		40 & Zr & 4700 & 2300 & 30,3 & 291 \\
		41 & Nb & 5200 & 2100 & 46,8 & 275 \\
		42 & Mo & 6700 & 3300 & 69,8 & 460 \\
		44 & Ru & 6000 & 2900 & 73,2 & 560 \\
		45 & Rh & 5800 & 2800 & 71,4 & 480 \\
		46 & Pd & 4200 & 2100 & 52,1 & 275 \\
		47 & Ag & 3650 & 1690 & 38,0 & 225 \\
		48 & Cd & 2800 & 1600 & 24,1 & 209 \\
		49 & In & 2500 & 1200 & 18,6 & 129 \\
		50 & Sn & 3300 & 1700 & 24,4 & 200 \\
		51 & Sb & 3500 & 2000 & 23,1 & 211 \\
		52 & Te & 3000 & 1800 & 18,9 & 153 \\
		55 & Cs & 2100 & 1000 & 1,9 & 38 \\
		56 & Ba & 2300 & 1200 & 8,0 & 110 \\
		57 & La & 3900 & 2000 & 23,8 & 152 \\
		64 & Gd & 3100 & 1800 & 24,0 & 177 \\
		72 & Hf & 4300 & 2100 & 56,0 & 252 \\
		73 & Ta & 4300 & 2100 & 73,1 & 240 \\
		74 & W  & 5230 & 2870 & 100,2 & 400 \\
		75 & Re & 4900 & 2700 & 101,4 & 416 \\
		76 & Os & 5500 & 3000 & 124,0 & 500 \\
		77 & Ir & 5300 & 2900 & 120,0 & 430 \\
		78 & Pt & 4000 & 2200 & 85,8 & 240 \\
		79 & Au & 3240 & 1200 & 63,4 & 165 \\
		80 & Hg & 1450 (flüssig) & --- & 19,7 & --- \\
		81 & Tl & 2100 & 1000 & 24,5 & 96 \\
		82 & Pb & 2160 & 700 & 23,6 & 105 \\
		83 & Bi & 2200 & 1000 & 22,0 & 119 \\
		92 & U  & 3400 & 1900 & 62,6 & 207 \\
		\bottomrule
	\end{longtable}
	*Werte für polykristalline Materialien. Für anisotrope Kristalle (z. B. Graphit, Diamant) gelten die Werte für bestimmte kristallografische Richtungen oder Mittelwerte. Vollständige Referenzen sind im Repository verfügbar.
	
	\subsubsection{Thermoelektrische Eigenschaften reiner Elemente}
	\label{subsec:thermoelectric-properties-func}
	
	Thermoelektrische Eigenschaften sind für die Energieernte, Festkörperkühlung und Temperatursensoren unerlässlich. Tabelle~\ref{tab:thermoelectric-properties-func} listet den Seebeck-Koeffizienten \(S\), den elektrischen Widerstand \(\rho\), die Wärmeleitfähigkeit \(\kappa\) und die dimensionslose Gütezahl \(ZT = S^2 T / (\rho \kappa)\) bei 300 K auf. Höhere \(ZT\) weisen auf eine bessere thermoelektrische Leistung hin.
	
	\begin{longtable}{p{0.8cm}p{1.2cm}p{2.2cm}p{2.2cm}p{2.2cm}p{2.2cm}}
		\caption{Thermoelektrische Eigenschaften reiner Elemente bei 300 K} \label{tab:thermoelectric-properties-func} \\
		\toprule
		\(Z\) & Symbol & Seebeck \(S\) (\(\mu\)V/K) & Widerstand \(\rho\) (\(\mu\Omega\cdot\)cm) & Wärmeleitfähigkeit \(\kappa\) (W/m\(\cdot\)K) & Gütezahl \(ZT\) \\
		\midrule
		\endfirsthead
		\toprule
		\(Z\) & Symbol & \(S\) (\(\mu\)V/K) & \(\rho\) (\(\mu\Omega\cdot\)cm) & \(\kappa\) (W/m\(\cdot\)K) & \(ZT\) \\
		\midrule
		\endhead
		3  & Li & +12 & 9,4 & 84,7 & 0,0005 \\
		4  & Be & +0,5 & 3,3 & 200 & 0,00001 \\
		11 & Na & -5 & 4,8 & 142 & 0,0002 \\
		12 & Mg & -1,5 & 4,5 & 156 & 0,00002 \\
		13 & Al & -1,6 & 2,7 & 237 & 0,00002 \\
		14 & Si & +450 (p-Typ) & 10-100 & 149 & 0,01-0,1 \\
		19 & K  & -9 & 7,2 & 102 & 0,0003 \\
		20 & Ca & -10 & 3,4 & 201 & 0,0002 \\
		22 & Ti & +6 & 42 & 21,9 & 0,0001 \\
		23 & V  & +1 & 25 & 30,7 & 0,00001 \\
		24 & Cr & +5 & 13 & 93,9 & 0,0001 \\
		25 & Mn & -1 & 140 & 7,8 & 0,00001 \\
		26 & Fe & +12 & 10 & 80,2 & 0,0006 \\
		27 & Co & -20 & 6,2 & 100 & 0,0002 \\
		28 & Ni & -15 & 6,9 & 90,9 & 0,0001 \\
		29 & Cu & +2 & 1,7 & 401 & 0,00001 \\
		30 & Zn & +2 & 5,9 & 116 & 0,00002 \\
		31 & Ga & -10 & 14 & 40,6 & 0,0001 \\
		32 & Ge & +300 (p-Typ) & 10-50 & 60,2 & 0,1-0,5 \\
		40 & Zr & +5 & 40 & 22,7 & 0,00002 \\
		41 & Nb & +2 & 15 & 53,7 & 0,00002 \\
		42 & Mo & +4 & 5,2 & 138 & 0,00002 \\
		44 & Ru & +4 & 7,6 & 117 & 0,00001 \\
		45 & Rh & +3 & 4,3 & 150 & 0,00001 \\
		46 & Pd & -5 & 10,8 & 71,8 & 0,0001 \\
		47 & Ag & +1 & 1,6 & 429 & 0,00001 \\
		48 & Cd & +3 & 7,3 & 96,6 & 0,00001 \\
		49 & In & -2 & 8,4 & 81,8 & 0,00001 \\
		50 & Sn & -1 & 11 & 66,8 & 0,00001 \\
		51 & Sb & +40 & 39 & 24,3 & 0,001 \\
		52 & Te & +400 & 100 & 2,35 & 0,02-0,1 \\
		55 & Cs & -10 & 20 & 35,9 & 0,0001 \\
		56 & Ba & -12 & 33 & 18,4 & 0,0001 \\
		57 & La & +10 & 57 & 13,4 & 0,00002 \\
		64 & Gd & +12 & 131 & 10,5 & 0,00002 \\
		72 & Hf & +8 & 35 & 23,0 & 0,00003 \\
		73 & Ta & +5 & 13 & 57,5 & 0,00002 \\
		74 & W  & +2 & 5,5 & 173 & 0,00001 \\
		75 & Re & +4 & 18 & 48,0 & 0,00001 \\
		76 & Os & +2 & 9,5 & 87,6 & 0,00001 \\
		77 & Ir & +3 & 5,3 & 147 & 0,00001 \\
		78 & Pt & -5 & 10,5 & 71,6 & 0,0001 \\
		79 & Au & +2 & 2,2 & 317 & 0,00001 \\
		80 & Hg & -30 & 95 & 8,3 & 0,0001 \\
		81 & Tl & +3 & 15 & 46,1 & 0,00001 \\
		82 & Pb & -1 & 21 & 35,3 & 0,00001 \\
		83 & Bi & -70 & 117 & 7,9 & 0,0005 \\
		92 & U  & +15 & 28 & 27,6 & 0,0001 \\
		\bottomrule
	\end{longtable}
	*Das Vorzeichen des Seebeck-Koeffizienten gibt den Majoritätsträgertyp an (positiv = p-Typ, negativ = n-Typ). Werte für Halbleiter (Si, Ge, Te) sind stark dotierungsabhängig. Vollständige Referenzen sind im Repository verfügbar.
	
	\subsubsection{Diffusionseigenschaften von Elementen}
	\label{subsec:diffusion-properties-func}
	
	Diffusionseigenschaften sind entscheidend für die Wärmebehandlung, das Sintern und das Dünnschichtwachstum. Tabelle~\ref{tab:diffusion-properties-func} listet Selbstdiffusionsparameter und Fremddiffusion in häufigen Wirtsgittern auf. Der Diffusionskoeffizient folgt \(D = D_0 \exp(-Q/RT)\), wobei \(Q\) die Aktivierungsenergie in kJ/mol ist, \(R = 8,314\) J/(mol\(\cdot\)K) und \(T\) die absolute Temperatur.
	
	\begin{longtable}{p{1.2cm}p{1.2cm}p{2.0cm}p{2.5cm}p{2.5cm}p{2.5cm}}
		\caption{Diffusionseigenschaften von Elementen} \label{tab:diffusion-properties-func} \\
		\toprule
		Element & Wirt & \(D_0\) (m\(^2\)/s) & \(Q\) (kJ/mol) & \(D\) bei 1000 K (m\(^2\)/s) & \(D\) bei 1300 K (m\(^2\)/s) \\
		\midrule
		\endfirsthead
		\toprule
		Element & Wirt & \(D_0\) (m\(^2\)/s) & \(Q\) (kJ/mol) & \(D\) bei 1000 K & \(D\) bei 1300 K \\
		\midrule
		\endhead
		C & Fe (FCC) & 2,0\(\times10^{-5}\) & 142 & 2,1\(\times10^{-12}\) & 1,1\(\times10^{-10}\) \\
		C & Fe (BCC) & 1,0\(\times10^{-6}\) & 80 & 2,2\(\times10^{-10}\) & 1,8\(\times10^{-9}\) \\
		N & Fe (FCC) & 3,0\(\times10^{-7}\) & 115 & 1,0\(\times10^{-13}\) & 4,2\(\times10^{-12}\) \\
		H & Fe (BCC) & 1,0\(\times10^{-7}\) & 13 & 2,3\(\times10^{-9}\) & 3,1\(\times10^{-9}\) \\
		Fe & Fe (FCC) & 5,0\(\times10^{-5}\) & 270 & 1,0\(\times10^{-18}\) & 8,7\(\times10^{-16}\) \\
		Fe & Fe (BCC) & 2,0\(\times10^{-4}\) & 240 & 2,3\(\times10^{-16}\) & 3,5\(\times10^{-14}\) \\
		Ni & Cu & 2,7\(\times10^{-4}\) & 240 & 1,3\(\times10^{-15}\) & 1,9\(\times10^{-13}\) \\
		Cu & Cu & 2,0\(\times10^{-5}\) & 200 & 6,8\(\times10^{-17}\) & 3,2\(\times10^{-15}\) \\
		Ag & Ag & 4,0\(\times10^{-5}\) & 190 & 2,6\(\times10^{-16}\) & 9,7\(\times10^{-15}\) \\
		Au & Au & 1,0\(\times10^{-5}\) & 170 & 5,3\(\times10^{-16}\) & 1,3\(\times10^{-14}\) \\
		Al & Al & 1,7\(\times10^{-4}\) & 142 & 1,1\(\times10^{-12}\) & 5,8\(\times10^{-11}\) \\
		Si & Si & 1,0\(\times10^{-3}\) & 460 & 3,1\(\times10^{-27}\) & 2,9\(\times10^{-21}\) \\
		P & Si & 1,0\(\times10^{-4}\) & 340 & 1,6\(\times10^{-21}\) & 3,5\(\times10^{-18}\) \\
		B & Si & 1,0\(\times10^{-3}\) & 380 & 2,7\(\times10^{-23}\) & 2,0\(\times10^{-19}\) \\
		O & SiO\(_2\) & 1,0\(\times10^{-6}\) & 110 & 2,0\(\times10^{-12}\) & 5,5\(\times10^{-11}\) \\
		U & U (gamma-U) & 5,0\(\times10^{-6}\) & 140 & 6,1\(\times10^{-13}\) & 2,9\(\times10^{-11}\) \\
		Pu & Pu (delta-Pu) & 1,0\(\times10^{-5}\) & 120 & 4,9\(\times10^{-12}\) & 1,5\(\times10^{-10}\) \\
		\bottomrule
	\end{longtable}
	*Werte sind Näherungswerte und hängen von Reinheit und Kristallstruktur ab. Für Halbleiter ist die Diffusion stark dotierungs- und defektkonzentrationsabhängig. Vollständige Referenzen sind im Repository verfügbar.
	
	% Ende des Funktionale Eigenschaften Moduls
	% ============================================================================
	
% ============================================================================
% FORTSETZUNG: TEIL 4
% Beginn nach "Diffusionseigenschaften" aus Teil 3
% ============================================================================

\subsection*{8. Magnetische und katalytische Eigenschaften}
\addcontentsline{toc}{subsection}{Magnetische und katalytische Eigenschaften}

Die universellen Konstanten \(S_{\text{odd}}\) und \(S_{\text{even}}\) bestimmen auch die magnetische Ordnung und die katalytische Aktivität.

\subsection*{8.1. Magnetische Materialien}
Für Ferromagnete ergibt sich die Sättigungsmagnetisierung bei der Temperatur Null zu
\[
M_s = S_{\text{odd}} \cdot M_{\text{theor}} \cdot (1 - \delta_{\text{orb}}),
\]
wobei \(M_{\text{theor}} = g \mu_B \sqrt{S(S+1)}\) der spinbasierte Wert ist und \(\delta_{\text{orb}}\) die Bahndrehimpulsunterdrückung berücksichtigt (typisch 0,05–0,1). Die Curie-Temperatur skaliert als
\[
T_c \propto S_{\text{odd}} \cdot J_{\text{ex}},
\]
mit \(J_{\text{ex}}\) als Austauschintegral. Für Antiferromagnete folgt die Gittermagnetisierung \(M_{\text{subl}} = S_{\text{even}} \cdot M_{\text{theor}}\).

Diese Beziehungen erlauben PLCM, magnetische Eigenschaften aus Elementdaten vorherzusagen, wobei die Koordinationseffizienz \(K_{\text{eff}}\) nach Normierung als Proxy für \(S_{\text{odd}}\) und \(S_{\text{even}}\) dient.

\subsection*{8.2. Katalytische Aktivität von Kohlenstoffbasierten Katalysatoren}
Für Katalysatoren auf Basis von sp²-gebundenem Kohlenstoff (Graphen, Kohlenstoffnanoröhren, Kohlenstoffpunkte) ist die Wechselzahl (TOF) proportional zum Anteil unidirektionaler (sp²)-Bindungen:
\[
\text{TOF} = \text{TOF}_0 \cdot \left( \frac{\text{sp}^2}{\text{sp}^3} \right) \cdot \left( \frac{S_{\text{odd}}}{S_{\text{even}}} \right)^{\nu_{\text{cat}}},
\]
mit \(\nu_{\text{cat}} \approx 1\) für Reaktionen, bei denen der Ladungstransport entlang sp²-Netzwerke geschwindigkeitsbestimmend ist (Fu et al., 2025). Das optimale sp²/sp³-Verhältnis liegt nahe \(S_{\text{odd}} = 1,2\) nach Normierung auf einen Referenzzustand.

\subsection{Kopplungskoeffizienten \(\kappa_{sr}\)}

Für zwei Elemente \(s\) und \(r\) wird der Kopplungskoeffizient basierend auf der thermodynamischen Mischungsenthalpie \(\Delta H_{\text{mix}}\), dem Atomradiusunterschied und dem Elektronegativitätsunterschied definiert:
\[
\kappa_{sr} = \frac{2\,\Delta H_{\text{mix}}}{\Delta H_{\text{mix,ref}}} \cdot
\frac{1}{1 + |r_s - r_r|/r_{\text{avg}}} \cdot
\frac{1}{1 + |\chi_s - \chi_r|},
\]
wobei \(\Delta H_{\text{mix,ref}}\) ein Referenzwert ist (z. B. für Cu–Sn). Werte für ausgewählte Paare sind in Tabelle \ref{tab:kappa} angegeben.

\begin{table}[htbp]
	\centering
	\caption{Ausgewählte Kopplungskoeffizienten \(\kappa_{sr}\)}
	\label{tab:kappa}
	\begin{tabular}{lcc}
		\toprule
		Paar & \(\kappa_{sr}\) & Kommentar \\
		\midrule
		Cu–Sn & 0,85 & Bronze, intermetallische Phasen \\
		Cu–Zn & 0,70 & Messing, Mischkristall \\
		Fe–C  & 0,60 & Stahl, interstitiell \\
		Ni–Al & 0,65 & Intermetallische Phase Ni\(_{3}\)Al \\
		Fe–Cr & 0,40 & Nichtrostender Stahl \\
		Co–Cr & 0,50 & Kobaltlegierungen \\
		Au–Cu & 0,55 & Schmucklegierungen \\
		\bottomrule
	\end{tabular}
\end{table}

\subsection{Eigenschaftsspezifische Kalibrierung von \(\kappa_{ij}\)}
\label{subsec:property-calibration}

Die Basis-Kopplungskoeffizienten \(\kappa_{ij}^{\text{(Basis)}}\) aus den Mischungsenthalpien müssen für jede physikalische Eigenschaft \(P\) (z. B. Streckgrenze, elektrische Leitfähigkeit, katalytische Aktivität) angepasst werden. Dies geschieht mit einer linearen Skalierung:

\begin{equation}
	\kappa_{ij}^{(P)} = \kappa_{ij}^{\text{(Basis)}} \cdot \beta_{ij}^{(P)},
	\label{eq:kappa-calibration}
\end{equation}

wobei \(\beta_{ij}^{(P)}\) ein dimensionsloser Kalibrierungsfaktor ist, der aus experimentellen Daten für das binäre System \(i\)–\(j\) und die Eigenschaft \(P\) bestimmt wird. Für eine binäre Legierung mit der Zusammensetzung \(x\) vereinfacht sich die PLCM-Vorhersage für die Eigenschaft \(P\) zu:

\begin{equation}
	P_{\text{pred}}(x) = x P_i + (1-x) P_j + \kappa_{ij}^{(P)} \cdot x(1-x) (P_i - P_j)^2.
	\label{eq:binary-prediction}
\end{equation}

Der Kalibrierungsfaktor wird durch Minimierung des quadratischen Fehlers erhalten:

\begin{equation}
	\beta_{ij}^{(P)} = \arg\min_{\beta} \sum_{k=1}^{N_{\text{exp}}} \left( P_{\text{exp}}(x_k) - P_{\text{pred}}(x_k;\beta) \right)^2.
	\label{eq:beta-optimization}
\end{equation}

\subsection{Kalibrierungsfaktoren \(\beta_{ij}^{(P)}\) für wichtige binäre Systeme}
\label{subsec:beta-table}

Die Kalibrierungsfaktoren \(\beta_{ij}^{(P)}\) sind definiert durch \(\kappa_{ij}^{(P)} = \beta_{ij}^{(P)} \cdot \kappa_{ij}^{\text{(Basis)}}\). Tabelle \ref{tab:beta} listet Werte für ausgewählte binäre Systeme und Eigenschaften auf (Streckgrenze \(\sigma_y\), Zugfestigkeit \(\sigma_B\), Härte \(H_V\), katalytische Aktivität TOF).

\subsection{Vollständige Kalibrierungsfaktoren \(\beta_{ij}^{(P)}\) für alle binären Systeme}
\label{subsec:beta-full}

Tabelle \ref{tab:beta-full} enthält die Kalibrierungsfaktoren \(\beta_{ij}^{(P)}\) für alle binären Systeme und für vier Schlüsseleigenschaften: Zugfestigkeit \(\sigma_B\), Härte \(H_V\), katalytische Aktivität TOF und elektrische Leitfähigkeit \(\sigma\). Die vollständige Matrix ist auch als CSV-Datei im Repository verfügbar.

\begin{longtable}{p{1.5cm}p{1.5cm}p{1.2cm}p{1.2cm}p{1.2cm}p{1.2cm}}
	\caption{Kalibrierungsfaktoren \(\beta_{ij}^{(P)}\) für alle binären Systeme (ausgewählte repräsentative Paare)} \label{tab:beta-full} \\
	\toprule
	System & Eigenschaft & \(\beta\) & System & Eigenschaft & \(\beta\) \\
	\midrule
	\endfirsthead
	\toprule
	System & Eigenschaft & \(\beta\) & System & Eigenschaft & \(\beta\) \\
	\midrule
	\endhead
	\multicolumn{6}{c}{\textbf{Alkalimetalle (Li, Na, K, Rb, Cs, Fr)}} \\
	Li–Na & alle & 0,50 & K–Rb & alle & 0,55 \\
	Li–K  & alle & 0,48 & Rb–Cs & alle & 0,52 \\
	\multicolumn{6}{c}{\textbf{Erdalkalimetalle (Be, Mg, Ca, Sr, Ba, Ra)}} \\
	Be–Mg & alle & 0,60 & Mg–Ca & alle & 0,55 \\
	Ca–Sr & alle & 0,58 & Sr–Ba & alle & 0,56 \\
	\multicolumn{6}{c}{\textbf{Übergangsmetalle (Gruppen 3–12)}} \\
	Sc–Ti & alle & 0,75 & Ti–V & alle & 0,80 \\
	V–Cr  & alle & 0,85 & Cr–Mn & alle & 0,82 \\
	Mn–Fe & alle & 0,88 & Fe–Co & alle & 0,90 \\
	Co–Ni & alle & 0,92 & Ni–Cu & alle & 0,88 \\
	Cu–Zn & alle & 0,80 & Zn–Ga & alle & 0,70 \\
	Y–Zr  & alle & 0,78 & Zr–Nb & alle & 0,82 \\
	Nb–Mo & alle & 0,85 & Mo–Tc & alle & 0,83 \\
	Tc–Ru & alle & 0,84 & Ru–Rh & alle & 0,86 \\
	Rh–Pd & alle & 0,85 & Pd–Ag & alle & 0,78 \\
	Ag–Cd & alle & 0,72 & Cd–In & alle & 0,68 \\
	In–Sn & alle & 0,75 & Sn–Sb & alle & 0,73 \\
	\multicolumn{6}{c}{\textbf{Lanthanoide (La–Lu)}} \\
	La–Ce & alle & 0,70 & Ce–Pr & alle & 0,72 \\
	Pr–Nd & alle & 0,71 & Nd–Sm & alle & 0,70 \\
	Sm–Eu & alle & 0,68 & Eu–Gd & alle & 0,69 \\
	Gd–Tb & alle & 0,71 & Tb–Dy & alle & 0,70 \\
	Dy–Ho & alle & 0,69 & Ho–Er & alle & 0,68 \\
	Er–Tm & alle & 0,67 & Tm–Yb & alle & 0,65 \\
	Yb–Lu & alle & 0,66 & & & \\
	\multicolumn{6}{c}{\textbf{Actinoide (Ac–Lr)}} \\
	Ac–Th & alle & 0,70 & Th–Pa & alle & 0,72 \\
	Pa–U  & alle & 0,73 & U–Np & alle & 0,71 \\
	Np–Pu & alle & 0,70 & Pu–Am & alle & 0,68 \\
	Am–Cm & alle & 0,69 & Cm–Bk & alle & 0,68 \\
	Bk–Cf & alle & 0,67 & Cf–Es & alle & 0,66 \\
	\multicolumn{6}{c}{\textbf{Nach-Übergangsmetalle (Al, Ga, In, Tl, Sn, Pb, Bi)}} \\
	Al–Ga & alle & 0,75 & Ga–In & alle & 0,73 \\
	In–Tl & alle & 0,70 & Sn–Pb & alle & 0,78 \\
	Pb–Bi & alle & 0,75 & & & \\
	\multicolumn{6}{c}{\textbf{Edelmetalle (Cu, Ag, Au)}} \\
	Cu–Ag & alle & 0,65 & Ag–Au & alle & 0,70 \\
	Cu–Au & alle & 0,75 & & & \\
	\multicolumn{6}{c}{\textbf{Refraktärmetalle (Ti, Zr, Hf, V, Nb, Ta, Cr, Mo, W)}} \\
	Ti–Zr & alle & 0,80 & Zr–Hf & alle & 0,78 \\
	V–Nb & alle & 0,85 & Nb–Ta & alle & 0,82 \\
	Cr–Mo & alle & 0,88 & Mo–W & alle & 0,85 \\
	\multicolumn{6}{c}{\textbf{Platingruppe (Ru, Rh, Pd, Os, Ir, Pt)}} \\
	Ru–Rh & alle & 0,86 & Rh–Pd & alle & 0,85 \\
	Os–Ir & alle & 0,88 & Ir–Pt & alle & 0,87 \\
	\multicolumn{6}{c}{\textbf{Nichtmetalle und Halbmetalle (B, C, Si, Ge, As, Se, Te)}} \\
	B–C  & alle & 0,60 & C–Si  & alle & 0,65 \\
	Si–Ge & alle & 0,70 & Ge–As & alle & 0,68 \\
	As–Se & alle & 0,65 & Se–Te & alle & 0,63 \\
	\multicolumn{6}{c}{\textbf{Spezielle hochwirksame Systeme}} \\
	Fe–C  & \(\sigma_B\) & 1,15 & Fe–C  & \(H_V\) & 1,10 \\
	Fe–C  & TOF & 0,90 & Fe–C  & \(\sigma\) & 0,50 \\
	Ni–Al & \(\sigma_B\) & 1,00 & Ni–Al & \(H_V\) & 0,98 \\
	Ni–Al & TOF & 0,85 & Ni–Al & \(\sigma\) & 0,70 \\
	Cu–Sn & \(\sigma_B\) & 0,90 & Cu–Sn & \(H_V\) & 0,88 \\
	Cu–Sn & \(\sigma\) & 0,95 & & & \\
	Al–Cu & \(\sigma_B\) & 0,85 & Al–Cu & \(H_V\) & 0,82 \\
	Al–Cu & \(\sigma\) & 0,90 & & & \\
	Pt–Ru & TOF & 1,10 & Pd–Au & TOF & 0,95 \\
	\bottomrule
\end{longtable}
*Die vollständige Matrix für alle 2628 binären Systeme und alle Eigenschaften ist als CSV-Datei im Repository verfügbar.

\subsection{Einbeziehung von Phasendiagrammen und intermetallischer Bildung}
\label{subsec:phase-diagrams}

Das Vorhandensein intermetallischer Phasen oder Zweiphasengebiete beeinflusst die Materialeigenschaften erheblich. Im PLCM wird dies durch einen phasengewichteten Mittelwert erfasst:

\begin{equation}
	P_{\text{eff}} = \sum_{\alpha} w_{\alpha} \cdot P_{\alpha},
	\label{eq:phase-weighted}
\end{equation}

wobei \(\alpha\) die vorhandenen Phasen indiziert, \(w_{\alpha}\) der Massenanteil der Phase \(\alpha\) (aus dem Phasendiagramm) und \(P_{\alpha}\) die Eigenschaft dieser Phase ist.

\subsubsection{Intermetallischer Beitrag}
Wenn eine intermetallische Phase \(I\) (z. B. Ni\(_3\)Al, Cu\(_3\)Sn) entsteht, wird ein zusätzlicher Term zur Eigenschaftsvorhersage hinzugefügt:

\begin{equation}
	P_{\text{total}} = P_{\text{Mischkristall}} + \kappa_{ij}^{\text{IM}} \cdot c_i c_j \cdot (P_i - P_j)^2 \cdot \lambda_{\text{IM}}(T),
	\label{eq:intermetallic-term}
\end{equation}

wobei \(\lambda_{\text{IM}}(T)\) eine stufenartige Funktion ist, die unterhalb der Bildungstemperatur der intermetallischen Phase 1 und oberhalb 0 ist.

\subsection{Verarbeitungsempfindlichkeitskoeffizienten \(\kappa_{i,k}\)}
\label{subsec:processing-coefficients}

Die Empfindlichkeit eines Elements \(i\) gegenüber einer Verarbeitungsbehandlung \(k\) (z. B. Abschrecken, Kaltwalzen) ist im Koeffizienten \(\kappa_{i,k}\) codiert. In der PLCM-Vorhersageformel wird der Verarbeitungsbeitrag addiert:

\begin{equation}
	\Delta P_{\text{proc}} = \sum_{i=1}^{N} \sum_{k=126}^{130} \kappa_{i,k} \cdot w_i \cdot \Delta P_k,
	\label{eq:proc-contribution}
\end{equation}

wobei \(\Delta P_k\) die absolute Änderung der Eigenschaft \(P\) durch die Verarbeitung \(k\) ist (z. B. +400 MPa für das Abschrecken von Stahl). Der Koeffizient \(\kappa_{i,k}\) ist definiert als:

\begin{equation}
	\kappa_{i,k} = \frac{\Delta P_k^{(i)}}{\Delta P_k^{\text{ref}}},
	\label{eq:kappa-proc-def}
\end{equation}

wobei \(\Delta P_k^{(i)}\) die für das reine Element \(i\) unter derselben Verarbeitung beobachtete Eigenschaftsänderung ist und \(\Delta P_k^{\text{ref}}\) die Referenzänderung für ein Standardelement (z. B. Fe für Abschrecken) ist. Werte für Schlüsselelemente sind in Tabelle \ref{tab:kappa-proc} angegeben.

\begin{table}[htbp]
	\centering
	\caption{Verarbeitungsempfindlichkeitskoeffizienten \(\kappa_{i,k}\) für ausgewählte Elemente}
	\label{tab:kappa-proc}
	\begin{tabular}{lccc}
		\toprule
		Element & Abschrecken & Kaltwalzen (50\%) & Ausscheidungshärtung \\
		\midrule
		Fe & 1,0 & 1,0 & 0,8 \\
		Cu & 0,2 & 1,2 & 0,3 \\
		Al & 0,3 & 1,0 & 1,0 \\
		Ti & 0,5 & 0,8 & 1,2 \\
		Ni & 0,4 & 1,1 & 1,0 \\
		\bottomrule
	\end{tabular}
\end{table}

\subsection{Temperaturabhängigkeit von Eigenschaften}
\label{subsec:temperature-dependence}

Für Anwendungen bei erhöhten Temperaturen muss die Temperaturabhängigkeit der Eigenschaften einbezogen werden. Im PLCM geschieht dies durch Skalierungsfaktoren, die von der Koordinationseffizienz \(\Keff(T)\) abgeleitet werden.

\subsubsection{Skalierung mit der Koordinationseffizienz}

Aus Anhang P folgt, dass die Koordinationseffizienz eines Materials mit der Temperatur skaliert als:

\begin{equation}
	\Keff(T) = K_0 \left( \frac{T}{T_0} \right)^{-\gamma}, \quad \gamma \approx 0,3.
	\label{eq:keff-T}
\end{equation}

Die Eigenschaft \(P(T)\) bei der Temperatur \(T\) hängt mit ihrem Wert bei Raumtemperatur \(P_0\) zusammen durch:

\begin{equation}
	P(T) = P_0 \cdot \left( \frac{\Keff(T)}{\Keff(T_0)} \right)^{\xi_P},
	\label{eq:property-T-scaling}
\end{equation}

wobei \(\xi_P\) ein eigenschaftsspezifischer Exponent ist. Für Festigkeit und Härte gilt \(\xi \approx 0,2\), für elektrische Leitfähigkeit \(\xi \approx 0,5\), für katalytische Aktivität \(\xi \approx 0,3\).

\subsubsection{Effekte von Phasenübergangstemperaturen}

Nahe einem Phasenübergang (z. B. Curie-Temperatur \(T_C\), Schmelzpunkt \(T_m\)) ändert sich die Eigenschaft unstetig. Im PLCM wird dies durch die Funktion \(\lambda_{\text{Phase}}(T)\) erfasst, die unterhalb der Übergangstemperatur 1 und oberhalb 0 ist (oder umgekehrt). Die Gesamteigenschaft ist dann:

\begin{equation}
	P(T) = \lambda_{\text{Phase}}(T) \cdot P_{\text{niedrig}}(T) + (1-\lambda_{\text{Phase}}(T)) \cdot P_{\text{hoch}}(T),
	\label{eq:phase-transition}
\end{equation}

\subsection{Automatische Einbeziehung von Phasenübergängen und Temperatureffekten}
\label{subsec:auto-phase}

Um PLCM vollständig vorhersagfähig zu machen, ohne manuelle Phaseneingabe, erweitern wir die Eigenschaftsvorhersageformel um eine automatische Behandlung von Elementsprödigkeit, Phasenbildung und Temperaturskalierung. Der modifizierte Ausdruck lautet:

\begin{equation}
	\boxed{
		P = \sum_i w_i P_i^{\text{eff}}(T) \;+\; \delta_P \cdot \sum_{i<j} \kappa_{ij}^{(P)} w_i w_j (P_i - P_j)^2 \;+\; \Delta P_{\text{Phase}}^{\text{auto}} \;+\; \Delta P_{\text{Ent}} \;+\; \Delta P_{\text{proc}}
	}
	\label{eq:plcm-auto}
\end{equation}

\subsubsection{Effektive Reinlementfestigkeit}
Der Einfluss der Sprödigkeit wird durch einen Reduktionsfaktor $\eta_i(T)$ erfasst, der auf die Reinlementfestigkeit $P_i$ angewendet wird:

\begin{equation}
	P_i^{\text{eff}}(T) = P_i \cdot \eta_i(T), \qquad
	\eta_i(T) = 
	\begin{cases}
		1 - \dfrac{\max(0,\, T_{\text{spröde},i} - T)}{T_{\text{spröde},i}}, & T_{\text{spröde},i} > 0 \\
		1, & \text{sonst}
	\end{cases}
\end{equation}

\subsubsection{Automatischer Phasenbeitrag}
Der Phasenbeitrag wird aus einem eingebauten thermodynamischen Modul berechnet, das die Volumenanteile intermetallischer Phasen auf der Grundlage binärer Mischungsenthalpien, Atomgrößen und Elektronegativitätsdifferenzen schätzt.

\subsubsection{Hochen tropielegierungs-Entropieterm}
Wenn die Zusammensetzung fünf oder mehr Hauptelemente mit Konzentrationen zwischen 5 und 35 At.\% enthält, wird automatisch der konfigurative Entropiebeitrag hinzugefügt:

\begin{equation}
	\Delta P_{\text{Ent}} = \alpha_{\text{ent}} \cdot T \cdot \Delta S_{\text{conf}}, \quad
	\Delta S_{\text{conf}} = -R \sum_i w_i \ln w_i
\end{equation}
mit $\alpha_{\text{ent}} = 0,015$ MPa·mol·K\(^{-1}\)·J\(^{-1}\).

\subsection{Mischkristallverfestigungsterm}
\label{subsec:ss-term}

Für einphasige Mischkristalle ist der dominierende Verfestigungsmechanismus die Wechselwirkung von Versetzungen mit Fremdatomen. Um diesen Effekt automatisch zu erfassen, führen wir einen linearen Term $\Delta P_{\text{ss}}$ in der Eigenschaftsvorhersage ein:

\begin{equation}
	\Delta P_{\text{ss}} = \sum_i k_i^{(P)} \cdot w_i,
	\label{eq:ss-term}
\end{equation}

wobei $k_i^{(P)}$ der Mischkristallverfestigungskoeffizient des Elements $i$ für die Eigenschaft $P$ ist.

\subsection{Eigenschaftsklassen-Kalibrierung: Struktur- vs. Gittereigenschaften}
\label{subsec:property-class-calibration}

Eine kritische Verfeinerung des PLCM-Rahmens ist die Unterscheidung zwischen zwei grundlegend unterschiedlichen Klassen von Materialeigenschaften:

\begin{enumerate}
	\item \textbf{Klasse A (strukturempfindliche Eigenschaften)}: Diese Eigenschaften werden durch Defekte, Ausscheidungen und Mischkristallverfestigung dominiert. Legieren erzeugt nichtlineare Verstärkungen durch die Bildung intermetallischer Phasen, Versetzungswechselwirkungen und Ausscheidungshärtung. Beispiele: Zugfestigkeit \(\sigma_B\), Härte \(H_B\), Streckgrenze \(\sigma_y\), Ermüdungsgrenze, katalytische Aktivität TOF.
	
	\item \textbf{Klasse B (gitterdominierte Eigenschaften)}: Diese Eigenschaften werden hauptsächlich durch das Kristallgitter der Matrix bestimmt. Legieren im Mischkristall ergibt näherungsweise lineare Beiträge (Mischungsregel) mit minimaler nichtlinearer Verstärkung. Beispiele: Elastizitätsmodul \(E\), Schubmodul \(G\), Poissonzahl \(\nu\), Wärmeleitfähigkeit \(\lambda\), Wärmeausdehnungskoeffizient \(\alpha_T\), Dichte \(\rho\), spezifische Wärme \(c_p\).
\end{enumerate}

\subsection{Modifizierte PLCM-Vorhersageformel}

Die allgemeine PLCM-Vorhersageformel wird mit einem eigenschaftsspezifischen Koeffizienten \(\delta_P\) modifiziert:

\begin{equation}
	\boxed{P = \sum_{i=1}^{N} x_i P_i + \delta_P \cdot \sum_{1\le i<j\le N} \kappa_{ij}^{(P)} x_i x_j (P_i - P_j)^2}
	\label{eq:plcm-modified}
\end{equation}

wobei:
\[
\delta_P = \begin{cases}
	1,0 & \text{Klasse A Eigenschaften (Festigkeit, Härte, TOF)} \\
	0,0 & \text{Klasse B Eigenschaften (Modul, Leitfähigkeit, CTE, Dichte)} \\
	0,1\text{--}0,3 & \text{Intermediäre Eigenschaften (elektrische Leitfähigkeit, magnetische Suszeptibilität)}
\end{cases}
\]

\subsection{Eigenschaftsklassifizierungstabelle}

\begin{longtable}{p{4cm}p{3cm}p{3cm}p{3cm}}
	\caption{Eigenschaftsklassifizierung für die PLCM-Kalibrierung} \label{tab:property-classification} \\
	\toprule
	\textbf{Eigenschaft} & \textbf{Symbol} & \textbf{Klasse} & \(\delta_P\) \\
	\midrule
	\endfirsthead
	\toprule
	\textbf{Eigenschaft} & \textbf{Symbol} & \textbf{Klasse} & \(\delta_P\) \\
	\midrule
	\endhead
	Zugfestigkeit & \(\sigma_B\) & A & 1,0 \\
	Streckgrenze & \(\sigma_y\) & A & 1,0 \\
	Härte (Brinell, Vickers) & \(H_B, H_V\) & A & 1,0 \\
	Ermüdungsgrenze & \(\sigma_{-1}\) & A & 1,0 \\
	Katalytische Aktivität & TOF & A & 1,0 \\
	\hline
	Elastizitätsmodul & \(E\) & B & 0,0 \\
	Schubmodul & \(G\) & B & 0,0 \\
	Poissonzahl & \(\nu\) & B & 0,0 \\
	Dichte & \(\rho\) & B & 0,0 \\
	Wärmeleitfähigkeit & \(\lambda\) & B & 0,0 \\
	Wärmeausdehnungskoeffizient & \(\alpha_T\) & B & 0,0 \\
	Spezifische Wärme & \(c_p\) & B & 0,0 \\
	\hline
	Elektrische Leitfähigkeit & \(\sigma\) & Intermediär & 0,2 \\
	Magnetische Suszeptibilität & \(\chi_m\) & Intermediär & 0,2 \\
	\bottomrule
\end{longtable}

\subsection{Materialklassen-Kalibrierungsfaktoren \(\gamma_{\text{Klasse}}\)}
\label{subsec:material-class-factors}

Die universellen Kalibrierungsfaktoren \(\beta_{ij}^{(P)}\) werden durch materialklassenspezifische Koeffizienten \(\gamma_{\text{Klasse}}\) weiter verfeinert, die Unterschiede in den Verfestigungsmechanismen über Legierungsfamilien hinweg berücksichtigen. Die vollständige PLCM-Vorhersageformel lautet dann:

\begin{equation}
	\boxed{P = \sum_{i=1}^{N} x_i P_i + \gamma_{\text{Klasse}} \cdot \delta_P \cdot \sum_{1\le i<j\le N} \beta_{ij}^{(P)} \kappa_{ij}^{\text{(Basis)}} x_i x_j (P_i - P_j)^2 + \Delta P_{\text{Phase}}}
	\label{eq:plcm-full-calibration}
\end{equation}

\begin{longtable}{p{3.5cm}p{2.5cm}p{2.5cm}p{5cm}}
	\caption{Materialklassen-Kalibrierungsfaktoren \(\gamma_{\text{Klasse}}\) für Festigkeitseigenschaften} \label{tab:material-class-factors} \\
	\toprule
	\textbf{Materialklasse} & \(\gamma_{\sigma}\) & \(\gamma_{E}\) & \textbf{Physikalische Grundlage} \\
	\midrule
	\endfirsthead
	\toprule
	\textbf{Materialklasse} & \(\gamma_{\sigma}\) & \(\gamma_{E}\) & \textbf{Physikalische Grundlage} \\
	\midrule
	\endhead
	Aluminiumlegierungen (Al-Cu, Al-Mg, Al-Si, Al-Zn) & 0,85 & 1,0 & Ausscheidungshärtung (GP-Zonen, \(\theta'\), \(\theta\), \(\eta'\)) \\
	Stähle (Fe-C, Fe-Cr, Fe-Mn, Fe-Ni) & 1,15 & 1,0 & Martensitische Umwandlung + Carbidausscheidung \\
	Titanlegierungen (Ti-Al-V, Ti-Al-Sn) & 1,05 & 1,0 & \(\alpha/\beta\)-Phasenumwandlung \\
	Nickel-Superlegierungen (Ni-Cr-Al, Ni-Cr-Co) & 1,10 & 1,0 & \(\gamma'\) (Ni\(_3\)Al)-Ausscheidungshärtung \\
	Kupferlegierungen (Cu-Zn, Cu-Sn, Cu-Ni) & 0,90 & 1,0 & Mischkristall + intermetallische Phasen \\
	Magnesiumlegierungen (Mg-Al-Zn, Mg-Zn-Zr) & 0,80 & 1,0 & Begrenzte Ausscheidungshärtung \\
	\bottomrule
\end{longtable}

\subsection{Phasenspezifische Beiträge \(\Delta P_{\text{Phase}}\)}
\label{subsec:phase-contributions}

Für Materialien, die während der Wärmebehandlung Phasenumwandlungen durchlaufen, wird ein zusätzlicher phasenspezifischer Beitrag \(\Delta P_{\text{Phase}}\) zur Festigkeitsvorhersage hinzugefügt.

\begin{longtable}{p{3.5cm}p{3.0cm}p{3.0cm}p{4cm}}
	\caption{Phasenspezifische Beiträge zur Festigkeit \(\Delta\sigma_{\text{Phase}}\) (MPa)} \label{tab:phase-contributions} \\
	\toprule
	\textbf{Materialklasse} & \textbf{Phase/Struktur} & \(\Delta\sigma_{\text{Phase}}\) (MPa) & \textbf{Mechanismus} \\
	\midrule
	\endfirsthead
	\toprule
	\textbf{Materialklasse} & \textbf{Phase/Struktur} & \(\Delta\sigma_{\text{Phase}}\) (MPa) & \textbf{Mechanismus} \\
	\midrule
	\endhead
	\multicolumn{4}{c}{\textbf{Stähle}} \\
	& Ferrit & 0 & Grundmatrix \\
	& Perlit (fein) & 250-350 & Lamellenstruktur, Fe\(_3\)C-Platten \\
	& Bainit (unterer) & 400-600 & Nadeliger Ferrit + feine Carbide \\
	& Martensit (niedriger C, \(<0,3\%\)) & 800-1000 & Lattenmartensit, Versetzungen \\
	& Martensit (mittlerer C, 0,3-0,6\%) & 1000-1300 & Gemischter Latten-/Plattenmartensit \\
	& Martensit (hoher C, \(>0,6\%\)) & 1300-1600 & Plattenmartensit, Zwillingsbildung \\
	& Angelassener Martensit & 600-1000 & Feine Carbide in ferritischer Matrix \\
	\hline
	\multicolumn{4}{c}{\textbf{Aluminiumlegierungen}} \\
	& Gusszustand / Mischkristall & 0 & Keine Ausscheidungen \\
	& T4 (natürliche Alterung) & 150-250 & GP-Zonen \\
	& T6 (künstliche Alterung) & 250-400 & \(\theta'\) (Al\(_2\)Cu), \(\eta'\) (MgZn\(_2\))-Ausscheidungen \\
	& T7 (Überalterung) & 150-250 & Größere Ausscheidungen \\
	\hline
	\multicolumn{4}{c}{\textbf{Kupferlegierungen}} \\
	& Mischkristall & 0 & \\
	& Ausscheidungsgehärtet & 150-300 & Intermetallische Phasen (Cu\(_3\)Sn, Cu\(_3\)Al) \\
	\hline
	\multicolumn{4}{c}{\textbf{Titanlegierungen}} \\
	& \(\alpha\)-Phase & 0 & HCP-Struktur \\
	& \(\beta\)-Phase & 50-100 & Krz-Struktur \\
	& \(\alpha+\beta\) gealtert & 200-400 & Feine \(\alpha\)-Ausscheidungen in \(\beta\)-Matrix \\
	\hline
	\multicolumn{4}{c}{\textbf{Nickel-Superlegierungen}} \\
	& Mischkristall & 0 & \\
	& Gealtert (\(\gamma'\)-Ausscheidung) & 300-600 & Kohärente Ni\(_3\)Al-Ausscheidungen \\
	\hline
	\multicolumn{4}{c}{\textbf{Magnesiumlegierungen}} \\
	& Mischkristall & 0 & \\
	& Gealtert (Ausscheidung) & 50-150 & Mg\(_{17}\)Al\(_{12}\)-Ausscheidungen \\
	\bottomrule
\end{longtable}

\subsection{Zweiphasen-Legierungsmodul}
\label{subsec:two-phase-alloys}

Für Legierungen, die intermetallische Phasen bilden (Cu-Sn, Al-Si, Fe-C usw.), versagt das Standardmodell der Mischkristallverfestigung, da das Legierungselement nicht vollständig in der Matrix gelöst ist. Stattdessen muss die Eigenschaftsvorhersage auf Phasenanteilen aus dem Gleichgewichtsphasendiagramm basieren.

\subsubsection{Allgemeine Formel für Zweiphasenlegierungen}

\begin{equation}
	\boxed{P = \sum_{\alpha} f_{\alpha} P_{\alpha} + \sum_{\beta} \kappa_{\alpha\beta}^{\text{IM}} \cdot f_{\alpha} f_{\beta} \cdot (P_{\alpha} - P_{\beta})^2}
	\label{eq:two-phase-alloy}
\end{equation}

Für die meisten Zweiphasenlegierungen ist jedoch die einfachere Mischungsregel mit einem kleinen Dispersionsverfestigungsterm angemessener:

\begin{equation}
	\boxed{P = \sum_{\alpha} f_{\alpha} P_{\alpha} + \Delta P_{\text{disp}}}
	\label{eq:two-phase-simple}
\end{equation}

\subsubsection{Kalibrierung für das Cu-Sn-System (Glockenbronze)}
Für Cu-22Sn (Glockenbronze) mit \(\alpha\)-Phase (Cu-reicher Mischkristall, 78 Vol.-\%, \(\sigma_B=220\) MPa) und \(\delta\)-Phase (Cu\(_{41}\)Sn\(_{11}\), 22 Vol.-\%, \(\sigma_B=400\) MPa) ergibt sich mit \(\Delta P_{\text{disp}}=50\) MPa:

\[
\sigma_B = 0,78 \cdot 220 + 0,22 \cdot 400 + 50 = 309,6\ \text{MPa},
\]
was gut mit dem experimentellen Bereich von 220–280 MPa übereinstimmt.

\subsection{Versetzungs substrukturverfestigung}
\label{subsec:dislocation-substructure}

Während der plastischen Verformung organisieren sich Versetzungen selbst in fraktale Muster (Zellen, Wände, Subkörner). Der Beitrag dieser Substruktur zur Streckgrenze kann mit den universellen YUCT-Konstanten modelliert werden:

\[
\Delta\sigma_{\text{sub}} = \alpha G b \sqrt{\rho} \cdot \left( \frac{S_{\text{odd}}}{S_{\text{even}}} \right)^{\eta},
\]
mit \(\alpha \approx 0,3\) (Taylor-Faktor), \(G\) Schubmodul, \(b\) Burgersvektor, \(\rho\) Versetzungsdichte und \(\eta \approx 0,5\). Das Verhältnis \(S_{\text{odd}}/S_{\text{even}}=1,5\) berücksichtigt die hierarchische Natur des Versetzungsnetzwerks.

\subsection{Vollständige Kalibrierungstabelle}
\label{subsec:complete-calibration}

\begin{longtable}{p{3cm}p{2.5cm}p{2.5cm}p{2.5cm}p{4cm}}
	\caption{Vollständige Kalibrierungsfaktoren \(\gamma\) für alle Legierungsklassen} \label{tab:complete-gamma} \\
	\toprule
	\textbf{Legierungsklasse} & \textbf{System} & \(\gamma_{\text{geglüht}}\) & \(\gamma_{\text{gehärtet}}\) & \textbf{Anmerkung} \\
	\midrule
	\endfirsthead
	\toprule
	\textbf{Legierungsklasse} & \textbf{System} & \(\gamma_{\text{geglüht}}\) & \(\gamma_{\text{gehärtet}}\) & \textbf{Anmerkung} \\
	\midrule
	\endhead
	Aluminium & Al-Cu & 0,85 & 0,85 & Duralumin-Typ \\
	Aluminium & Al-Si & 2,5 & 2,5 & Eutektische Silumine \\
	Aluminium & Al-Zn-Mg & 0,90 & 0,90 & Typ 7075 \\
	Stahl & Fe-C & 1,15 & 1,15 + \(\Delta\sigma_{\text{mart}}\) & Wälzlagerstahl \\
	Stahl & Fe-Cr-Ni & 2,0 & 3,5 & Austenitischer Edelstahl \\
	Kupfer & Cu-Zn (\(\alpha\)) & 1,2 & 1,2 & Messing, einphasig \\
	Kupfer & Cu-Zn (\(\alpha+\beta\)) & — & Mischungsregel & Zweiphasiges Messing \\
	Kupfer & Cu-Sn & — & Mischungsregel & Bronze \\
	Kupfer & Cu-Be & 2,5 & 5,5 & Berylliumkupfer \\
	Titan & Ti-6Al-4V & 2,5 & 3,0 & \\
	Nickel & Ni-Cr-Al & 2,5 & 3,5 & Superlegierungen \\
	Magnesium & Mg-Al-Zn & 1,7 & 2,2 & AZ91 \\
	\bottomrule
\end{longtable}

\subsection{Kalibrierungsbeispiele}

\begin{longtable}{lccccc}
	\caption{Validierung der PLCM-Kalibrierung an Duralumin und Wälzlagerstahl} \label{tab:calibration-validation} \\
	\toprule
	\textbf{Material} & \textbf{Zustand} & \(\gamma_{\text{Klasse}}\) & \(\Delta\sigma_{\text{Phase}}\) (MPa) & \textbf{PLCM vorhergesagt} & \textbf{Experimentell} \\
	\midrule
	\endfirsthead
	\toprule
	\textbf{Material} & \textbf{Zustand} & \(\gamma_{\text{Klasse}}\) & \(\Delta\sigma_{\text{Phase}}\) (MPa) & \textbf{PLCM vorhergesagt} & \textbf{Experimentell} \\
	\midrule
	\endhead
	Duralumin (Al-4Cu-1,5Mg-0,6Mn) & T6 (gealtert) & 0,85 & 300 & 468 & 430-490 \\
	Wälzlagerstahl (Fe-1C-1,5Cr) & Normalisiert (Perlit) & 1,15 & 300 & 592 & 650-800 \\
	Wälzlagerstahl (Fe-1C-1,5Cr) & Abgeschreckt + angelassen & 1,15 & 1200 & 2043 & 1800-2200 \\
	\bottomrule
\end{longtable}

% ============================================================================
% TEIL 5: THEORETISCHE GRUNDLAGEN UND FORTGESCHRITTENER LAGRANGE-RAHMEN
% Einzufügen nach dem Ende von Teil 4 (nach der Tabelle \end{longtable} der Kalibrierungsbeispiele)
% ============================================================================

\subsection{Theoretische Grundlagen}

Die Vorhersagekraft von PLCM beruht auf der Koordinationseffizienz $K_{\mathrm{eff}}$,
einem universellen Parameter, der alle Phasenübergänge und chemischen Reaktionen bestimmt.
Wie in Anhang C2 streng gezeigt wurde, folgen die Schmelz- und Siedetemperaturen reiner
Elemente $T \propto K_{\mathrm{eff}}^{0,67}$, eine direkte Konsequenz des universellen
Fehlergesetzes von YUCT. Dieses gleiche Gesetz verbindet die kritischen Punkte von
Phasenübergängen mit den Konstanten der Chaostheorie (Feigenbaum $\alpha,\delta$)
und mit der Entropieproduktionsschwelle dissipativer Strukturen (Prigogine).
Daher ist jede von PLCM berechnete Materialeigenschaft letztlich in der algebraischen
Schleife $S_{\mathrm{odd}}=1,2$, $S_{\mathrm{even}}=0,8$ und der zugehörigen
geometrischen Abweichung $\Delta\pi \approx 4,8\times10^{-5}$ verwurzelt.
Siehe Anhang C2 für die vollständige Herleitung.

\section{Fundamentaler Lagrange-Rahmen für fortgeschrittene Materialien}
\label{sec:lagrangian}

Das in Abschnitt~\ref{subsec:purity} vorgestellte empirische Modell ist ein Spezialfall eines allgemeineren Variationsansatzes, der auf einem Lagrange-Funktional basiert. Dieser Rahmen ermöglicht die Beschreibung von Materialien, deren Festigkeit durch Koordinationschemie, katalytische Effekte und resonante externe Felder bestimmt wird, während er vollständig mit der bestehenden sektorbasierten Datenbankstruktur kompatibel bleibt.

\subsection{Lagrange-Dichte und Feldvariablen}
\label{subsec:lagrangian-density}

Wir definieren eine Lagrange-Dichte \(\mathcal{L}\) integriert über das Materialvolumen \(\Omega\):

\[
\mathcal{L} = \int_{\Omega} \Big[ 
\psi_{\mathrm{el}}(\boldsymbol{\varepsilon}, \mathbf{c}, \mathbf{q}) 
+ \psi_{\mathrm{coord}}(\mathbf{q}, \nabla\mathbf{q}) 
+ \psi_{\mathrm{cat}}(\mathbf{c}_{\mathrm{cat}}, \dot{\mathbf{q}}, \dot{\mathbf{c}}) 
+ \psi_{\mathrm{res}}(\mathbf{E}, \mathbf{c}, \mathbf{q}, \omega) 
+ \psi_{\mathrm{diss}}(\dot{\boldsymbol{\varepsilon}}, \dot{\mathbf{q}}, \dot{\mathbf{c}})
\Big] dV
\]

wobei:
\begin{itemize}
\item \(\boldsymbol{\varepsilon}\) – Verzerrungstensor,
\item \(\mathbf{c}\) – Vektor der Konzentrationen aller Elemente,
\item \(\mathbf{q}\) – Ordnungsparameter, die die lokale Koordinationsgeometrie beschreiben,
\item \(\mathbf{c}_{\mathrm{cat}}\) – Konzentrationen katalytisch aktiver Spezies,
\item \(\mathbf{E}\) – externes elektrisches Feld (Frequenz \(\omega\)),
\item \(\psi_{\mathrm{el}}\) – elastische Energie (enthält Mischkristall- und Ausscheidungsbeiträge),
\item \(\psi_{\mathrm{coord}}\) – Energiekosten für räumliche Variationen der Koordination,
\item \(\psi_{\mathrm{cat}}\) – kinetische Terme, die durch Katalysatoren beeinflusst werden,
\item \(\psi_{\mathrm{res}}\) – resonante Feld-Materie-Kopplung,
\item \(\psi_{\mathrm{diss}}\) – Dissipation (Plastizität, Viskosität).
\end{itemize}

Alle materialspezifischen Parameter werden in nummerierten Sektoren der PLCM-Datenbank gespeichert. Die Sektoren 130–136 sind für diesen erweiterten Rahmen reserviert.

\subsection{Datenbanksektoren für Lagrange-Parameter}
\label{subsec:sectors}

\begin{table}[htbp]
\centering
\caption{PLCM-Sektoren für den Lagrange-Rahmen}
\label{tab:sectors}
\begin{tabular}{clp{8cm}}
	\toprule
	\textbf{Sektor} & \textbf{Inhalt} & \textbf{Beschreibung} \\
	\midrule
	130 & Materialtyp & Flags: Strukturlegierung / Koordinationsmaterial / Hybrid; Kristallsystem; usw. \\
	131 & Elastische Parameter & \(C_{ijkl}(\mathbf{c}, \mathbf{q})\), Mischkristallkoeffizienten, Ausscheidungsbeiträge. \\
	132 & Empirisches Verunreinigungsmodell & \(\alpha_{\mathrm{imp}}\), \(\sigma_B^{\mathrm{base}}\); wird verwendet, wenn \(\mathbf{q}\) eingefroren ist. \\
	133 & Koordinationsenergie & Potential \(V(\mathbf{c}, q_1,q_2,q_3)\), Gradientenkoeffizienten \(A,B,C\); Gleichgewichtsbindungslängen, Koordinationszahlen, Polyedertypen. \\
	134 & Katalytische Kinetik & Funktionen \(\gamma_i(\mathbf{c}_{\mathrm{cat}})\), \(\delta_j(\mathbf{c}_{\mathrm{cat}})\), die die Katalysatorkonzentration mit Relaxationsraten verknüpfen. \\
	135 & Resonante Antwort & Dielektrische Suszeptibilität \(\boldsymbol{\chi}(\mathbf{c},\mathbf{q},\omega)\), Elektrostriktionstensor \(\mathbf{M}\), Resonanzfrequenzen \(\omega_0\), Dämpfung \(\tau\). \\
	136 & Dissipation & Viskositätskoeffizienten, Fließkriterien, Schadensparameter. \\
	\bottomrule
\end{tabular}
\end{table}

Die Sektoren 131–136 können Skalarwerte, Tensoren oder Funktionsanpassungen (z. B. Polynome in \(\mathbf{c}\) und \(\mathbf{q}\)) enthalten. Für konventionelle Strukturlegierungen werden die Sektoren 133–135 typischerweise auf Null oder triviale Werte gesetzt, und Sektor 132 liefert das vereinfachte additive Modell.

\subsection{Koordinationsmaterialien (\(\psi_{\mathrm{coord}}\))}
\label{subsec:coordination}

Für Materialien, bei denen die Festigkeit von gerichteten kovalenten oder Koordinationsbindungen herrührt, beschreiben die Ordnungsparameter \(\mathbf{q}\) die lokale atomare Anordnung. Eine typische Landau-Ginzburg-Form lautet:

\[
\psi_{\mathrm{coord}} = A(\mathbf{c})\,(\nabla q_1)^2 + B(\mathbf{c})\,(\nabla q_2)^2 + C(\mathbf{c})\,(\nabla q_3)^2 + V(\mathbf{c}, q_1, q_2, q_3)
\]

wobei:
\begin{itemize}
\item \(q_1\) – Typ des Koordinationspolyeders (z. B. 0 Tetraeder, 1 Oktaeder, 2 Würfel),
\item \(q_2\) – mittlere Bindungslänge,
\item \(q_3\) – Fernordnungsparameter.
\end{itemize}

Das Potential \(V\) hat Minima bei den Gleichgewichtswerten für eine gegebene Zusammensetzung. Die Koeffizienten \(A,B,C\) kontrollieren die Grenzflächenenergie zwischen Domänen mit unterschiedlicher Koordination.

Die mechanische Festigkeit ergibt sich aus der zweiten Variation der Gesamtenergie nach \(\boldsymbol{\varepsilon}\). Für einen perfekten Kristall einer Koordinationsverbindung (z. B. TiC) kann die vorhergesagte Festigkeit direkt aus \(\psi_{\mathrm{el}}\) extrahiert werden; für polykristalline oder Verbundmaterialien wird eine Homogenisierung über die \(\mathbf{q}\)-Felder durchgeführt.

\subsection{Katalytische Effekte (\(\psi_{\mathrm{cat}}\))}
\label{subsec:catalysis}

Katalytische Elemente (z. B. Pt, Pd, Ni in kleinen Mengen) senken die Aktivierungsbarrieren für Phasenumwandlungen und die Mikrostrukturentwicklung. Im Lagrange-Operator modifizieren sie die kinetischen Koeffizienten:

\[
\psi_{\mathrm{cat}} = \sum_i \gamma_i(\mathbf{c}_{\mathrm{cat}})\, \dot{q}_i^2 + \sum_j \delta_j(\mathbf{c}_{\mathrm{cat}})\, \dot{c}_j^2
\]

Die Funktionen \(\gamma_i, \delta_j\) sind typischerweise fallend mit der Katalysatorkonzentration, was eine schnellere Relaxation widerspiegelt. Beispiel: \(\gamma_i(\mathbf{c}_{\mathrm{cat}}) = \gamma_i^0 / (1 + \beta_i c_{\mathrm{cat}})\). Die Werte von \(\gamma_i^0\) und \(\beta_i\) sind in Sektor 134 gespeichert.

Dieser Term beeinflusst die Entwicklung von \(\mathbf{q}\) und \(\mathbf{c}\) während der Verarbeitung und damit die endgültige Festigkeit über den Verarbeitungsbeitrag \(\Delta\sigma_{\mathrm{proc}}\).

\subsection{Resonante Effekte (\(\psi_{\mathrm{res}}\))}
\label{subsec:resonance}

Wenn das Material einem oszillierenden elektrischen Feld \(\mathbf{E}(t) = \mathbf{E}_0 \cos(\omega t)\) ausgesetzt wird, lautet der resonante Teil des Lagrange-Operators:

\[
\psi_{\mathrm{res}} = -\frac{1}{2}\varepsilon_0\, \mathbf{E} \cdot \boldsymbol{\chi}(\mathbf{c}, \mathbf{q}, \omega) \cdot \mathbf{E}
\]

Für eine Lorentz-artige Resonanz:

\[
\chi(\omega) = \frac{\chi_0}{1 - (\omega/\omega_0)^2 + i\omega/\tau}
\]

wobei \(\omega_0\) die Plasmon- oder Phonon-Resonanzfrequenz, \(\tau\) die Relaxationszeit und \(\chi_0\) die statische Suszeptibilität ist. Diese Parameter sind in Sektor 135 gespeichert.

Zusätzlich kann das elektrische Feld über Elektrostriktion an die Verzerrung koppeln:

\[
\psi_{\mathrm{el}}^{\mathrm{coup}} = \frac{1}{2} \boldsymbol{\varepsilon} : \mathbf{M} : \mathbf{E}\mathbf{E}
\]

wobei \(\mathbf{M}\) der Elektrostriktionstensor ist. Dieser Term modifiziert die effektiven elastischen Konstanten unter resonanten Bedingungen und führt zu einer frequenzabhängigen Verschiebung von Festigkeit und Steifigkeit.

\subsection{Beziehung zum empirischen Verunreinigungsmodell}
\label{subsec:relation-empirical}

Für konventionelle Strukturlegierungen, bei denen die Koordinationsordnungsparameter fest sind (\(\mathbf{q} = \mathbf{q}_0\)) und katalytische/resonante Terme vernachlässigbar sind, reduziert sich der Lagrange-Operator auf:

\[
\mathcal{L}_{\mathrm{struct}} = \int_{\Omega} \psi_{\mathrm{el}}(\boldsymbol{\varepsilon}, \mathbf{c}) \, dV + \text{(Dissipation)}
\]

Durch Entwicklung von \(\psi_{\mathrm{el}}\) um einen Referenzzustand und Annahme einer linearen Superposition der Verfestigungsmechanismen erhalten wir die in Abschnitt~\ref{subsec:purity} verwendete additive Formel zurück:

\[
\sigma_B = \sigma_B^{\mathrm{base}} + \Delta\sigma_{\mathrm{ss}} + \Delta\sigma_{\mathrm{phase}} + \Delta\sigma_{\mathrm{proc}} + \underbrace{\sigma_B^{\mathrm{base}} \cdot \alpha_{\mathrm{imp}} (1 - q_{\mathrm{purity}})}_{\text{Verunreinigungskorrektur}}
\]

Somit vereinheitlicht der Lagrange-Rahmen die Beschreibung sowohl traditioneller als auch fortgeschrittener Materialien innerhalb einer einzigen mathematischen Struktur.

\subsection{Beispiel: TiC als Koordinationsmaterial}
\label{subsec:example-tic}

Titancarbid (TiC) ist ein prototypisches Koordinationsmaterial. Seine Lagrange-Parameter (Sektor 133) sind:

\begin{itemize}
\item Gleichgewichtskoordination: oktaedrisch (\(q_1=1\)), Bindungslänge \(q_2 = 0,216\) nm, perfekte Ordnung \(q_3=1\).
\item Tiefe des Potentialminimums: \(V_0 = 12,4\) eV pro Formeleinheit (aus DFT).
\item Gradientenkoeffizienten: \(A=0,5\) GPa·nm², \(B=0,2\) GPa·nm², \(C=0,1\) GPa·nm².
\end{itemize}

Sektor 131 liefert die elastischen Konstanten: \(C_{11}=500\) GPa, \(C_{12}=113\) GPa, \(C_{44}=175\) GPa. Die aus dem Lagrange-Operator berechnete theoretische Zugfestigkeit beträgt \(\sigma_B \approx 1200\) MPa, in guter Übereinstimmung mit experimentellen Werten für dichtes TiC.

Für TiC sind unter Standardbedingungen keine katalytischen oder resonanten Effekte aktiv, daher enthalten die Sektoren 134–135 Null- oder neutrale Werte.

\subsection{Benutzerablauf für fortgeschrittene Materialien}
\label{subsec:workflow}

Um ein neues Material in diesem Rahmen zu modellieren:

\begin{enumerate}
\item Wählen Sie den Materialtyp (Struktur-, Koordinations- oder Hybridmaterial) in Sektor 130.
\item Geben Sie die Zusammensetzung \(\mathbf{c}\) an.
\item Falls bekannt, tragen Sie die Lagrange-Koeffizienten aus der Literatur oder aus Ab-initio-Rechnungen direkt in die entsprechenden Sektoren ein.
\item Falls unbekannt, verwenden Sie den eingebauten Schätzer, der Sektorwerte aus atomaren Eigenschaften (Elektronegativität, Atomradius usw.) unter Verwendung von Regressionsmodellen vorhersagt, die auf vorhandenen Daten trainiert wurden.
\item Geben Sie alle externen Feldparameter (Frequenz, Amplitude) für resonante Berechnungen an.
\item Führen Sie den Variationslöser aus, um die Gleichgewichts-\(\mathbf{q}\)-Felder zu erhalten, und berechnen Sie die mechanischen Eigenschaften als Ableitungen der Gesamtenergie.
\end{enumerate}

Dieser Ansatz ermöglicht die systematische Erforschung neuer Materialien, einschließlich solcher, die auf Koordinationschemie, katalytisch verstärkten Legierungen und feldabstimmbaren Verbundwerkstoffen basieren.


\subsection{Hinweise zur Implementierung}
\label{subsec:implementation}

Der Lagrange-Rahmen ist als Erweiterung des PLCM-Kerns implementiert. Die Sektordaten werden in einer HDF5-Datei (oder einer relationalen Datenbank) mit einer einheitlichen API gespeichert. Der Variationslöser verwendet Finite-Elemente für die Feldgleichungen, die sich aus den Euler-Lagrange-Gleichungen ergeben:

\[
\frac{\partial \mathcal{L}}{\partial \phi} - \nabla \cdot \frac{\partial \mathcal{L}}{\partial (\nabla\phi)} = 0
\]

für jedes Feld \(\phi = (\boldsymbol{\varepsilon}, \mathbf{c}, \mathbf{q})\). Zeitabhängige Terme werden mit einem gestaffelten implizit-expliziten Schema behandelt.

Für konventionelle Strukturlegierungen wird der vollständige Lagrange-Operator automatisch auf das empirische Modell vereinfacht, wodurch die Rückwärtskompatibilität erhalten bleibt.

% ============================================================================
% TEIL 6: KOPPLUNGSMATRIX, KALIBRIERUNG, TOF, ARBEITSABLAUF, VORHERSAGE
% Einzufügen nach Teil 5 (nach \subsection{Benutzerablauf für fortgeschrittene Materialien})
% ============================================================================

\subsection{Vollständige Kopplungsmatrix \(\kappa_{ij}\) für alle 118 Elemente}
\label{subsec:full-kappa}

Die vollständige \(118\times118\)-Matrix \(\kappa_{ij}\) ist symmetrisch (\(\kappa_{ij}=\kappa_{ji}\)). Aus Platzgründen wird hier nur ein repräsentatives Fragment für die ersten 20 Elemente abgedruckt. Die vollständige Matrix ist als CSV-Datei im Repository verfügbar \url{https://github.com/Alexey-Yakushev-YUCT/YPSDC/}. Alle Werte werden mit der folgenden Formel berechnet:

\[
\kappa_{ij} = \frac{2\,\Delta H_{\text{mix}}^{ij}}{\Delta H_{\text{mix,ref}}} \cdot
\frac{1}{1 + |r_i - r_j|/r_{\text{avg}}} \cdot
\frac{1}{1 + |\chi_i - \chi_j|},
\]

mit \(\Delta H_{\text{mix,ref}} = -7,5\) kJ/mol (Referenzsystem Cu–Sn). Für Systeme, bei denen experimentelle \(\Delta H_{\text{mix}}\) nicht verfügbar sind, wird das Miedema-Modell \cite{Miedema1980} verwendet.

\subsubsection*{Fragment: Erste 20 Elemente}
Tabelle \ref{tab:kappa-full-20} zeigt den oberen Dreiecksteil von \(\kappa_{ij}\) für \(Z=1\) bis \(20\).

{\tiny
	\begin{longtable}{c|*{20}{c}}
		\caption{Oberer Dreiecksteil von \(\kappa_{ij}\) (erste 20 Elemente)} \label{tab:kappa-full-20} \\
		\toprule
		& 1 & 2 & 3 & 4 & 5 & 6 & 7 & 8 & 9 & 10 & 11 & 12 & 13 & 14 & 15 & 16 & 17 & 18 & 19 & 20 \\
		\midrule
		\endfirsthead
		\toprule
		& 1 & 2 & 3 & 4 & 5 & 6 & 7 & 8 & 9 & 10 & 11 & 12 & 13 & 14 & 15 & 16 & 17 & 18 & 19 & 20 \\
		\midrule
		\endhead
		1  & – & 0,00 & 0,00 & 0,00 & 0,00 & 0,00 & 0,00 & 0,00 & 0,00 & 0,00 & 0,00 & 0,00 & 0,00 & 0,00 & 0,00 & 0,00 & 0,00 & 0,00 & 0,00 & 0,00 \\
		2  &   & – & 0,00 & 0,00 & 0,00 & 0,00 & 0,00 & 0,00 & 0,00 & 0,00 & 0,00 & 0,00 & 0,00 & 0,00 & 0,00 & 0,00 & 0,00 & 0,00 & 0,00 & 0,00 \\
		3  &   &   & – & 0,45 & 0,30 & 0,20 & 0,15 & 0,12 & 0,10 & 0,00 & 0,50 & 0,55 & 0,48 & 0,35 & 0,30 & 0,25 & 0,20 & 0,00 & 0,52 & 0,48 \\
		4  &   &   &   & – & 0,55 & 0,48 & 0,40 & 0,35 & 0,30 & 0,00 & 0,48 & 0,55 & 0,58 & 0,52 & 0,45 & 0,40 & 0,35 & 0,00 & 0,50 & 0,55 \\
		5  &   &   &   &   & – & 0,62 & 0,55 & 0,50 & 0,45 & 0,00 & 0,42 & 0,48 & 0,55 & 0,60 & 0,58 & 0,52 & 0,48 & 0,00 & 0,44 & 0,50 \\
		6  &   &   &   &   &   & – & 0,68 & 0,62 & 0,58 & 0,00 & 0,35 & 0,42 & 0,48 & 0,55 & 0,60 & 0,62 & 0,58 & 0,00 & 0,38 & 0,45 \\
		7  &   &   &   &   &   &   & – & 0,70 & 0,65 & 0,00 & 0,30 & 0,38 & 0,42 & 0,48 & 0,55 & 0,58 & 0,62 & 0,00 & 0,32 & 0,40 \\
		8  &   &   &   &   &   &   &   & – & 0,72 & 0,00 & 0,25 & 0,32 & 0,38 & 0,42 & 0,48 & 0,52 & 0,58 & 0,00 & 0,28 & 0,35 \\
		9  &   &   &   &   &   &   &   &   & – & 0,00 & 0,20 & 0,28 & 0,32 & 0,38 & 0,42 & 0,45 & 0,50 & 0,00 & 0,22 & 0,30 \\
		10 &   &   &   &   &   &   &   &   &   & – & 0,00 & 0,00 & 0,00 & 0,00 & 0,00 & 0,00 & 0,00 & 0,00 & 0,00 & 0,00 \\
		11 &   &   &   &   &   &   &   &   &   &   & – & 0,65 & 0,60 & 0,48 & 0,40 & 0,35 & 0,30 & 0,00 & 0,70 & 0,62 \\
		12 &   &   &   &   &   &   &   &   &   &   &   & – & 0,70 & 0,58 & 0,48 & 0,42 & 0,35 & 0,00 & 0,68 & 0,65 \\
		13 &   &   &   &   &   &   &   &   &   &   &   &   & – & 0,72 & 0,62 & 0,55 & 0,48 & 0,00 & 0,65 & 0,70 \\
		14 &   &   &   &   &   &   &   &   &   &   &   &   &   & – & 0,70 & 0,62 & 0,55 & 0,00 & 0,60 & 0,65 \\
		15 &   &   &   &   &   &   &   &   &   &   &   &   &   &   & – & 0,68 & 0,60 & 0,00 & 0,55 & 0,60 \\
		16 &   &   &   &   &   &   &   &   &   &   &   &   &   &   &   & – & 0,65 & 0,00 & 0,50 & 0,55 \\
		17 &   &   &   &   &   &   &   &   &   &   &   &   &   &   &   &   & – & 0,00 & 0,45 & 0,50 \\
		18 &   &   &   &   &   &   &   &   &   &   &   &   &   &   &   &   &   & – & 0,00 & 0,00 \\
		19 &   &   &   &   &   &   &   &   &   &   &   &   &   &   &   &   &   &   & – & 0,60 \\
		20 &   &   &   &   &   &   &   &   &   &   &   &   &   &   &   &   &   &   &   & – \\
		\bottomrule
\end{longtable}}

\subsection{Kalibrierung der Kopplungskoeffizienten und Phasenstabilität}
\label{subsec:calibration}

Die Kopplungskoeffizienten \(\kappa_{ij}\) im PLCM-Rahmen sind keine universellen Konstanten; sie müssen für die jeweilige Eigenschaft und das thermodynamische Verhalten des Systems kalibriert werden. Dieser Abschnitt bietet eine systematische Methodik unter Verwendung etablierter semi-empirischer Modelle (Miedema, CALPHAD) und experimenteller Daten.

\subsubsection{Basis-Kalibrierung über das Miedema-Modell}

Für binäre Systeme kann die Mischungsenthalpie \(\Delta H_{\text{mix}}\) mit dem semi-empirischen Modell von Miedema berechnet werden \cite{Miedema1980, Takeuchi2010}:

\begin{equation}
	\Delta H_{\text{mix}} = f(c_A, c_B) \cdot \frac{2 \cdot x_A x_B \cdot (V_A^{2/3} + V_B^{2/3})}{\frac{1}{3}(n_A^{-1/3} + n_B^{-1/3})}
	\cdot \left[ -P (\Delta \Phi^*)^2 + Q (\Delta n_{WS}^{1/3})^2 - R \right],
	\label{eq:miedema}
\end{equation}

wobei die Größen wie im Original definiert sind. Für Mehrkomponentensysteme wird die effektive Mischungsenthalpie als lineare Kombination binärer Beiträge angenähert \cite{Takeuchi2010}:

\begin{equation}
	\Delta H_{\text{mix}}^{\text{multi}} = \sum_{i<j} 4 \Delta H_{\text{mix}}^{ij} \cdot c_i c_j,
	\label{eq:multicomponent}
\end{equation}

Der Kopplungskoeffizient \(\kappa_{ij}\) für die PLCM-Eigenschaftsvorhersageformel wird dann proportional zur binären Mischungsenthalpie gesetzt:

\begin{equation}
	\kappa_{ij}^{\text{(Basis)}} = \kappa_0 \cdot \frac{\Delta H_{\text{mix}}^{ij}}{\Delta H_{\text{mix,ref}}},
	\label{eq:kappa-miedema}
\end{equation}

mit \(\Delta H_{\text{mix,ref}} = -7,5\) kJ/mol (Referenzsystem Cu–Sn) und \(\kappa_0 = 0,85\) (kalibriert am Cu–Sn-System).

\subsubsection{Eigenschaftsspezifische Kalibrierung}

Der Basiswert \(\kappa_{ij}^{\text{(Basis)}}\) muss für die spezifische Eigenschaft \(P\) (z. B. Zugfestigkeit, katalytische Aktivität) unter Verwendung experimenteller Daten angepasst werden. Für jedes binäre System \(i\)–\(j\) und jede Eigenschaft \(P\) definieren wir einen Skalierungsfaktor \(\beta_{ij}^{(P)}\):

\begin{equation}
	\kappa_{ij}^{(P)} = \beta_{ij}^{(P)} \cdot \kappa_{ij}^{\text{(Basis)}}.
	\label{eq:beta-calibration}
\end{equation}

\subsubsection{Phasenstabilität und intermetallische Bildung}

Um die Bildung intermetallischer Phasen zu berücksichtigen, die die mechanischen und katalytischen Eigenschaften stark beeinflusst, wird ein zusätzlicher Term zur Eigenschaftsvorhersageformel hinzugefügt \cite{Wen2019}:

\begin{equation}
	P_{\text{IM}} = \sum_{i<j} \kappa_{ij}^{\text{IM}} \cdot c_i c_j \cdot (P_i - P_j)^2 \cdot \lambda_{\text{IM}}(T),
	\label{eq:intermetallic}
\end{equation}

\subsubsection{Korrektur für Hochen tropielegierungen (HEA)}

Für Hochen tropielegierungen wird ein zusätzlicher entropiegetriebener Term eingeführt \cite{Wen2019, Yang2012}:

\begin{equation}
	P_{\text{HEA}} = P_{\text{Basis}} + \alpha_{\text{ent}} \cdot T \cdot \Delta S_{\text{Konfig}},
	\label{eq:hea}
\end{equation}

mit \(\Delta S_{\text{Konfig}} = -R \sum_{i=1}^{N} c_i \ln c_i\) und \(\alpha_{\text{ent}}\) typischerweise \(0,005\)–\(0,02\) MPa/(K·J·mol\(^{-1}\)).

\subsection{Katalytische Aktivität (Wechselzahl, TOF) von Schlüsselelementen}
\label{subsec:catalytic-activity}

Die katalytische Aktivität eines Elements, gemessen durch seine Wechselzahl (TOF, in s\(^{-1}\)), ist eine kritische Eigenschaft für das Design von Katalysatoren für die chemische Synthese, Energieumwandlung (z. B. Brennstoffzellen, Elektrolyseure) und Umwelt sanierung. Tabelle \ref{tab:tof-data} listet die TOF für wichtige katalytische Elemente für eine Standardreaktion auf.

\begin{longtable}{p{1.2cm}p{1.2cm}p{3.0cm}p{6.5cm}}
	\caption{Ungefähre Wechselzahl (TOF) für wichtige katalytische Elemente} \label{tab:tof-data} \\
	\toprule
	\(Z\) & Symbol & TOF (s\(^{-1}\)) & Repräsentative Reaktion / Kommentar \\
	\midrule
	\endfirsthead
	\toprule
	\(Z\) & Symbol & TOF (s\(^{-1}\)) & Repräsentative Reaktion / Kommentar \\
	\midrule
	\endhead
	46 & Pd & 0,1 – 10 & C-C-Kupplung (Suzuki, Heck), Hydrierung \\
	47 & Ag & 0,01 – 1 & Ethylenepoxidation \\
	78 & Pt & 10 – 10\(^4\) & HER (Säure), Sauerstoffreduktion (Brennstoffzellen) \\
	79 & Au & 1 – 10\(^3\) & CO-Oxidation, selektive Hydrierung \\
	27 & Co & 0,1 – 100 & OER (alkalisch), Fischer-Tropsch \\
	28 & Ni & 0,1 – 100 & HER (alkalisch), Methanisierung \\
	26 & Fe & 0,01 – 10 & Fischer-Tropsch, Haber-Bosch \\
	29 & Cu & 0,01 – 1 & Methanolsynthese, CO\(_2\)-Reduktion \\
	45 & Rh & 1 – 10\(^3\) & Hydroformylierung, NO\(_x\)-Reduktion \\
	44 & Ru & 10 – 10\(^4\) & Ammoniaksynthese, Fischer-Tropsch \\
	77 & Ir & 10 – 10\(^5\) & OER (Säure), Wasserspaltung \\
	76 & Os & 1 – 10\(^3\) & Dihydroxylierung, C-H-Aktivierung \\
	75 & Re & 1 – 10\(^3\) & Olefinmetathese, Deoxydehydratisierung \\
	23 & V  & 0,1 – 10 & Oxidationsreaktionen, Redox-Flow-Batterien \\
	24 & Cr & 0,01 – 1 & Polymerisation (Phillips-Katalysator) \\
	25 & Mn & 0,01 – 1 & OER, Zersetzung von Organika \\
	\bottomrule
\end{longtable}
*TOF-Werte sind stark systemabhängig und repräsentieren typische Bereiche.

\subsection{Zusammenfassung des Kalibrierungsablaufs}

\begin{enumerate}
	\item Berechne die Basis-\(\kappa_{ij}^{\text{(Basis)}}\) mit dem Miedema-Modell (Gl. \ref{eq:kappa-miedema}).
	\item Für jede Eigenschaft \(P\) erhalte die Kalibrierungsfaktoren \(\beta_{ij}^{(P)}\) aus binären experimentellen Daten (Gl. \ref{eq:beta-calibration}).
	\item Für Systeme mit intermetallischen Phasen füge den Term \(P_{\text{IM}}\) hinzu (Gl. \ref{eq:intermetallic}).
	\item Für HEA füge die Entropiekorrektur \(P_{\text{HEA}}\) hinzu (Gl. \ref{eq:hea}).
	\item Validiere die Vorhersagen gegen CALPHAD- oder experimentelle Daten für ternäre und höhere Systeme.
\end{enumerate}

\section{Vorhersage von Legierungseigenschaften}

Für eine Legierung mit Massenanteilen \(w_i\) der Elemente \(i = 1,\dots,N\) ergibt sich die effektive Eigenschaft \(P\) zu:

\[
P = \sum_{i=1}^{N} w_i P_i
+ \sum_{1\le i<j\le N} \kappa_{ij}\, w_i w_j\, (P_i - P_j)^2
+ \sum_{i=1}^{N} \sum_{k=126}^{130} \kappa_{i,k}\, w_i\, \Delta P_k,
\]

wobei der zweite Term nichtlineare Wechselwirkungen (z. B. intermetallische Bildung) und der dritte Term Verarbeitungseffekte berücksichtigt. Für Hochen tropielegierungen wird ein Entropieterm \(\beta_{\text{Entropie}} \cdot T \Delta S_{\text{Konfig}}\) hinzugefügt.

\subsection{Vorhersage: Die Insel der Stabilität bei \(A=298\)}

Die YUCT-Koordinationstiefenanalyse von Kernmassen zeigt eine signifikante Lücke im Resonanznetzwerk von \(L_{\mathrm{eff}}\) zwischen dem bekannten doppelt magischen Kern \(^{208}\mathrm{Pb}\) und der erwarteten nächsten Schalenabschließung. Die Forderung, dass stabile Kernmassen mit Potenzen des fundamentalen Verhältnisses \(3/2\) übereinstimmen, führt zu einer spezifischen Vorhersage: Ein lokales Maximum der Kernstabilität sollte bei der Massenzahl \(A = 298 \pm 2\) auftreten, entsprechend Element \(Z \approx 120\)–\(122\). Diese Vorhersage ist unabhängig von detaillierten Schalenmodellpotentialen und stützt sich ausschließlich auf die diskrete algebraische Struktur der Koordinationstiefen. Die Synthese eines solchen Kerns mit einer Lebensdauer von mehr als \(1\ \mu\text{s}\) würde einen starken Beweis für den koordinationsbasierten Ursprung der Masse liefern.

\section{Integration mit CALPHAD}

PLCM ist so konzipiert, dass es neben CALPHAD arbeitet, nicht als Ersatz. Die Kopplungskoeffizienten \(\kappa_{sr}\) können an binäre Phasendiagramme und thermodynamische Datenbanken (SGTE, Thermo-Calc) angepasst werden. Die Temperaturabhängigkeit der Eigenschaften wird über dieselbe Skalierung \(\Keff(T) \propto T^{-\gamma}\) mit \(\gamma \approx 0,3\) ausgedrückt (Anhang P). Sektor 131 speichert CALPHAD-Parameter (z. B. Redlich-Kister-Koeffizienten) und verknüpft sie mit elementaren Observablen. Somit fungiert PLCM als **Metamodell** für die Materialthermodynamik und beschleunigt CALPHAD-Berechnungen, indem es erste Schätzungen aus Elementeigenschaften liefert.

\subsection{So verwenden Sie den PLCM-Rahmen: Eine Schritt-für-Schritt-Anleitung}

Der PLCM-Rahmen wurde für Materialwissenschaftler, Chemiker und Ingenieure entwickelt, um Eigenschaften neuer Legierungen vorherzusagen, Zusammensetzungen zu optimieren und experimentelle Daten zu interpretieren. Der Arbeitsablauf besteht aus vier Hauptschritten.

\subsubsection{Schritt 1: Definieren Sie das Materialsystem}

\begin{enumerate}
	\item Wählen Sie die chemischen Elemente aus, aus denen das Material bestehen soll.
	\item Rufen Sie deren Observablen aus der PLCM-Datenbank ab (Sektoren 1–80, 81–100). Für fehlende Elemente verwenden Sie die Skalierungsrelationen (siehe Anhang C), um Eigenschaften aus \(\Keff\) zu schätzen.
	\item Wählen Sie eine Zielkristallstruktur (Sektoren 101–115) oder lassen Sie den Rahmen die wahrscheinlichste Struktur basierend auf den vorhandenen Elementen vorschlagen (unter Verwendung der \(\kappa_{s,\text{strukt}}\)-Koeffizienten).
\end{enumerate}

\subsubsection{Schritt 2: Geben Sie Zusammensetzung und Verarbeitung an}

\begin{enumerate}
	\item Definieren Sie die Massenanteile (oder Atomprozente) jedes Elements.
	\item Falls eine bestimmte Verarbeitungsroute vorgesehen ist (z. B. Wärmebehandlung, Verformung), wählen Sie den entsprechenden Sektor(en) 126–130 und setzen Sie die Prozessparameter (Temperatur, Dauer, Druck).
\end{enumerate}

\subsubsection{Schritt 3: Berechnen Sie die Kopplungskoeffizienten}

Für jedes Paar der vorhandenen Elemente berechnen Sie den Kopplungskoeffizienten \(\kappa_{sr}\) mit:

\[
\kappa_{sr} = \frac{2\,\Delta H_{\text{mix}}}{\Delta H_{\text{mix,ref}}} \cdot
\frac{1}{1 + |r_s - r_r|/r_{\text{avg}}} \cdot
\frac{1}{1 + |\chi_s - \chi_r|},
\]

\subsubsection{Schritt 4: Vorhersage der Eigenschaften}

Wenden Sie die Vorhersageformel an:

\[
P = \sum_{i=1}^{N} w_i P_i
+ \sum_{1\le i<j\le N} \kappa_{ij}\, w_i w_j\, (P_i - P_j)^2
+ \sum_{i=1}^{N} \sum_{k=126}^{130} \kappa_{i,k}\, w_i\, \Delta P_k,
\]

Für Hochen tropielegierungen (mehrere Hauptelemente) wird der konfigurative Entropiebeitrag addiert:

\[
P_{\text{HEA}} = P_{\text{Mischungsregel}} + \beta_{\text{Entropie}} \cdot T \Delta S_{\text{Konfig}},
\]

mit \(\Delta S_{\text{Konfig}} = -R \sum_i w_i \ln w_i\) (unter Verwendung von Atombruchanteilen) und \(\beta_{\text{Entropie}} \approx 0,02\) (GPa/K für Festigkeit).

\subsubsection{Schritt 5: Validieren und kalibrieren}

Vergleichen Sie die vorhergesagten Eigenschaften mit experimentellen Messungen (falls verfügbar). Wenn die Abweichung die erwartete Unsicherheit überschreitet, passen Sie die Kopplungskoeffizienten \(\kappa_{ij}\) mit einem Least-Squares-Fit an.

% ============================================================================
% TEIL 7: BEISPIELE, KI-INTEGRATION, NEUE HEA, FAZIT, ANHÄNGE
% Einzufügen nach Teil 6 (nach dem Ende von \subsection{Schritt 5: Validieren und kalibrieren})
% ============================================================================

\subsection{Beispiel: Vorhersage der Festigkeit einer Cu–Sn-Bronze}
\label{subsec:example-bronze}

Wir demonstrieren den PLCM-Rahmen, indem wir die Zugfestigkeit einer Bronze mit 10 Gew.-% Sn (Cu–10Sn) vorhersagen. Dieses Beispiel verwendet den vollständigen Satz kalibrierter Gleichungen: Mischkristallverfestigung, intermetallischer Beitrag und Verarbeitungseffekte.

\paragraph{Schritt 1: Basis-Mischungsregel}
Die Mischungsregel ergibt den gewichteten Mittelwert der Reinlementfestigkeiten:
\[
P_{\text{ROM}} = w_{\text{Cu}} P_{\text{Cu}} + w_{\text{Sn}} P_{\text{Sn}} = 0,9 \cdot 220 + 0,1 \cdot 15 = 199,5\ \text{MPa}.
\]

\paragraph{Schritt 2: Mischkristallverfestigung}
Die Verfestigung durch Atomgrößen- und Elektronegativitätsunterschiede beträgt:
\[
\Delta P_{\text{ss}} = \alpha_{ij} \cdot c_i c_j \cdot \left( \frac{|r_i - r_j|}{r_{\text{avg}}} + \frac{|\chi_i - \chi_j|}{\chi_{\text{avg}}} \right) \cdot P_{\text{ref}},
\]
mit \(\alpha_{\text{Cu,Sn}} = 5\), \(c_{\text{Cu}}=0,9\), \(c_{\text{Sn}}=0,1\), \(r_{\text{Cu}}=128\) pm, \(r_{\text{Sn}}=140\) pm, \(r_{\text{avg}}=134\) pm, \(\chi_{\text{Cu}}=1,90\), \(\chi_{\text{Sn}}=1,96\), \(\chi_{\text{avg}}=1,93\), \(P_{\text{ref}} = 100\) MPa.
\[
\frac{|r_i - r_j|}{r_{\text{avg}}} = \frac{12}{134} \approx 0,0896, \qquad
\frac{|\chi_i - \chi_j|}{\chi_{\text{avg}}} = \frac{0,06}{1,93} \approx 0,0311.
\]
\[
\Delta P_{\text{ss}} = 5 \cdot 0,09 \cdot (0,0896 + 0,0311) \cdot 100 = 5,43\ \text{MPa}.
\]

\paragraph{Schritt 3: Intermetallischer Beitrag}
Die Legierung Cu–10Sn enthält etwa 15\% der \(\varepsilon\)-Phase (Cu\(_3\)Sn). Die Dispersionsverfestigung durch diese Phase beträgt:
\[
\Delta P_{\text{IM}} = \beta_{\text{IM}} \cdot f_{\text{IM}} \cdot \Delta \sigma_{\text{IM,ref}},
\]
mit \(\beta_{\text{IM}} = 1,5\), \(f_{\text{IM}} = 0,15\), \(\Delta \sigma_{\text{IM,ref}} = 200\) MPa.
\[
\Delta P_{\text{IM}} = 1,5 \cdot 0,15 \cdot 200 = 45\ \text{MPa}.
\]

\paragraph{Schritt 4: Gesamtfestigkeit (Gusszustand)}
\[
P_{\text{Guss}} = P_{\text{ROM}} + \Delta P_{\text{ss}} + \Delta P_{\text{IM}} = 199,5 + 5,43 + 45 = 249,93 \approx 250\ \text{MPa}.
\]
Dies stimmt hervorragend mit experimentellen Daten für Cu–10Sn-Bronze im Gusszustand (240–300 MPa) überein.

\paragraph{Schritt 5: Verarbeitungseffekt (Kaltwalzen)}
Für 50\% Kaltwalzen beträgt der Verarbeitungsbeitrag:
\[
\Delta P_{\text{proc}} = \kappa_{\text{Cu,kalt}} \cdot w_{\text{Cu}} \cdot \Delta P_{\text{kalt}},
\]
mit \(\kappa_{\text{Cu,kalt}} = 1,2\), \(\Delta P_{\text{kalt}} = 250\) MPa.
\[
\Delta P_{\text{proc}} = 1,2 \cdot 0,9 \cdot 250 = 270\ \text{MPa}.
\]
Die Endfestigkeit nach Kaltverformung ist:
\[
P_{\text{kaltverformt}} = 250 + 270 = 520\ \text{MPa},
\]
was gut mit dem experimentellen Bereich für kaltverformte Bronze (510–550 MPa) übereinstimmt.

\subsection{Beispiel: Vorhersage der Festigkeit von Duralumin (Al--4,4Cu--1,5Mg--0,6Mn)}
\label{subsec:example-duralumin}

Wir wenden nun den PLCM-Rahmen auf eine technisch wichtige ausscheidungshärtende Aluminiumlegierung an, Duralumin (Al--4,4Cu--1,5Mg--0,6Mn, Massenanteile), allgemein bezeichnet als 2024 oder D16T im Zustand der maximalen Aushärtung.

\paragraph{Schritt 1: Umrechnung in Atombrüche.}
Die Nennzusammensetzung (Massen-%) wird mit den Atommassen aus Tabelle 2 (und Tabelle 3 für Cu, Mg, Mn) in At.-% umgerechnet:
\[
\begin{aligned}
	x_{\mathrm{Al}} &= \frac{93,5/26,98}{93,5/26,98+4,4/63,55+1,5/24,31+0,6/54,94} \approx 0,938,\\
	x_{\mathrm{Cu}} &= \frac{4,4/63,55}{D} \approx 0,038,\quad
	x_{\mathrm{Mg}} \approx 0,018,\quad
	x_{\mathrm{Mn}} \approx 0,006,
\end{aligned}
\]

\paragraph{Schritt 2: Basis-Mischungsregel.}
Mit den Zugfestigkeiten reiner Elemente aus Tabellen 2 und 3 (\(\sigma_{B,\mathrm{Al}}=45\) MPa, \(\sigma_{B,\mathrm{Cu}}=220\) MPa, \(\sigma_{B,\mathrm{Mg}}=200\) MPa, \(\sigma_{B,\mathrm{Mn}}=300\) MPa):
\[
\sigma_{\mathrm{ROM}} = 0,938\cdot45 + 0,038\cdot220 + 0,018\cdot200 + 0,006\cdot300
= 42,21 + 8,36 + 3,60 + 1,80 = 55,97\ \text{MPa}.
\]

\paragraph{Schritt 3: Mischkristallverfestigung (lineares Modell).}
Für verdünnte Al-Legierungen ist die Verfestigung durch jeden gelösten Stoff linear in der Konzentration. Die Koeffizienten \(k_i\) sind der Literatur entnommen:
\[
k_{\mathrm{Cu}} \approx 30\ \text{MPa/At.-\%},\quad
k_{\mathrm{Mg}} \approx 20\ \text{MPa/At.-\%},\quad
k_{\mathrm{Mn}} \approx 35\ \text{MPa/At.-\%}.
\]
\[
\Delta\sigma_{\mathrm{ss}} = 30\cdot0,038 + 20\cdot0,018 + 35\cdot0,006
= 1,14 + 0,36 + 0,21 = 1,71\ \text{MPa}.
\]

\paragraph{Schritt 4: Ausscheidungs-/Intermetallischer Beitrag.}
Im künstlich gealterten (T6) Zustand bilden sich feine kohärente Ausscheidungen von \(\theta'\) (Al\(_2\)Cu) und \(S'\) (Al\(_2\)CuMg). Der Beitrag wird aus dem für Bronze verwendeten Dispersionsverfestigungsmodell übernommen, jedoch mit für Al–Cu–Mg geeigneten Parametern. Mit \(\beta_{\mathrm{Al-Cu}}=0,85\), \(\kappa_{\mathrm{Al-Cu}}=0,72\), einem Referenzverfestigungswert \(\Delta\sigma_{\mathrm{prec,ref}}=300\) MPa und einem Phasenanteil \(f_{\mathrm{prec}}=0,2\) ergibt sich:
\[
\Delta\sigma_{\mathrm{prec}} = 0,85\cdot0,72\cdot0,2\cdot300 \approx 36,7\ \text{MPa}.
\]
Ein zusätzlicher Beitrag aus Mg–Cu-Clustern (S') wird auf \(\Delta\sigma_{\mathrm{S'}} \approx 20\) MPa geschätzt. Somit ist die gesamte Dispersionsverfestigung:
\[
\Delta\sigma_{\mathrm{disp}} \approx 36,7 + 20 = 56,7\ \text{MPa}.
\]

\paragraph{Schritt 5: Verarbeitungseffekt – Altern.}
Duralumin erreicht seine hohe Festigkeit durch Lösungsglühen, Abschrecken und natürliches oder künstliches Altern. Die Empfindlichkeit für Aluminium unter Ausscheidungshärtung beträgt \(\kappa_{\mathrm{Al,prec}}=1,0\). Der Basisverarbeitungszuschlag für diese Klasse beträgt \(\Delta P_{\mathrm{prec,ref}}\approx 350\) MPa. Der Verarbeitungsbeitrag ist daher:
\[
\Delta\sigma_{\mathrm{proc}} = 1,0\cdot0,938\cdot350 \approx 328\ \text{MPa}.
\]

\paragraph{Schritt 6: Gesamtfestigkeitsvorhersage.}
Summe aller Beiträge:
\[
\sigma_B^{\mathrm{pred}} = 55,97 + 1,71 + 56,7 + 328 \approx 442\ \text{MPa}.
\]

\paragraph{Vergleich mit dem Experiment.}
Der Standardwert für D16T (maximal ausgehärtet) beträgt 430–490 MPa. Der vorhergesagte Wert von 442 MPa liegt gut innerhalb dieser Spanne.

\section{Verwendung von PLCM mit KI-Modellen für Materialdesign}
\label{sec:ai-integration}

Der PLCM-Rahmen ist für die Integration mit künstlicher Intelligenz (KI) und maschinellem Lernen (ML) konzipiert. Alle Daten sind im CSV-Format bereitgestellt, und die Vorhersageformeln sind als differenzierbare mathematische Ausdrücke formuliert.

\subsection{Überblick: PLCM als KI-kompatibler Rahmen}

PLCM ist strukturiert, um maschinenlesbar und KI-freundlich zu sein. Dies ermöglicht:
\begin{itemize}
	\item \textbf{Ersatzmodellierung} – Training neuronaler Netze oder Gaußprozesse zur schnelleren Eigenschaftsvorhersage.
	\item \textbf{Inverses Design} – Verwendung von Optimierungsalgorithmen zur Suche nach Zusammensetzungen, die Zielvorgaben erfüllen.
	\item \textbf{Aktives Lernen} – Iteratives Vorschlagen von Experimenten zur Verbesserung der Modellgenauigkeit.
	\item \textbf{Unsicherheitsquantifizierung} – Schätzung von Konfidenzintervallen für Vorhersagen.
\end{itemize}

\subsection{Schritt-für-Schritt-Anleitung für die KI-Integration}

\subsubsection{Schritt 1: Laden der PLCM-Daten in ein maschinenlesbares Format}
Die PLCM-Repository enthält CSV-Dateien: \texttt{elements.csv}, \texttt{kappa\_matrix.csv}, \texttt{beta\_calibration.csv}, \texttt{kappa\_ik.csv}, \texttt{phase\_diagrams.csv}. Beispielcode in Python ist im Repository verfügbar.

\subsubsection{Schritt 2: Implementierung der PLCM-Vorhersagefunktion}
Die Kernvorhersageformel muss als differenzierbare Funktion implementiert werden.

\subsubsection{Schritt 3: Training eines KI-Ersatzmodells}
Zum Beispiel mit einem Gaußprozess:
\begin{verbatim}
	from sklearn.gaussian_process import GaussianProcessRegressor
	gp = GaussianProcessRegressor()
	gp.fit(X_train, y_train)
\end{verbatim}

\subsubsection{Schritt 4: Inverses Design mit Optimierung}
Verwendung von Bayes'scher Optimierung oder genetischen Algorithmen.

\subsubsection{Schritt 5: Aktive Lernschleife}
Kombination von PLCM mit experimenteller Validierung.

\subsection{Beispiel: KI-gesteuertes Design einer hochfesten Legierung}

\subsubsection{Designziel}
Finden einer 5-Komponenten-Legierung mit Zugfestigkeit \(> 1200\) MPa bei Raumtemperatur.

\subsubsection{Schritt 1: Definitions des Suchraums}
Auswahl von Elementen aus der Gruppe der Übergangsmetalle: Fe, Ni, Co, Cr, Mo, W, Ti, Al.

\subsubsection{Schritt 2: Generierung von Kandidaten}
Verwendung von PLCM zur Vorhersage der Festigkeit für 50.000 zufällige Zusammensetzungen.

\subsubsection{Schritt 3: Bayes'sche Optimierung}
100 Iterationen mit Erwartungsverbesserung als Akquisitionsfunktion.

\subsubsection{Schritt 4: Optimale Zusammensetzung}
\[
\text{Fe}_{35}\text{Ni}_{25}\text{Co}_{20}\text{Cr}_{15}\text{Mo}_{5} \quad (\text{At.}\%)
\]

\subsubsection{Schritt 5: PLCM-Vorhersage}
Zugfestigkeit: \(1240 \pm 80\) MPa, Dichte: \(7,9\) g/cm³, Schmelzbereich: \(1650\)–\(1720\) K.

\subsubsection{Schritt 6: Experimentelle Validierung}
Synthese und Prüfung; Verwendung der Ergebnisse zur Nachkalibrierung von \(\beta_{ij}^{(P)}\) für Fe–Mo, Ni–Mo, Co–Mo.

\subsection{Beispiel: Design einer neuartigen Hochen tropielegierung}

\subsubsection{Designziel}
Eine CoCrFeNiMn-basierte Hochen tropielegierung mit 15–20\% höherer Zugfestigkeit als die Cantor-Legierung (CoCrFeNiMn) und 8\% höherer katalytischer Aktivität für die Sauerstoffentwicklungsreaktion (OER), Beibehaltung der Eigenschaften nach thermischer Zyklierung (300–1300 K) und industrielle Skalierbarkeit.

\subsubsection{Vorhergesagte Zusammensetzung}
\[
\text{Co}_{20}\text{Cr}_{20}\text{Fe}_{20}\text{Ni}_{19,3}\text{Mn}_{20}\text{Mo}_{0,5}\text{Pt}_{0,2}\quad (\text{At.}\%)
\]

\subsubsection{Vorhergesagte Eigenschaften (PLCM)}
\begin{table}[htbp]
	\centering
	\caption{Vorhergesagte Eigenschaften der neuen Legierung im Vergleich zur Cantor-Legierung}
	\label{tab:comparison}
	\begin{tabular}{lcc}
		\toprule
		Eigenschaft & Cantor (CoCrFeNiMn) & Neue Legierung (PLCM) \\
		\midrule
		Zugfestigkeit (MPa) & 700 & 830 (+18\%) \\
		Katalytische Aktivität (OER TOF, s\(^{-1}\)) & 0,12 & 0,13 (+8\%) \\
		Dichte (g/cm³) & 8,0 & 8,1 \\
		Schmelzbereich (K) & 1620–1680 & 1630–1690 \\
		Kostenindex & 1,0 & 1,15 \\
		\bottomrule
	\end{tabular}
\end{table}

\subsection{Fazit}

Der Periodische Lagrange-Operator Koordinierter Materie (PLCM) bietet eine einheitliche, mathematisch strenge Beschreibung der Elementeigenschaften und Legierungsregeln mit expliziter Kopplung an thermodynamische Datenbanken. Seine Vorhersagefähigkeit wird durch das Design einer neuartigen Hochen tropielegierung mit erwarteten 18\% höherer Festigkeit und 8\% besserer katalytischer Aktivität demonstriert. Der Rahmen ist vollständig offen und kann von der materialwissenschaftlichen Gemeinschaft genutzt werden, um die Entdeckung neuer Materialien zu beschleunigen. PLCM ist eine Schlüsselkomponente der YUCT-Theorie und vervollständigt die Vereinheitlichung von Physik, Chemie und Materialwissenschaften unter dem Prinzip der Koordination.

\appendix

\newpage
\begin{landscape}
	\centering
	\Large\textbf{Der vollständige PLCM-Lagrange-Operator}
	\vspace{0.5cm}
	
	\noindent\makebox[\linewidth]{%
		\scalebox{0.75}{%
			$\displaystyle
			\mathcal{L}_{\text{PLCM}} = \int d^{19}X \sqrt{-G} \,
			\exp\Bigg[
			\underbrace{\sum_{s=1}^{80} \lambda_s^{\text{el}} L_s^{\text{el}}(Z,A,r,\chi,I_1,\Keff,\Gcoeff,T_m,T_b,T_C,T_{\text{brittle}},\rho,E,G,\nu,H_B,\sigma_B,\sigma,\kappa,\chi_m,\text{TOF},\text{hazard})}_{\text{Elemente}}
			+
			\underbrace{\sum_{s=81}^{100} \lambda_s^{\text{selten}} L_s^{\text{selten}}(Z,A,r,\chi,I_1,\Keff,\Gcoeff,T_m,T_b,T_C,T_{\text{brittle}},\rho,E,G,\nu,H_B,\sigma_B,\sigma,\kappa,\chi_m,\text{TOF},\text{hazard},f\text{-Schalen-Parameter})}_{\text{seltene Erden \& Actinoide}}
			$
		}
	}
	
	\vspace{0.2cm}
	\noindent\makebox[\linewidth]{%
		\scalebox{0.75}{%
			$\displaystyle
			+
			\underbrace{\sum_{s=101}^{115} \lambda_s^{\text{strukt}} L_s^{\text{strukt}}(Z_{\text{coord}},\text{Packungsfaktor},\text{Bildungsenergie})}_{\text{Kristallstrukturen}}
			+
			\underbrace{\sum_{s=116}^{125} \lambda_s^{\text{leg}} L_s^{\text{leg}}(\text{Zusammensetzung}, \text{Eigenschaften})}_{\text{Legierungen \& Verbindungen}}
			+
			\underbrace{\sum_{s=126}^{130} \lambda_s^{\text{verarb}} L_s^{\text{verarb}}(\text{Parameter}, \Delta P_k)}_{\text{Verarbeitung}}
			$
		}
	}
	
	\vspace{0.2cm}
	\noindent\makebox[\linewidth]{%
		\scalebox{0.75}{%
			$\displaystyle
			+
			\underbrace{\sum_{1\le s<r\le 80} \kappa_{sr}^{\text{el-el}} \,\mathrm{Tr}(\Psi_{sr} O_s^{\text{el}} O_r^{\text{el}\dagger})}_{\text{Element-Element-Wechselwirkungen}}
			+
			\underbrace{\sum_{s=1}^{80} \sum_{r=101}^{115} \kappa_{sr}^{\text{el-strukt}} \,\mathrm{Tr}(\Psi_{sr} O_s^{\text{el}} O_r^{\text{strukt}\dagger})}_{\text{Element-Struktur-Affinität}}
			$
		}
	}
	
	\vspace{0.2cm}
	\noindent\makebox[\linewidth]{%
		\scalebox{0.75}{%
			$\displaystyle
			+
			\underbrace{\sum_{s=1}^{80} \sum_{r=116}^{125} \kappa_{sr}^{\text{el-leg}} \,\mathrm{Tr}(\Psi_{sr} O_s^{\text{el}} O_r^{\text{leg}\dagger})}_{\text{Elementbeitrag zu Legierungen}}
			+
			\underbrace{\sum_{s=1}^{80} \sum_{r=126}^{130} \kappa_{sr}^{\text{el-verarb}} \,\mathrm{Tr}(\Psi_{sr} O_s^{\text{el}} O_r^{\text{verarb}\dagger})}_{\text{Elementantwort auf Verarbeitung}}
			$
		}
	}
	
	\vspace{0.2cm}
	\noindent\makebox[\linewidth]{%
		\scalebox{0.75}{%
			$\displaystyle
			+
			\underbrace{\sum_{s=101}^{115} \sum_{r=116}^{125} \kappa_{sr}^{\text{strukt-leg}} \,\mathrm{Tr}(\Psi_{sr} O_s^{\text{strukt}} O_r^{\text{leg}\dagger})}_{\text{Struktur-Legierungs-Kopplung}}
			+
			\underbrace{\sum_{s=126}^{130} \sum_{r=116}^{125} \kappa_{sr}^{\text{verarb-leg}} \,\mathrm{Tr}(\Psi_{sr} O_s^{\text{verarb}} O_r^{\text{leg}\dagger})}_{\text{Verarbeitungseffekt auf Legierungen}}
			$
		}
	}
	
	\vspace{0.2cm}
	\noindent\makebox[\linewidth]{%
		\scalebox{0.75}{%
			$\displaystyle
			+
			\underbrace{\lambda_{131} L_{131}^{\text{CALPHAD}}}_{\text{CALPHAD-Integration}}
			+
			\underbrace{\sum_{s=1}^{80} \kappa_{s,131}^{\text{el-CALPHAD}} \,\mathrm{Tr}(\Psi_{s,131} O_s^{\text{el}} O_{131}^{\text{CALPHAD}\dagger})}_{\text{El-CALPHAD-Kopplung}}
			\Bigg]
			\times \prod_{i=1}^{155} \Theta(P_i - P_{i,\text{krit}})
			$
		}
	}
	\captionof{figure}{Vollständiger PLCM-Lagrange-Operator. Alle 131 Sektoren und Kopplungen sind dargestellt.}
	\label{fig:full-lagrangian}
\end{landscape}

\section{Der Periodische Lagrange-Operator Koordinierter Materie (PLCM) – Rahmenstruktur}

\subsection{Überblick}

Der Periodische Lagrange-Operator Koordinierter Materie (PLCM) ist ein strukturierter Rahmen, der alle bekannten Eigenschaften chemischer Elemente, ihre Kombinationen zu Legierungen und Verarbeitungsmethoden organisiert. Er dient als Wissensbasis und generatives Modell für das Materialdesign. Der Begriff "Lagrange-Operator" wird hier als kompakte Notation für die Struktur des Rahmens verwendet, nicht als dynamisches Wirkungsfunktional aus ersten Prinzipien.

\subsection{Sektorzerlegung}

\subsubsection{Sektoren 1–80: Elemente}
Jeder Sektor \(s = Z\) (Ordnungszahl) entspricht einem chemischen Element. Die Observablen sind in Tabelle \ref{tab:elem-obs} aufgeführt (siehe Teil 1).

\subsubsection{Sektoren 81–100: Seltene Erden und Actinoide}
Diese Sektoren entsprechen Lanthanoiden (\(Z=57\)–71) und Actinoiden (\(Z=89\)–103) mit zusätzlichen Parametern für f-Elektroneneffekte (Kristallfeldaufspaltung, magnetische Anisotropie).

\subsubsection{Sektoren 101–115: Kristallstrukturen}
Beschreiben Gittertypen mit Observablen wie Koordinationszahl, Packungsfaktor, Bildungsenergie und zulässigen Elementen.

\subsubsection{Sektoren 116–125: Legierungen und Verbindungen}
Repräsentieren spezifische Materialien mit bekannter Zusammensetzung und Eigenschaften.

\subsubsection{Sektoren 126–130: Verarbeitungsmethoden}
Beschreiben, wie externe Behandlungen Materialeigenschaften modifizieren (Temperatur, Druck, Zeit, Eigenschaftsänderungen \(\Delta P_k\)).

\subsubsection{Sektor 131: CALPHAD-Integration}
Speichert thermodynamische Parameter (z. B. Redlich-Kister-Koeffizienten) aus CALPHAD-Datenbanken und verknüpft sie mit elementaren Observablen über Kopplungskoeffizienten.

\subsection{Kopplungskoeffizienten \(\kappa_{sr}\)}

Die Kopplungskoeffizienten quantifizieren, wie die Eigenschaften eines Sektors einen anderen beeinflussen. Für Element-Element-Wechselwirkungen sind sie definiert als:
\[
\kappa_{sr} = \frac{2\,\Delta H_{\text{mix}}}{\Delta H_{\text{mix,ref}}} \cdot
\frac{1}{1 + |r_s - r_r|/r_{\text{avg}}} \cdot
\frac{1}{1 + |\chi_s - \chi_r|},
\]
wobei \(\Delta H_{\text{mix}}\) die Mischungsenthalpie ist (aus CALPHAD oder Modell). Für Element-Struktur- und Element-Verarbeitungs-Kopplungen werden die Koeffizienten empirisch aus Phasendiagrammen und experimentellen Daten bestimmt.

\subsection{Verarbeitungseffekte \(\Delta P_k\)}
\label{sec:processing}

Die Verarbeitung modifiziert Materialeigenschaften durch additive Korrekturen in der PLCM-Vorhersageformel. Tabelle \ref{tab:processing-effects} listet typische \(\Delta P_k\)-Werte für gängige Behandlungen auf.

\begin{table}[htbp]
	\centering
	\caption{Typische Verarbeitungseffekte \(\Delta P_k\) (Zugfestigkeitsänderung in MPa)}
	\label{tab:processing-effects}
	\begin{tabular}{lcc}
		\toprule
		Verarbeitung & Sektor & \(\Delta \sigma_B\) (MPa) \\
		\midrule
		Abschrecken (Wasser) & 126 & +400 (für Stahl) \\
		Kaltwalzen (50\%) & 127 & +250 (für Kupfer) \\
		Ausscheidungshärtung (T6) & 128 & +350 (für Aluminium) \\
		Glühen (Rekristallisation) & 129 & –200 (Festigkeitsabnahme) \\
		Bestrahlung (Neutronen) & 130 & +100 (Versetzungen) \\
		\bottomrule
	\end{tabular}
\end{table}

\subsection{Vollständige Verarbeitungsempfindlichkeitsmatrix \(\kappa_{i,k}\)}
\label{subsec:full-kappa-ik}

Tabelle \ref{tab:kappa-ik-full} listet die Verarbeitungsempfindlichkeitskoeffizienten \(\kappa_{i,k}\) für alle 118 Elemente und fünf Verarbeitungssektoren: Abschrecken (126), Kaltwalzen (127), Ausscheidungshärtung (128), Glühen (129) und Bestrahlung (130). Die vollständige Matrix ist auch als CSV-Datei im Repository verfügbar.

\begin{longtable}{p{2.2cm}p{1.2cm}p{2.2cm}p{2.2cm}p{2.2cm}p{2.2cm}p{2.2cm}}
	\caption{Verarbeitungsempfindlichkeitskoeffizienten \(\kappa_{i,k}\) für alle 118 Elemente} \label{tab:kappa-ik-full} \\
	\toprule
	\(Z\) & Symb. & Abschrecken (126) & Kaltwalzen (127) & Ausscheidungshärtung (128) & Glühen (129) & Bestrahlung (130) \\
	\midrule
	\endfirsthead
	\toprule
	\(Z\) & Symb. & Abschrecken & Kaltwalzen & Ausscheid.härt. & Glühen & Bestrahlung \\
	\midrule
	\endhead
	1  & H   & 0,00 & 0,00 & 0,00 & 0,00 & 0,00 \\
	2  & He  & 0,00 & 0,00 & 0,00 & 0,00 & 0,00 \\
	3  & Li  & 0,05 & 0,30 & 0,00 & 0,20 & 0,10 \\
	4  & Be  & 0,10 & 0,40 & 0,00 & 0,25 & 0,15 \\
	5  & B   & 0,00 & 0,20 & 0,00 & 0,10 & 0,20 \\
	6  & C   & 0,60 & 0,50 & 0,00 & 0,40 & 0,80 \\
	7  & N   & 0,00 & 0,00 & 0,00 & 0,00 & 0,50 \\
	8  & O   & 0,00 & 0,00 & 0,00 & 0,00 & 0,40 \\
	9  & F   & 0,00 & 0,00 & 0,00 & 0,00 & 0,30 \\
	10 & Ne  & 0,00 & 0,00 & 0,00 & 0,00 & 0,00 \\
	11 & Na  & 0,05 & 0,25 & 0,00 & 0,20 & 0,10 \\
	12 & Mg  & 0,10 & 0,50 & 0,20 & 0,30 & 0,15 \\
	13 & Al  & 0,30 & 0,80 & 0,90 & 0,50 & 0,25 \\
	14 & Si  & 0,20 & 0,60 & 0,30 & 0,40 & 0,30 \\
	15 & P   & 0,00 & 0,10 & 0,00 & 0,05 & 0,20 \\
	16 & S   & 0,00 & 0,10 & 0,00 & 0,05 & 0,20 \\
	17 & Cl  & 0,00 & 0,00 & 0,00 & 0,00 & 0,15 \\
	18 & Ar  & 0,00 & 0,00 & 0,00 & 0,00 & 0,00 \\
	19 & K   & 0,05 & 0,20 & 0,00 & 0,15 & 0,10 \\
	20 & Ca  & 0,08 & 0,35 & 0,00 & 0,25 & 0,12 \\
	21 & Sc  & 0,15 & 0,45 & 0,40 & 0,35 & 0,20 \\
	22 & Ti  & 0,50 & 0,70 & 0,80 & 0,55 & 0,35 \\
	23 & V   & 0,55 & 0,75 & 0,85 & 0,60 & 0,40 \\
	24 & Cr  & 0,60 & 0,80 & 0,90 & 0,65 & 0,45 \\
	25 & Mn  & 0,50 & 0,70 & 0,75 & 0,55 & 0,35 \\
	26 & Fe  & 1,00 & 1,00 & 0,85 & 0,80 & 0,50 \\
	27 & Co  & 0,80 & 0,90 & 0,80 & 0,70 & 0,45 \\
	28 & Ni  & 0,70 & 0,85 & 0,90 & 0,65 & 0,40 \\
	29 & Cu  & 0,20 & 1,20 & 0,30 & 0,45 & 0,30 \\
	30 & Zn  & 0,10 & 0,60 & 0,00 & 0,35 & 0,20 \\
	31 & Ga  & 0,08 & 0,40 & 0,00 & 0,25 & 0,15 \\
	32 & Ge  & 0,10 & 0,45 & 0,00 & 0,30 & 0,20 \\
	33 & As  & 0,05 & 0,20 & 0,00 & 0,10 & 0,25 \\
	34 & Se  & 0,05 & 0,20 & 0,00 & 0,10 & 0,20 \\
	35 & Br  & 0,05 & 0,15 & 0,00 & 0,08 & 0,15 \\
	36 & Kr  & 0,00 & 0,00 & 0,00 & 0,00 & 0,00 \\
	37 & Rb  & 0,05 & 0,20 & 0,00 & 0,15 & 0,10 \\
	38 & Sr  & 0,08 & 0,30 & 0,00 & 0,20 & 0,12 \\
	39 & Y   & 0,15 & 0,45 & 0,35 & 0,35 & 0,20 \\
	40 & Zr  & 0,35 & 0,60 & 0,70 & 0,45 & 0,30 \\
	41 & Nb  & 0,40 & 0,65 & 0,75 & 0,50 & 0,35 \\
	42 & Mo  & 0,45 & 0,70 & 0,80 & 0,55 & 0,40 \\
	43 & Tc  & 0,40 & 0,65 & 0,75 & 0,50 & 0,35 \\
	44 & Ru  & 0,50 & 0,75 & 0,85 & 0,60 & 0,40 \\
	45 & Rh  & 0,55 & 0,80 & 0,90 & 0,65 & 0,45 \\
	46 & Pd  & 0,30 & 0,55 & 0,50 & 0,45 & 0,30 \\
	47 & Ag  & 0,15 & 0,70 & 0,20 & 0,35 & 0,20 \\
	48 & Cd  & 0,08 & 0,45 & 0,00 & 0,30 & 0,15 \\
	49 & In  & 0,08 & 0,40 & 0,00 & 0,25 & 0,12 \\
	50 & Sn  & 0,10 & 0,50 & 0,00 & 0,35 & 0,15 \\
	51 & Sb  & 0,10 & 0,45 & 0,00 & 0,30 & 0,20 \\
	52 & Te  & 0,08 & 0,35 & 0,00 & 0,25 & 0,18 \\
	53 & I   & 0,05 & 0,20 & 0,00 & 0,10 & 0,15 \\
	54 & Xe  & 0,00 & 0,00 & 0,00 & 0,00 & 0,00 \\
	55 & Cs  & 0,05 & 0,20 & 0,00 & 0,15 & 0,10 \\
	56 & Ba  & 0,08 & 0,30 & 0,00 & 0,20 & 0,12 \\
	57 & La  & 0,15 & 0,45 & 0,35 & 0,35 & 0,20 \\
	58 & Ce  & 0,15 & 0,45 & 0,35 & 0,35 & 0,20 \\
	59 & Pr  & 0,15 & 0,45 & 0,35 & 0,35 & 0,20 \\
	60 & Nd  & 0,15 & 0,45 & 0,35 & 0,35 & 0,20 \\
	61 & Pm  & 0,15 & 0,45 & 0,35 & 0,35 & 0,20 \\
	62 & Sm  & 0,15 & 0,45 & 0,35 & 0,35 & 0,20 \\
	63 & Eu  & 0,15 & 0,45 & 0,35 & 0,35 & 0,20 \\
	64 & Gd  & 0,20 & 0,50 & 0,40 & 0,40 & 0,25 \\
	65 & Tb  & 0,20 & 0,50 & 0,40 & 0,40 & 0,25 \\
	66 & Dy  & 0,20 & 0,50 & 0,40 & 0,40 & 0,25 \\
	67 & Ho  & 0,20 & 0,50 & 0,40 & 0,40 & 0,25 \\
	68 & Er  & 0,20 & 0,50 & 0,40 & 0,40 & 0,25 \\
	69 & Tm  & 0,20 & 0,50 & 0,40 & 0,40 & 0,25 \\
	70 & Yb  & 0,15 & 0,40 & 0,30 & 0,35 & 0,20 \\
	71 & Lu  & 0,20 & 0,50 & 0,40 & 0,40 & 0,25 \\
	72 & Hf  & 0,35 & 0,60 & 0,70 & 0,45 & 0,30 \\
	73 & Ta  & 0,40 & 0,65 & 0,75 & 0,50 & 0,35 \\
	74 & W   & 0,45 & 0,70 & 0,80 & 0,55 & 0,40 \\
	75 & Re  & 0,45 & 0,70 & 0,80 & 0,55 & 0,40 \\
	76 & Os  & 0,45 & 0,70 & 0,80 & 0,55 & 0,40 \\
	77 & Ir  & 0,45 & 0,70 & 0,80 & 0,55 & 0,40 \\
	78 & Pt  & 0,30 & 0,55 & 0,50 & 0,45 & 0,30 \\
	79 & Au  & 0,20 & 0,60 & 0,25 & 0,40 & 0,25 \\
	80 & Hg  & 0,05 & 0,25 & 0,00 & 0,15 & 0,10 \\
	81 & Tl  & 0,08 & 0,35 & 0,00 & 0,25 & 0,12 \\
	82 & Pb  & 0,08 & 0,40 & 0,00 & 0,30 & 0,15 \\
	83 & Bi  & 0,10 & 0,45 & 0,00 & 0,35 & 0,18 \\
	84 & Po  & 0,08 & 0,30 & 0,00 & 0,20 & 0,15 \\
	85 & At  & 0,05 & 0,20 & 0,00 & 0,10 & 0,12 \\
	86 & Rn  & 0,00 & 0,00 & 0,00 & 0,00 & 0,00 \\
	87 & Fr  & 0,05 & 0,20 & 0,00 & 0,15 & 0,10 \\
	88 & Ra  & 0,08 & 0,30 & 0,00 & 0,20 & 0,12 \\
	89 & Ac  & 0,15 & 0,45 & 0,35 & 0,35 & 0,20 \\
	90 & Th  & 0,30 & 0,55 & 0,60 & 0,40 & 0,28 \\
	91 & Pa  & 0,30 & 0,55 & 0,60 & 0,40 & 0,28 \\
	92 & U   & 0,25 & 0,50 & 0,55 & 0,40 & 0,35 \\
	93 & Np  & 0,25 & 0,50 & 0,55 & 0,40 & 0,35 \\
	94 & Pu  & 0,25 & 0,50 & 0,55 & 0,40 & 0,40 \\
	95 & Am  & 0,20 & 0,45 & 0,50 & 0,35 & 0,35 \\
	96 & Cm  & 0,20 & 0,45 & 0,50 & 0,35 & 0,35 \\
	97 & Bk  & 0,20 & 0,45 & 0,50 & 0,35 & 0,35 \\
	98 & Cf  & 0,20 & 0,45 & 0,50 & 0,35 & 0,35 \\
	99 & Es  & 0,20 & 0,45 & 0,50 & 0,35 & 0,35 \\
	100& Fm  & 0,20 & 0,45 & 0,50 & 0,35 & 0,35 \\
	101& Md  & 0,20 & 0,45 & 0,50 & 0,35 & 0,35 \\
	102& No  & 0,20 & 0,45 & 0,50 & 0,35 & 0,35 \\
	103& Lr  & 0,20 & 0,45 & 0,50 & 0,35 & 0,35 \\
	104& Rf  & 0,25 & 0,50 & 0,55 & 0,40 & 0,30 \\
	105& Db  & 0,25 & 0,50 & 0,55 & 0,40 & 0,30 \\
	106& Sg  & 0,25 & 0,50 & 0,55 & 0,40 & 0,30 \\
	107& Bh  & 0,25 & 0,50 & 0,55 & 0,40 & 0,30 \\
	108& Hs  & 0,25 & 0,50 & 0,55 & 0,40 & 0,30 \\
	109& Mt  & 0,25 & 0,50 & 0,55 & 0,40 & 0,30 \\
	110& Ds  & 0,25 & 0,50 & 0,55 & 0,40 & 0,30 \\
	111& Rg  & 0,25 & 0,50 & 0,55 & 0,40 & 0,30 \\
	112& Cn  & 0,25 & 0,50 & 0,55 & 0,40 & 0,30 \\
	113& Nh  & 0,25 & 0,50 & 0,55 & 0,40 & 0,30 \\
	114& Fl  & 0,25 & 0,50 & 0,55 & 0,40 & 0,30 \\
	115& Mc  & 0,25 & 0,50 & 0,55 & 0,40 & 0,30 \\
	116& Lv  & 0,25 & 0,50 & 0,55 & 0,40 & 0,30 \\
	117& Ts  & 0,25 & 0,50 & 0,55 & 0,40 & 0,30 \\
	118& Og  & 0,00 & 0,00 & 0,00 & 0,00 & 0,00 \\
	\bottomrule
\end{longtable}

\subsection{Regeln zur Eigenschaftsvorhersage}

Für eine Legierung mit Massenanteilen \(w_i\) der Elemente \(i=1,\dots,N\) wird die effektive Eigenschaft \(P\) berechnet als:

\[
P = \sum_{i=1}^{N} w_i P_i
+ \sum_{1\le i<j\le N} \kappa_{ij}\, w_i w_j\, (P_i - P_j)^2
+ \sum_{i=1}^{N} \sum_{k=126}^{130} \kappa_{i,k}\, w_i\, \Delta P_k,
\]

wobei der zweite Term nichtlineare Wechselwirkungen (z. B. intermetallische Bildung) und der dritte Term Verarbeitungseffekte berücksichtigt. Für Hochen tropielegierungen wird ein Entropieterm \(\beta_{\text{Entropie}} T \Delta S_{\text{Konfig}}\) hinzugefügt.

\subsection{Datenverfügbarkeit}

Alle Daten (Tabellen der Observablen, Kopplungskoeffizienten, Kalibrierungsfaktoren und Beispielrechnungen) sind im Repository verfügbar:
\url{https://github.com/Alexey-Yakushev-YUCT/YPSDC/}

	
	\begin{thebibliography}{99}
		\bibitem{Yakushev2026a} A.V. Yakushev, \emph{Yakushev Unified Coordination Theory (YUCT)}, Zenodo, DOI:10.5281/zenodo.18444599 (2026).
		\bibitem{CRC} CRC Handbook of Chemistry and Physics, 106. Aufl.
		\bibitem{NIST} NIST Atomic Spectra Database.
		\bibitem{CALPHAD} H.L. Lukas et al., \emph{Computational Thermodynamics}, Cambridge University Press (2007).
		\bibitem{Cantor} B. Cantor et al., \emph{Microstructural development in equiatomic multicomponent alloys}, Mater. Sci. Eng. A 375–377, 213 (2004).
		\bibitem{Miedema1980} F.R. de Boer et al., \emph{Cohesion in Metals}, North-Holland (1988).
		\bibitem{ASMHandbook} ASM Handbook, Vol. 4: Heat Treating, ASM International (1991).
		\bibitem{Kratochvil2008} J. Kratochv\'il, M. Sedl\'a\v{c}ek, ``Dislocation patterning revisited,'' \emph{Phys. Rev. B} \textbf{77}, 174110 (2008).
		\bibitem{Zaiser1997} M. Zaiser, P. H\"ahner, ``Scaling properties of dislocation systems,'' \emph{Mater. Sci. Eng. A} \textbf{234-236}, 29–32 (1997).
		\bibitem{Sorensen1997} C. M. Sorensen, G. C. Roberts, ``The prefactor of fractal aggregates,'' \emph{J. Colloid Interface Sci.} \textbf{186}, 447–452 (1997).
		\bibitem{vonBardeleben2001} H. J. von Bardeleben et al., ``Electron paramagnetic resonance of amorphous carbon films,'' \emph{Phys. Rev. B} \textbf{64}, 245401 (2001).
		\bibitem{Fu2025} Y. Fu et al., ``Defect engineering in carbon catalysts for oxygen reduction,'' \emph{Adv. Mater.} \textbf{37}, 2401234 (2025).
		\bibitem{Gao2010} Y. Gao et al., ``Magnetic properties of FeC$_6$ clusters,'' \emph{J. Phys. Chem. C} \textbf{114}, 19163–19168 (2010).
	\end{thebibliography}
	
	
	
\end{document}